% ======================================================================
% NOTES FOR THE AUTHOR (rewrite of 2026-book-orig.tex; original untouched)
% ----------------------------------------------------------------------
% Length: about 93,000 words -> 65,000.  All 176 citation keys are kept,
% none added; the file compiles with references.bib and svozil.bib.
%
% STRUCTURE
%  * 13 chapters.  The original chapter "Too Late to Change Course:
%    Demography, Migration ..." held two unrelated texts.  The essay on
%    Dirac and equal opportunity is now its own chapter 7, with the label
%    chap:inheritance-paradox that the last chapter already refers to.
%    Its title ("The Inheritance Paradox") is my proposal.
%  * chapter_outsourced_goodness.tex is rewritten and inlined as ch. 13,
%    so this file stands alone.
%  * Labels changed so that the \ref's of the last chapter resolve:
%    chap:last-permissible-thought, chap:moving-climate,
%    chap:small-wickedness, chap:inheritance-paradox.
%
% PASSAGES I ADDED (not in the original; cut what you dislike)
%  * Ch. 1: opening that ties Universe 25 to the book's thesis; the
%    "grandparents were poorer" argument; the "ring lights" line.
%  * Ch. 2: "favorite weapon of people who are certain they are good";
%    the paragraph on the chuckle having "changed hands".
%  * Ch. 6: reception of Sarrazin's book (chancellor's "not helpful",
%    loss of post and party); the concession that his heredity passages
%    were weak; the Frisch quotation; the UN term "replacement
%    migration"; the section "Who was at the station"; the closing
%    pointer to the stress-test section of ch. 13.
%  * Ch. 7: the paragraph linking "equal opportunity" to sentimentality.
%
% DROPPED OR CHANGED
%  * "Equilibrium of agony" is no longer attributed to Calhoun; I could
%    not confirm the phrase is his.
%  * References to the old working title ("Too Late to Change Course" as
%    the book's title) are rephrased.
%  * Manual section numbers in the climate chapter; the German raw notes
%    after \end{document}.
%  * Ch. 6: "replacement on a civilizational scale", "cognitive
%    preparation" and the "minority in its own cities" wording are
%    replaced by statements of the same facts without those phrases.
%
% OPEN POINTS
%  * Ch. 6 states empirical claims (schools, welfare, crime, city
%    demographics) without citations.  The sourced figures are in ch. 13
%    (Destatis, IAB, BKA, IQB).  Ch. 13 also calls "replacement" one of
%    two over-finished stories; ch. 6 now says it states the case "as its
%    defenders make it" and points forward.  You may want to decide how
%    far ch. 6 should commit.
%  * Overlap remains: Popper's paradox (ch. 8, 9, 13), Nietzsche (9, 13),
%    Shen Te (9, 13), Sinn's green paradox (5, 13), wolves (12, 13).
%    Shortened in ch. 13 with cross-references, not removed.
%  * Section titles: ch. 1-2 use Title Case, the rest sentence case, as
%    in the original.  Spelling is American throughout.
% ======================================================================
\documentclass[12pt,a4paper]{book}
\usepackage[T1]{fontenc}
\usepackage[utf8]{inputenc}
\usepackage{lmodern}
\usepackage[margin=2.5cm]{geometry}
\usepackage{setspace}
\usepackage{titlesec}
\usepackage{microtype}
\usepackage{csquotes}
\usepackage{epigraph}
\usepackage{hyperref}
\usepackage{enumitem}

\setstretch{1.5}
\setlength{\epigraphwidth}{0.88\textwidth}
\setlength{\emergencystretch}{3em}
\titleformat{\chapter}[display]{\normalfont\huge\bfseries}{\chaptername\ \thechapter}{20pt}{\Huge}

% \title{Disregard This Book---It's So Over}
\title{Sentimental Anti-Realism\\[0.6em]
\large Lifting All Four Limbs and Expecting to Float}
\author{Katherine Vane}
\date{}

\begin{document}
\maketitle
\tableofcontents


% ===================== Chapter file ch01 =====================
\chapter[Universe 25 and the Possibility of Collapse]{The Architecture of Paradise: Universe 25 and the Possibility of Collapse}

\epigraph{For an animal so simple as a mouse, the most complex behaviors involve the interrelated set of courtship, maternal care, territorial defense and hierarchical organization. When these fail to develop, the organization that supports life dies.}{John B. Calhoun, \emph{Death Squared}, 1973 \cite{calhoun1973}}

\section{Paradise, 101 Inches Square}

On 9 July 1968, while half the Western world was in the streets demanding paradise, a man in a government laboratory quietly built one. It was a box about 101 inches on a side. He put eight mice in it.

John B. Calhoun had removed everything a mouse has ever had reason to fear. No cats. No cold. No hunger, no thirst, no epidemic. Food arrived without effort and never ran out. There were more nesting places than the colony would ever fill. Then he stepped back to watch what a species does with a world from which hardship has been engineered away.

It died out. Hunger had nothing to do with it; the feeders were full to the end. Nor did the mice run out of room, since the box never came close to holding the number it had beds for. They simply stopped doing the things that make more mice. They stopped courting. They stopped defending nests. They stopped raising their young. The population climbed to about 2,200, stalled, and ran down like a clock nobody bothered to wind \cite{calhoun1973}.

Calhoun called the enclosure Universe 25. It is the most famous cage in behavioral science, and it has been abused accordingly. One camp reads it as a prophecy of our doom. The other waves it away as a mouse story of no interest to creatures with smartphones. Both are being lazy.

I open with it for a narrower reason. The modern West runs on one large sentimental assumption: take away want, and flourishing follows by itself. Feed people, house them, keep them safe and entertained, and the rest---family, purpose, loyalty, the wish to carry on---will take care of itself. This is what it means to lift all four limbs and expect to float. Universe 25 shows what happens when somebody actually runs the test, on animals too simple to lie about the result.

The experiment does not show that we are mice. It shows something smaller and much harder to wave away: a population can be fed, sheltered, and physically healthy, and still fail to pass on the habits that keep it alive. Provision is one thing. A life is another.

\section{The Man Who Went Looking for a Limit}

Calhoun did not stumble onto this. He had been circling it for twenty years.

In his earlier rat colonies he expected the textbook story. Give rats food and water and they multiply until the food runs short. They never got that far. Growth slowed long before the calories gave out. The rats bunched around particular feeders, and among them everything went wrong at once: feeding, fighting, mating, mothering. Infant mortality soared in a colony that was not short of anything. In 1962 he gave the phenomenon a name that has outlived him, the ``behavioral sink'' \cite{calhoun1962}.

What he took from the rats is a distinction worth keeping. Density is arithmetic: bodies per square foot. Crowding is experience: the number of encounters you cannot get out of. A pen can have plenty of floor and still be unbearable if every route to the food runs past somebody who wants a fight. An empty nest box is not available if a dominant male sits on the only ramp that leads to it.

Universe 25 was designed to isolate that second thing. Strip out hunger, weather, and predators, and leave the animals with one problem only: each other.

\section{How the Box Was Built}

The box was a finite structure, and that matters. It held nesting compartments, feeding stations, and water bottles, arranged on several levels so the mice could move between sectors. Food was replenished continuously. The temperature was controlled. The founders were healthy laboratory stock. Calhoun called it a ``mortality-inhibiting environment,'' and by the standards of mouse history it was the Garden of Eden with room service.

Two things the designer did not supply, and could not. The first was a way out. There were places to occupy but no frontier beyond the walls. A mouse that lost its position could shuffle to another compartment; it could not strike out and found a colony somewhere else. The second was any means of changing the arrangements. Every mouse depended on the conduct of every other mouse, permanently, with no appeal.

That makes Universe 25 a poor model of a planet. It is a rather good model of any institution in which all the chairs are taken and nobody is allowed to leave.

\section{Act One: Work to Do}

Calhoun called the first 104 days the Strive Period, and it is the only happy stretch in the story. Note why. There was work.

The eight founders had a world to make. They explored the box, staked out territories, chose nesting sites, courted, and produced litters. They fought, too. Nobody should picture a mouse golden age. But a fight in those days settled something. The loser had somewhere to go and something still to win. An empty box could be taken. A pecking order that had not yet hardened could be challenged.

I am not going to tell you the mice felt purposeful. Nobody knows what a mouse feels. What can be counted is behavior, and the founders displayed the full repertoire of their species: courtship, nest-building, nursing, defense. Their descendants would not.

One more thing happened in Act One, and it is the hinge of the whole experiment. The founders learned the box directly. Their offspring learned it from their mothers. From the first litter on, the colony's future depended less on the steel and the food supply than on what each generation managed to teach the next.

\section{Act Two: Boom}

Then came what Calhoun called the Exploit Period. The population doubled roughly every 55 days. By any measure a ministry would use, the colony was a triumph: numbers up, mortality down, consumption steady.

Underneath, the positions were running out. Young mice matured into a world in which every territory was held and every role filled. Calhoun did not count vacancies the way an economist counts jobs. What he saw was simpler. There were more grown animals than the social order could absorb, and the surplus had nowhere to go.

Here the behavioral sink returned. Mice packed themselves around certain feeders although identical feeders stood free a short walk away. The jammed spots bred fights and disturbance, and the animals kept coming back to them anyway. Whole sectors of the box sat half empty. The colony had more room than it would ever use, and could not make that room usable.

Any country with empty houses and homeless people, unfilled vacancies and unemployed graduates, will find the picture familiar. A resource that exists in the statistics is not yet a resource anyone can reach.

\section{Act Three: The Rot}

Around Day 315 growth slowed sharply, and the behavior turned strange.

Some males became violent in ways that no longer served any territorial purpose; there was nothing left to win and they attacked anyway. Others gave up and sat motionless for hours. Some females defended their nests savagely. Others abandoned them, or dropped and forgot their pups, or injured them. Courtship, which in a mouse is an ordered sequence, dissolved into indiscriminate mounting. Later retellings have milked that detail for scandal. What mattered was that the whole chain had come apart, from courting through pregnancy and nursing to the defense of the young, and a colony with a broken chain has no next generation.

The mothers were the heart of it. A pup needs more than milk. It needs warmth, protection, and the daily example of adults behaving like adults. Pups that got none of this grew up without knowing how to court, how to defend, or how to mother. So the damage was passed on. The colony's trouble stopped being a quarrel among adults over space and became a defect of upbringing, handed down and deepened with every litter.

There was no dramatic crash. That is the most frightening part. The colony settled into a long, grinding stability in which violence was routine and withdrawal was normal, and anyone reading only the food-consumption figures would have reported that all was well. A society can look perfectly steady while the thing that renews it has already stopped.

\section{Act Four: The Beautiful Ones}

At about Day 560 the population peaked near 2,200. The last pups that would survive were born around Day 600. After that, nothing. The colony aged and thinned toward extinction \cite{calhoun1973}.

Among the survivors Calhoun picked out a class of males he called the ``beautiful ones.'' They never fought, so their coats were perfect: no scars, no bitten tails. They ate, slept, and groomed. They showed no interest in females and no interest in territory. They looked like the healthiest animals in the box. They were its obituary.

Calhoun spoke of a ``first death,'' the death of behavior that comes before the death of the body. This is a metaphor, and nobody can weigh it on a scale. But it puts the right question. When do we call a society healthy? When its members are clean, fed, and unscarred? Or when they can still form attachments, take on obligations, and bring up children who can do the same?

\subsection{The Rescue That Failed}

Could the beautiful ones be saved? The story usually told is that Calhoun moved some of them into fresh, uncrowded quarters with normal companions, and that it made no difference. They did not court. They did not claim territory. They did not breed. They went on grooming.

I repeat the story with a warning. These relocation trials are far less well documented than the population curve, and they cannot bear the words people like to hang on them, such as ``permanent'' and ``hard-wired.''

The modest version is bad enough. Once a generation has grown up without learning how to live, more space and more food do not teach it. You cannot repair with supplies what was broken in the upbringing. Recovery, if it comes, needs models, teachers, and time. It needs a generation that still works to bring up the one that does not yet exist.

This is where the mouse box turns into a theory of inheritance. A society hands down more than genes and property. It hands down habits of trust, expectations about marriage and friendship, ways of settling a quarrel, and some picture of a future worth the trouble. When that handover fails, the heirs receive a fully furnished world that they cannot read.

\section{Now the Cold Water}

A book against wishful thinking cannot make an exception for its favorite parable. So here is what Universe 25 does not show.

\emph{It does not show that comfort kills.} Humans build institutions, laws, religions, and trades. We migrate. We redesign our cages. We teach on purpose. No mouse ever founded a school.

\emph{It does not explain falling birth rates.} People have fewer children for many reasons, including housing, money, education, contraception, and plain preference, and the mix differs from country to country.

\emph{It does not show that decline is irreversible.} The claim that social behavior was erased for good rests mostly on Calhoun's own reading of his late generations. Institutions can be rebuilt. Newcomers bring working habits with them. A generation can decide to teach what it was never taught.

\emph{It was never properly replicated.} There was one laboratory, one strain, and one celebrated run, and no program of independent repetition by modern standards. Popular accounts make things worse by splicing the rats of the 1950s onto the mice of 1968 and presenting observations from different colonies as a single story.

\emph{It does not isolate a cause.} Density, forced contact, the layout of the box, failed mothering, stress, and the quirks of one laboratory strain are all candidates. The experiment shows an outcome and leaves the mechanism open.

Notice who needs the experiment to say more than it does. The prophet of doom wants a law of nature, because a law of nature excuses him from doing anything. That is sentimentality as well. Despair is the cheapest of the comfortable feelings.

What survives the cold water is still plenty. Abundance does not produce order by itself. And once the ordinary machinery of learning and taking part has failed, the failure feeds on itself.

\section{The Uncomfortable Resemblance}

Now look up from the box.

Japan has a word, \emph{hikikomori}, for people who have shut the bedroom door on society and stayed behind it for months or years; official surveys there track loneliness and isolation as a national problem \cite{cabinet2025}. Young Chinese have coined \emph{tang ping}, ``lying flat,'' and \emph{bai lan}, ``let it rot.'' Across the rich world birth rates have sunk below replacement level. Everywhere there are young people who are educated, fed, housed after a fashion, endlessly entertained, and unable to see a road into adult life.

None of these people is a mouse, and each case has local causes. Hikikomori does not come from population density. Lying flat is an answer to a particular labor market. A falling birth rate proves no ``first death.'' I grant all of it.

But the polite explanation, that the young simply cannot afford a life, is too easy. Our grandparents were far poorer and had more children, more neighbors, and more to do. The size of the purse has grown. The path has disappeared. A young worker holds a degree and cannot buy a home. A couple wants children and can find no arrangement of work, cost, and care that makes it possible. A citizen has all the information in the world and no institution he trusts to act on it. In each case the material capacity exists and the social route is blocked. That is Calhoun's finding in human dress: plenty of nest boxes, and no way to reach them.

Then there is the grooming. The beautiful ones kept their coats immaculate. We have ring lights. Display is as old as the species, and I do not sneer at it. But a feed turns the self into a permanent project of presentation, and it pays out the reactions of a crowd without the obligations of a single friend. One can be seen by thousands and needed by nobody. A system that offers the sensation of taking part while removing the risks of taking part is a very comfortable cage.

So the useful comparison is between a closed system and an open one, with mouse and man on either side. A closed system offers comfort without exit, contact without belonging, and information without authority. An open one lets people leave, form new groups, invent new roles, and teach new habits. Which of the two our institutions are building is the practical question. The answer is not obviously flattering.

\section{The First Death}

Universe 25 offers no proof that we are finished. It offers a demonstration that material provision and the failure of social life can sit side by side for a long time, each politely ignoring the other.

The death worth fearing comes first and makes no noise. It is the slow disappearance of the practices that hold a common life together: courtship, care, teaching, duty, trust, and the willingness to build something you will not live to enjoy.

Abundance is not the villain. The villain is the belief that abundance is enough, that a society is a feeding arrangement with entertainment attached. A civilization does not decline because its members are comfortable. It declines when comfort is cut loose from obligation, when no believable road leads from youth into adulthood, and when each generation inherits fewer living examples of how people belong to one another.

Can a society keep its purpose once survival is easy? Can it invent new duties when the old necessities fall away? Can it hand down competence, and not only information?

Calhoun's mice could not ask these questions. We can. If the answer is no, no quantity of provision will save us. If it is yes, there is work to do, and the mice were happiest when they had some.


% ===================== Chapter file ch02 =====================
\chapter[The Weaponized Chuckle]{The Weaponized Chuckle: Ridicule, Status, and the Refusal to Argue}

\epigraph{Ridicule is man's most potent weapon.}{Saul D. Alinsky, \emph{Rules for Radicals} \cite{alinsky1971}}

\section{The Laugh That Ends the Conversation}

You know the sound. Someone on a panel says a thing the room does not want to hear. It might be wrong. It might only be unfamiliar, or inconvenient, or unfashionable this year. Before anyone asks for evidence, before a single premise is examined, the person in the next chair gives a short laugh through the nose and glances at the host. The host smiles back. The matter is closed.

Nothing has been refuted. No fact was produced, no counterexample, no reason. The claim has merely been moved to the pile of things that decent people do not bother to answer.

There is a laugh that opens a door and a laugh that shuts one. The first invites you in. The second tells you that you have already been judged. It may be a snort, a theatrical guffaw, a little ``heh,'' or a raised eyebrow shared with the audience. The sound is small. What it does is not.

What it says, without saying it, is: \emph{you cannot be serious}; \emph{everyone can see how foolish this is}; \emph{I will not dignify that with an answer}. It asks the bystanders to reach a verdict before they have heard the case. It trades the question ``Is this true?'' for the question ``What sort of person would say such a thing?''

This chapter is about that laugh, and only that one. I have nothing against jokes. Satire exposes hypocrisy, and a good line can reveal a contradiction faster than a page of solemn prose. The trouble starts when amusement is used \emph{in place of} analysis: when the laugh shows no defect in the argument and simply instructs the audience not to look for one.

It is the favorite weapon of people who are certain they are good. If you already know you are on the right side, an opposing argument is not a problem to be solved. It is a breach of manners, and you deal with it as you deal with any breach of manners: you let the room see you are amused.

\section{A Laugh with No Joke in It}

Most laughter has little to do with jokes. Robert Provine, who went out and recorded the real thing, found that people laugh mostly in ordinary talk, at remarks that are not funny at all \cite{provine2000}. Laughter is social glue. It says ``I'm with you,'' ``I get it,'' ``we're safe here.''

The weaponized laugh steals that machinery. Watch where the laugher looks: at a third party, hardly ever at the person who spoke. The speaker becomes an exhibit, and the laugh a caption that reads \emph{we know what kind of person this is}. That is how it can go off when nobody has said anything funny. It is a frame. It tells the audience how to feel about a sentence before they have worked out what the sentence means.

And it can be denied, which is the beauty of it. An insult is a declaration of war; a laugh is only a laugh. Challenge it, and you will hear that the wording was funny, that no offense was meant, that you are being terribly sensitive. So the target has two choices. He can let it pass, and the frame stands. Or he can object, and confirm the picture of the humorless crank. ``You're too sensitive'' is the second half of the maneuver. First you wound the man, then you cite the wound as evidence against him.

Not every laugh in an argument is of this kind. People laugh from nerves, from relief, at their own mistakes. You tell the difference by timing and direction: who is looked at, what the room does, and above all what comes next. The weaponized laugh is aimed at a person or a position, and its job is to change the terms on which that person will be heard.

\section{The Name in the Manual}

Logic textbooks have a name for this: the \textbf{appeal to ridicule}, also known as the horse laugh or \emph{argumentum ad ridiculum}. You present a claim as laughable and then treat the laughter as evidence against it.

\begin{quote}
Person A: ``The evidence suggests that this policy may have unintended long-term effects.''\\
Person B: ``So you think the entire country should panic because of one spreadsheet?''
\end{quote}

B may have a real objection hidden somewhere. But what the audience is invited to laugh at is a cartoon of what A said. The horse laugh usually travels with a straw man for company.

Humor is not always beside the point. If a man announces that no human being has ever made a mistake and then corrects himself, laughing is a fair comment. And the character of a speaker is sometimes relevant: his expertise, his record, his undeclared interests can all bear on how far to trust him \cite{walton1998}. The fault lies in offering the social verdict as a substitute for the work.

The test is simple. Take the laugh away and see what is left. If the critic can name a false premise, a bad inference, a missing number, or a counterexample, then the laugh was decoration on an argument. If all that remains is ``that's ridiculous,'' the audience has been handed a verdict and no reasons.

Is it an \emph{ad hominem}? Nearly. The classic version says: you are a fool, therefore you are wrong. The laugh runs it backwards: that is absurd, and only a fool would say it. The insult arrives by way of the idea, but it arrives.

Erving Goffman explained why this works so well. Every encounter requires each of us to keep up a presentable face. A public laugh damages that face without any declaration of hostilities, and the victim has to stop and repair it. Before long he is defending his intelligence, his tone, his motives, and his emotional stability, and no longer his argument. Which was the point.

Is it passive aggression? In manner, yes: hostility with a built-in alibi. But the term is too small. It suggests a private sulk, and this is frequently a public technique, practiced and repeated. Researchers who study broadcast debates have shown politicians using snorts and long laughs to do real argumentative work, undercutting an opponent's standing as someone who knows the facts \cite{demasi2019}. When the laugh can be denied, call it passive aggression. When it is aimed at an audience, it is coalition enforcement. When it comes from a superior, it is plain domination. The same sound can be all three at once.

\section{Sudden Glory}

The oldest theory of laughter is the least flattering. Plato thought we laugh with pleasure at the weakness of others. Hobbes put it in two words. Laughter, he wrote in \emph{Leviathan}, comes from ``sudden glory,'' the jolt of finding oneself superior to somebody else \cite{hobbes1651}.

As an account of all humor this is too bleak. As an account of the sneering chuckle it is exact, and it explains the economics. The laugher does not have to demonstrate any competence of his own. The target supplies the contrast. Make the other fellow look naive, and you rise without having achieved anything. It is victory without the expense of combat.

The word ``sudden'' matters too. Following an argument takes effort and time. Ridicule reclassifies a person in a second, and the audience does not calculate the odds that he is wrong. It simply feels him drop.

This also explains why ridicule is addictive. A man who cannot put together a rebuttal can still win the room with a face, a snort, or a comic paraphrase. The applause arrives before anybody knows whether the mocked argument was sound.

Henri Bergson added the social side. Laughter, he said, is a corrective. The group uses it to punish whoever has stepped out of line, be he an eccentric, a pedant, or a man who has misread the room. Michael Billig took the thought to its unpleasant conclusion. Ridicule is one of the main ways ordinary people are trained to conform. We learn to anticipate humiliation and to adjust in advance, before any punishment is announced \cite{billig2005}.

This is why ridicule outperforms insult. An insult wounds one person. Ridicule teaches a whole audience what not to say. The man being laughed at is only the demonstration model. The lesson is for the bystanders, and they take it. Someone who has watched a colleague laughed out of a meeting will prune his own remarks for years.

\section{Who May Laugh at Whom}

Laughter is a status signal, and the research on voice and dominance says so in so many words: laughing lets people display superiority in a form less confrontational than open aggression \cite{bachorowski2016}. A shove invites a shove back. A laugh establishes the same ranking and looks like play.

The powerful do not even need to explain the joke. When the boss laughs, the room laughs, because they see the hierarchy, or want to be near it, or know that silence will be read as disloyalty. So a room may laugh at an argument because the laugh came from a strong chair, and for no other reason.

Laughter from below is a different animal. The prisoner laughing at the warden, the dissident at the Party, the subordinate at the pompous chief: this is how the powerless refuse the solemnity that power demands. The sound is the same and the direction is opposite. So ``Is it funny?'' is never the interesting question. Ask who is permitted to laugh at whom, and what happens to the person who declines to join in.

A laugh also counts heads. It says ``I am above you,'' and it says ``look who is with me.'' A man who laughs alone in a meeting loses a little. A man who draws laughter from three allies has built a coalition on the spot, and his target now faces what looks like a unanimous room.

One observation, since this is a book about the self-righteous. For most of history the knowing chuckle belonged to bishops, colonels, and headmasters. It has changed hands. Today it is most at home in the studios, faculties, and editorial offices where everyone agrees on what a decent person thinks. Much professional comedy has followed it there, and has gone from mocking the powerful to mocking whoever doubts them. The audience no longer laughs because something is funny. It applauds because something is correct.

\section{The Economics of a Snort}

A public argument is never a duel. There is a speaker, a respondent, and an audience that must decide whom to believe and usually cannot judge the matter on its merits. So the audience reads cues: confidence, fluency, titles, applause, silence, and laughter.

In that setting a real refutation is expensive. It takes time. It may expose the refuter's own ignorance. It concedes that the opponent is worth answering. A laugh costs nothing and commits the laugher to nothing.

In the economist's language the laugh is a \emph{cheap signal} dressed up as a costly one. It behaves as though a crushing rebuttal existed and were being withheld only because the point is too obvious to need it. The audience is asked to infer competence from a refusal to show any.

It also works as a focal point. The room is unsure how to take a claim, and the laugh supplies an answer everyone can converge on: this is absurd. If enough people converge, the speaker's standing collapses although nobody has looked at his evidence. The laugh has not changed the facts. It has changed what everyone expects everyone else to think of them.

Repeat the game and the laugh becomes a punishment. Somebody keeps raising awkward objections, and answering each one would be work. A laugh puts a tax on every future objection. He learns that speaking costs him, and the laughers learn that mockery is cheaper than engagement. After a while the group shows a fine consensus. Nobody has been persuaded. Dissent has been priced out of the market.

The trap has no good exit. If the target gets angry, he has lost his composure and therefore the exchange. If he explains at length, he is overreacting to a bit of fun. If he mocks back, the argument turns into a talent show. In all three cases the original question is gone. He brought evidence to a status contest, and he will go home wondering why evidence did not win.

\section{Rule Five}

The military comparison is more than decoration. Armies are beaten in morale before they are beaten in the field, and an enemy who believes resistance is hopeless need not be shot.

Saul Alinsky wrote the doctrine down. The fifth of his \emph{Rules for Radicals} holds that ridicule is the most potent weapon because it is almost impossible to counterattack and it infuriates the opposition \cite{alinsky1971}. He was describing an operation, with no pretense of argument. Provoke the target. Let the audience watch him lose his temper. Then the anger becomes the story, and the grievance that caused it is forgotten.

Today the trade speaks of delegitimizing and discrediting. A set of claims is tied to clowns, cranks, and extremists. The aim is no longer to convince the target. It is to stop him from convincing anyone else. Ridicule works as area denial: it makes the ground around an idea too dangerous to stand on.

The internet has industrialized the method. A laughing emoji, a clipped reaction video, or a screenshot with a smirking caption travels much further than the argument it mocks, and usually travels without it. Ten thousand people enjoy the ridicule and not one of them has read the original claim. A pile-on looks like overwhelming refutation when not a single participant has addressed the point. Whole communities keep a stock of phrases and images that mark an opponent as ridiculous on sight. Nobody has to think. Joining the laughter proves loyalty, and that is all that is wanted.

So anyone with something unwelcome to say now faces two questions. Is it true? And can I survive saying it? When the second question dominates, people censor themselves. No idea is banned. It is merely made too expensive to utter, and the public sphere narrows without a single prohibition on the books.

Neither political side owns the tactic. Activists mock institutions, and institutions mock activists. The danger is the same in both directions. A public that is laughed at often enough stops distinguishing between ``false'' and ``socially unsafe,'' and after that it is governed by the fear of being laughed at.

\section{An Old Trick}

None of this is new. In 423 BCE Aristophanes put a comic Socrates on the Athenian stage, dangling in a basket and running a school of intellectual swindles. The play is no report on Socratic philosophy. It is a caricature that lumps Socrates together with every sophist the audience disliked. A generation later, in Plato's \emph{Apology}, Socrates tells his jury that his most dangerous accusers are the old ones, the storytellers and comic poets who had shaped opinion about him for decades. Nobody can show that \emph{The Clouds} killed him, and historians rightly refuse so simple a story. But the mechanism is on display. By the time the man stood up to speak, the jury already carried a comic picture of what he was.

It is worth adding that Socrates mocked people himself, and mercilessly. The difference is in where the mockery points. Socratic irony drives the listener back to a contradiction he had not noticed. The horse laugh leads him away from it.

The Romans taught the technique as a craft. Quintilian devotes part of Book VI of his \emph{Institutio Oratoria} to laughter, and he is disarmingly frank. A well-timed jest relaxes the judges, weakens an opponent, and distracts attention from damaging evidence \cite{quintilian1920}. He also warns that too much mockery makes the speaker look petty and suggests he has no better answer. Two thousand years on, that remains the right question to ask about any laugh. Is it exposing a weak argument, or hiding a weak rebuttal?

The history of science contributes its own examples, and they need handling with care. ``They laughed at Galileo'' proves nothing, because they also laughed at a great many fools. Some mocked ideas are false. Some true ideas were opposed with serious arguments and no laughter at all. The pattern that holds is narrower: ridicule can delay inquiry by making a proposal embarrassing before its evidence has been looked at.

After \emph{On the Origin of Species}, cartoonists drew Darwin as an ape. The drawings did not touch natural selection. They converted a biological argument into a question of personal insult \cite{darwinproject2024}. Ignaz Semmelweis showed that doctors who disinfected their hands stopped killing new mothers, and much of his profession resisted and attacked him. The usual morality tale of stupid doctors laughing at an obvious truth leaves out the institutional rivalries, the theoretical disputes, and his own gift for making enemies \cite{semmelweis2024}. The fuller story makes the point better. Ridicule seldom works alone. It is the compact, cheerful expression of a larger refusal, by a guild, a faculty, or a class, to revise its opinion of itself.

In the eighteenth century Lord Shaftesbury turned the matter round and defended ridicule as a \emph{test} of truth. Sound ideas would survive mockery, he held, and humbug would wither under it \cite{shaftesbury2024}. It is a handsome thought with one obvious flaw. Laughter has no neutral jury. Audiences laugh at true things, at foreign accents, at complicated science, and at a speaker's ears. A bad argument can have a true conclusion, and an elegant joke can be false. Ridicule tests social fluency and membership of the group at least as well as it tests truth.

\section{Why the Room Goes Along}

Why do bystanders take a laugh for information? Five reasons, none of them to our credit.

\emph{Social proof.} When unsure, we look at other faces. A laugh tells us what reaction is expected.

\emph{Contagion.} Laughter is physical and it spreads. Even someone who has not followed the argument feels the pull to smile along.

\emph{Laziness.} Judging an argument is work. Ridicule delivers the conclusion, ``not serious,'' and spares us the reasoning.

\emph{Cowardice.} Being wrong in private is uncomfortable, and being wrong in public is costly. The laugh shows where safety lies.

\emph{Tribe.} An argument from outside threatens more than a belief. It threatens the border between people like us and people like them. Laughter is the immune response that rejects the foreign body before it gets into the bloodstream.

This is how a technically strong argument fails in public. It was lost before anyone weighed the evidence, at the moment it was classified as the kind of thing one does not weigh.

\section{When Mockery Is Earned}

A chapter like this must not end by asking for a world without jokes. That would be a worse world, and a more dishonest one.

Ridicule is earned when a claim really does contradict itself and the joke makes the contradiction visible. It is earned when satire holds a public figure to his own stated principle. It is earned when the powerless have no other weapon against an authority that cannot be confronted head on. And it is earned when cruelty and corruption have already been answered with patient argument and have ignored the answer.

To tell the two kinds apart, ask four things. Does the mockery point at a specific defect? Could the mocker state that defect without the joke? Is the target an argument, or merely a member of a disliked group? And when the laughter stops, is anyone still allowed to speak?

Satire is a way of arguing. The horse laugh is a way of not having to.

\section{What to Do About It}

\emph{Name it.} ``That was a laugh, not an answer. Which premise do you reject?'' There is no need to scold. You are only separating the performance from the reasoning it pretends to contain.

\emph{Restate the claim} in its strongest fair form. This takes the cartoon out of the laugher's hands and gives the bystanders a second chance to hear what was said.

\emph{Demand specifics.} A laugh implies that the whole thing is obviously wrong. Ask what, exactly: the evidence, the inference, the definition? That turns a vague air of certainty into a statement that can be disputed.

\emph{Stay calm.} You need not pretend it did not sting. But visible distress is the prize the laugher is playing for. ``I can see you find it funny. I am asking what you think is wrong with it'' is stronger than any insult you could return.

\emph{Leave}, sometimes. If the room has already decided against inquiry, every further sentence is fresh material for mockery. You are not obliged to win every room. Save the argument for one that can hear it.

\emph{Mock back} only with care, and mock the tactic and never the person: ``A powerful refutation. Perhaps you could publish the laugh.'' The danger is that you end up in a contest of sneers, which you may well lose and which nobody will learn anything from. The aim is to show what the laugh is doing. Laughing louder shows nothing.

\section{The Missing Sentence}

The derisive laugh is several things at once, and that is why it works. It is a fallacy when it stands in for evidence. It is an insult with an alibi. It damages a reputation, displays rank, and recruits bystanders, and it warns everyone else in the room. The victim usually answers on one level only, by offering more reasons to a gesture that never had anything to do with reasons.

Neither laughter nor solemnity detects truth. Serious arguments can be funny, and nonsense is often delivered with a straight face. What tells the two laughs apart is the few seconds afterwards. Does anybody explain? Does the room go back to the premises? Is the target invited to reply? Or does everyone move on with the pleasant feeling that a question has been settled, when it has only been made embarrassing to ask?

Laughter is among the oldest tools we have for holding a group together, which is why it can be warm, liberating, and cruel by turns. The weaponized chuckle is dangerous because it is efficient. It ends an argument before the argument starts, and it punishes a speaker without the bother of an accusation. It lets a group feel united when nobody in it knows whether the mocked claim was true.

So when you hear that laugh, listen for the sentence that is missing. What would the laugher have had to show if laughing were not allowed? Which premise would he have had to attack, and with what evidence? More often than not the laugh tells you nothing about the argument. It tells you that somebody found mockery cheaper than refutation, and expected the room to accept the discount.


% ===================== Chapter file ch03 =====================
\chapter{Who Owns the Science? Who Casts the Scientists?}
\label{chap:who-owns-the-science}

\begin{flushright}
\emph{The future is not cancelled by one bad decision.\\
It is made expensive by a thousand reasonable ones.}
\end{flushright}

\section{The lab coat as costume}

Science arrives in public wearing a costume: a white coat, a graph, and a calm voice saying ``the data are clear.'' Behind the costume stand laboratories, universities, ministries, foundations, corporations, publishers, regulators, and public-relations departments, together with a small army of people whose trade is turning uncertainty into a press release.

The costume is useful. It gives strangers a reason to trust one another. It also makes it very easy to mistake the look of knowledge for the possession of it. We are told that science speaks. In fact science is made to speak, through institutions, and institutions have landlords.

In 2022, at a World Economic Forum panel on disinformation, one of the landlords said so out loud. Melissa Fleming, then the United Nations' Under-Secretary-General for Global Communications, put it in a sentence that deserves to be kept whole: ``We own the science, and we think that the world should know it, and the platforms themselves also do'' \cite{FlemingMelissa}.

One can read this kindly. Perhaps she meant only that the United Nations feels responsible for getting reliable information out. But listen to the grammar. Science, in that sentence, is a property. An institution holds it and platforms distribute it. Nobody is being invited to test anything. People who talk like this in public generally think like this in private.

So who does own the science? The person who made the discovery? The laboratory that paid for it? The company that patented it? The journal that judged it fit to print? The minister who chose the advisory panel? The television editor who cut a three-hour argument to eight alarming seconds? Or the audience, which is expected to obey the conclusion without ever being shown the premises?

None of this is a conspiracy theory. It is bookkeeping. Follow the money, the appointments, the permissions, the headlines, and the career risks, and you find that nobody owns ``science'' in the singular. But many people own pieces of the machine that decides which science is seen, which scientist is heard, and which doubt gets filed under heresy.

The argument of this chapter is that public science is two things at once. It is a method for finding things out. It is also an institution for handing out attention, prestige, risk, and permission. The method can be open in principle while the institution is shut in practice. A paper can be public while its data are locked away. A debate can be formally free while one side has a grant, a press office, and the minister's phone number, and the other has a rented room and a damaged name.

This does not make scientific knowledge worthless. It makes the social life of knowledge everybody's business. If we refuse to look at that life, the word ``science'' is left lying around for anyone who needs a technical-sounding blessing.

\section{Five questions of ownership}

We usually think of ownership as a deed in a drawer. The ownership of science is stranger, and it is exercised without any deed at all. It is control of the pipeline: the power to decide which questions get funded, which instruments get built, which results get printed, which experts get invited, and which conclusions become rules.

The legal tangles are real enough. A university claims the intellectual property in research the public paid for. A publisher sells access to articles whose authors and referees worked for nothing. The author keeps the glory and loses the right to share his own paper. Physicists have been arguing for years about who ought to own a scientific article \cite{blume-aps-copyright}.

But the ownership that matters most lies upstream. A sponsor who can choose the question has no need to tamper with the answer. ``Can this treatment work?'' is a different question from ``For whom does it work, at what cost, and compared with what?'' ``Can this system predict behavior?'' is a different question from ``Should it be allowed to judge people?'' Whoever picks the first wording has made a political decision before the experiment starts. The researcher can then be perfectly neutral about the result. He is working inside a frame that somebody else chose, which fixed the subject, the time scale, the population, and the definition of success. The work can be flawless and still answer a question that serves another man's purposes.

The public meets science at the far end of this pipe, after the mess has been laundered into a sentence. Researchers ``prove.'' Science ``demands.'' Independent experts ``agree.'' Nobody shows the audience the questions that were turned down, the methods that were dropped, the negative results, or the names of the people who were never invited.

So the ownership question comes in five parts:

\begin{enumerate}
\item Who owns the question?
\item Who owns the instruments, the data, and the computing power?
\item Who owns the channels of publication and public attention?
\item Who owns the authority to turn evidence into rules?
\item Who pays when the forecast fails?
\end{enumerate}

The fifth is the one that never comes up at the press conference.

\section{The casting call}

Every crisis needs a cast. There is the Expert, whose authority must look good on camera. There is the Modeler, whose equations are asked to do the work of prophecy. There is the Dissenter, who is useful as a villain even when he is right. There is the Responsible Citizen, who nods on cue. And there is the public, which might have been a jury and is treated as a crowd to be managed.

Nobody announces the audition. It happens through invitations, contracts, funding priorities, editorial taste, and the quiet arithmetic of promotion. The scientist who says ``we do not yet know'' may be entirely honest, and he is useless on television. The one who can produce a single number before the evening news has a bright future. In an emergency, what counts is seldom knowing more. It is looking as though you know enough.

This is how a method turns into a pantomime. Science is supposed to expose its errors in public. The show needs them hidden until after the applause. Science is supposed to keep a measurement apart from a guess. The show needs a clean line, a clear villain, and a policy that can be announced before the footnotes are typeset.

What comes out need not be false. It is something more slippery: truth that has been selected, trimmed, and branded, and then released at the moment most convenient to power. Edward Bernays, who wrote the manual a century ago, never supposed that propaganda required lies. It requires a managed environment in which some truths arrive with microphones and others arrive after the audience has gone home \cite{bernais-propaganda}.

The scientist who becomes a public character also acquires a new master, and the master is the character. A cautious sentence is good science and bad television. A correction is honest and looks like weakness. Bit by bit he stops defending a proposition and starts defending a role.

Then the institution casts around him. The same faces are invited because they are known, and they are known because they are invited. Their familiarity is taken for consensus. Nobody has measured that consensus. It has been staged.

\section{The investor's microscope}

Money does not have to dictate results in order to shape a field. It only has to decide which questions can be afforded. A grant makes one hypothesis urgent and another eccentric. A government program produces a generation of specialists who are superb at the question that was funded and strangers to the one that was not.

None of this requires villains. Researchers need salaries, instruments, data, assistants, and above all the right to fail. Whoever controls those things also acquires a quieter power, which is to decide what failure may mean.

In a healthy scientific culture a failed prediction is a useful bruise. In a managed one it is a fire in the reputation department. The institution does not ask what was learned. It asks who can be blamed without endangering the budget. The failed model is quietly retired and the forecast is reworded, and the press release recalls only that the experts acted decisively.

The public is then offered a peculiar bargain. It must trust the system because the system is expert, and it must not inspect the system, because inspection is anti-expert. Ask who paid, who chose, who gained, or who was left out, and you are accused of attacking knowledge itself. A question about governance is answered with a lecture about facts. A question about values is buried under a chart.

But no chart can tell us whether a sacrifice is just. No model can decide whose liberty is expendable. No instrument measures the right balance between risk and dignity. These are political judgments. When they are smuggled into a technical recommendation, politics puts on the white coat and science is left holding the bill.

The arrangement is at its worst where a public agency and a private vendor merge into one working system. The agency has legitimacy and no capacity. The vendor has capacity and no legitimacy. Responsibility then travels like the baton in a relay. The vendor only supplied a tool, and the agency only followed technical advice. The citizen is left trying to sue a cloud.

Publishing runs on the same logic. Researchers are scored by counting their papers, citations, grants, patents, and media mentions. A system built to reward discovery ends up rewarding discoverability. The paper that survives is the one the metric can read. Everything else becomes a hobby.

\section{Peer review: gate, filter, oracle}

Peer review is sold to the public as the moment when science protects everyone from error. It is better understood as a filter, with real virtues and well-known faults. It catches the missing control and the careless calculation. It also reproduces fashion, polices the borders between disciplines, and rewards the manuscript that sounds like the last one. The evidence that it does these things is decades old \cite{1997-armstrong}.

Reviewers are not wicked. They are busy and unpaid, and they come from the same small world as the authors. They know which claims are prestigious, which methods are in season, and which names earn a premium or a discount. The system has a prior, and nobody declares it.

In principle a young researcher with a strong result can overturn a famous theory. In practice he must first persuade the people who built their careers on the old vocabulary. A paper can be rejected because it is wrong. It can also be rejected because it is awkward, impolite, unfashionable, or hard to classify.

The gate also controls memory. What gets published is cited, taught, funded, and repeated, and what gets rejected becomes an anecdote. That is how a field comes to look more settled than it is. The reader of the literature sees only the survivors.

Then comes the theater. ``Peer reviewed'' is used as though it meant ``true.'' All it means is that two or three people thought a manuscript fit to publish under the conditions of the day. It does not mean the result is correct, or important, or fit to govern millions of strangers.

Once that is forgotten, later criticism looks like an attack on science. It is nothing of the sort. Criticism is the second half of review, the half that cannot be scheduled and cannot be ticked off on a form.

\section{Models that take office}

A model is a disciplined act of imagination. It keeps some features of the world and drops others, and that is where its power comes from. A model that left nothing out would be the world over again and of no use to anybody.

The trouble begins when usefulness is promoted to authority. First the model is a map. Then it is a forecast. Then it is a target. Finally the target is treated as a fact, the institution rearranges its conduct to hit it, and the prediction is declared confirmed. A population forecast shapes a policy that changes the population. A risk score decides who gets attention, and the attention changes the measured risk. A school ranking steers the money, and the ranking is then presented as an objective account of quality. By this point the model is a player in the game it was built to describe.

Watch how the words move. ``Assume'' becomes ``estimate.'' ``Estimate'' becomes ``projection.'' ``Projection'' becomes ``expected outcome.'' And ``expected outcome'' becomes grounds for coercion. Somewhere along the way the adjective ``model-based'' drops out, usually at the door of the room where the power is. But information that was missing at the start is still missing at the end \cite{brillouin2}. A dashboard built around the gap does not fill it.

The same reversal runs through every institution that measures itself. The old phrase is the tail wagging the dog. Universities invent rankings and then rebuild themselves to climb them. Journals invent impact factors, and authors bend their titles to attract citations. A measure becomes a target and stops measuring anything. The dog is now wagged by the dashboard and by the bean-counters, the \emph{Erbsenz\"ahler}, who convert life into columns, compliance rates, and quarterly variance. The dashboard tells them what can be counted, and they tell everyone else what counts. Between them a human being becomes a cost center with a pulse. The tail has gone beyond wagging the dog. It sends the dog an invoice and calls the invoice evidence.

All of it looks voluntary, which is what makes it so effective. Nobody is forced to chase the metric. People are only encouraged to do what it rewards. The researcher who refuses keeps her integrity and loses her job. The system can say with a clear conscience that nobody was coerced, because the coercion was spread across a thousand perfectly polite decisions.

So a society learns to edit itself in advance. People anticipate the ranking, the grant panel, the moderator, the news cycle, and the future employer. They leave out the sentence that might be misread. They choose the project that can be counted.

This is one way a society can decay without a coup. No dictator has to forbid a thought. It is enough that some thoughts are expensive and others pay well. The population does the censoring itself, and then congratulates itself on its freedom.

The answer is not to throw the models out. Without them we have only anecdote and instinct, and those belong to whoever tells the most confident story. The answer is to be more demanding. Show the model. State its assumptions. Publish the sensitivity analysis and the plausible alternatives. Say what observation would prove it wrong. A model that cannot survive that treatment has no business being used to punish people.

\section{Two kinds of doubt}

There are two kinds of skepticism, and power would like us to confuse them.

The first is built into science. Show me the method. Show me the controls. Tell me what would change your mind. Let another team try to get the same result. This is the immune system of science, and there is no science without it.

The second says that every fact is an opinion, every measurement a swindle, and every expert the servant of a hidden master. It liberates nobody. It leaves people defenseless before the loudest advertiser, the richest lobby, and the most confident fraud.

Authority profits when the two are run together. If all doubt is irrational, nobody may check anything. If all knowledge is corrupt, there is no point in checking. Either way the population obeys, the institution in one case and, in the other, whichever influencer claims to have escaped it.

The remedy lies between blind trust and theatrical disbelief, and it is organized questioning. Who collected the data? What was not measured? Which costs were counted, and which were left for someone else to bear? What happens if the prediction fails? Who has the power to correct the record, and who profits while it stays uncorrected?

Those questions endanger only the systems that cannot survive them.

Disagreement among scientists is no sign of corruption. Fields have live disputes because evidence is incomplete and concepts are contested. The demand for unanimity is itself anti-scientific. It asks researchers to hide the unsettled edges of their subject so that the public can be handed one obedient conclusion.

And disagreement can be strangled while the outward forms of debate are kept up. A paper is reviewed forever. A dissenter is called ``not constructive.'' A result is left out of the meta-analysis because it uses the wrong vocabulary. A question about evidence is answered with a remark about character. Faking data is only one kind of scientific misconduct. Biased refereeing, smears, and the policing of approved opinion do their damage as well \cite{Cabbolet-2013}.

\section{``This is physics''}

We picture propaganda as a loudspeaker bellowing an obvious lie. The successful kind is quieter. It arranges words and pictures so that one conclusion feels like common sense before anyone has argued for it.

The recipe is short. Name the danger. Show a number. Put a clock in the background. Add a person in a uniform or a lab coat. Announce that responsible people agree. Describe objections as confusion, selfishness, or menace. Offer obedience as the only grown-up emotion.

No fact needs to be invented. The facts only need to be arranged into a moral climate in which any alternative looks indecent. The public is not told that a policy is one option among several. It is told that reality has issued an order.

Language is where the ownership is exercised. Whoever defines ``misinformation'' defines the limits of what may be considered real. Certainly some statements are false and some are meant to deceive. But a filter is not neutral merely because much of what it catches deserves catching. Somebody sets the threshold, the penalty, and the appeal procedure, and somebody appoints the judge.

``Follow the science'' is the perfect specimen. Science can describe relationships, estimate probabilities, and trace consequences. It cannot walk ahead of a society and choose where the society should go. Follow which finding, under which assumptions, at whose expense, and toward whose preferred future?

I have a small exhibit, a short video of a communications consultant \cite{CommunicationsConsultantCO2Video}. Faced with an objection about carbon dioxide, he explains that this is not his opinion, nor his opponent's, nor that of seven people outside the door. It is \emph{physics}. Picture a rod of gold and a rod of iron held in a fire. The gold, he says, gets hotter, and gets hot faster; mix the two metals, and the final heat depends on how much gold is in the mix. Gases are just the same: ``CO$_2$ gets hotter than oxygen. That is just how it is. It does not matter whether you or I believe it.''

As rhetoric it is splendid. It has all the gestures of a school demonstration: two rods, a fire, a mixture, and at the end a falling ball. As physics it is a wreck. How fast a metal object heats depends on its mass, shape, and surface, on how the heat reaches it, and on its heat capacity. Gold does have a lower specific heat per unit mass than iron, so equal energy put into equal masses raises the gold's temperature more. That is a long way from a law that gold in a fire ``gets hotter.'' Leave both rods in the same fire and both will approach the temperature of the fire. And an alloy does not take its behavior as a kind of moral average from the percentage of the nobler metal.

The gas version is worse. In a parcel of air near local thermodynamic equilibrium, the carbon-dioxide molecules do not run a private fever while the oxygen stays cool. They share one temperature. The greenhouse effect is a matter of radiation. CO$_2$ absorbs and emits infrared in wavelength bands where oxygen and nitrogen hardly interact at all. Adding CO$_2$ raises the altitude, and so lowers the temperature, at which part of the planet's heat escapes to space. The surface and troposphere then warm until the energy budget balances again. One can refine this with overlapping absorption bands, convection, water vapor, clouds, and lapse rates \cite{IPCCWGI2021TechnicalSummary}. No refinement turns carbon dioxide into a gas that is simply ``hotter than oxygen.''

As for the falling ball, it proves something true and irrelevant. Gravity does not take a vote. But the fact that reality ignores our beliefs does not certify every sentence that begins ``this is physics.'' You cannot promote a muddle to a law of nature by standing it next to gravity. You are only borrowing gravity's good name after your own explanation has failed.

The case is worth the space because it forces us to keep two things apart. More carbon dioxide in the atmosphere does contribute to warming. The consultant's account of why is nonsense. A critic who demolishes the gold-rod analogy has not refuted the greenhouse effect. A believer who invokes the real greenhouse effect has not rescued the analogy. Sentimental anti-realism lives in exactly this gap, where loyalty to the right conclusion excuses any amount of indifference to how one got there. In that gap the word ``science'' has stopped being an invitation to inspect the mechanism. It has become an order to stop inspecting.

Declining to ask may look like humility, but it is surrender.

\section{The emergency that never leaves}

The decisive step is always described as temporary. Just for now we will suspend the usual procedure, narrow the discussion, trust the approved voices, and postpone the audit. The situation is exceptional. It is dangerous.

Then the emergency finds that it is comfortable.

Once a society has built the database, the surveillance channel, the advisory committee, the procurement habit, and the censorship reflex, taking them apart is harder than building them was. Each layer is defended by the one before. The exception becomes a precedent. The precedent becomes infrastructure, and the infrastructure becomes common sense. Anyone who objects is told the machine is too big to stop.

Often that is true. It is history's most depressing trick: yesterday's decisions turn up today disguised as facts of nature. We did not choose the platform. It is simply there. We did not choose the dependency, which now goes by the name of efficiency. We did not choose the narrowing of debate, which is now called responsibility.

By the time the public sees where the road leads, the road has been privatized, the signposts are sponsored, and the return lane has been taken out. In the clean, heroic sense it is already too late. There will be no lever marked ``restore society.'' There will be only expensive friction: disclosure, repair, legal limits, independent replication, and the unglamorous labor of saying no after the money has been spent.

``Too late'' is not a date on the calendar. It is a condition. It has arrived when everybody can see the damage and the people able to stop it would first have to admit that they profited from starting it.

The permanent emergency lives on two stories. In the first, nobody could have foreseen the danger. In the second, nobody can now afford to undo the response. Between the two the real question goes missing: who made the choices that turned a temporary measure into a permanent power?

There is a further thing the emergency state forgets, which is that the patient is awake. Policy changes what it measures. A population is no sample in a jar. It adapts, resists, hides, and learns the rules of the test. A risk label changes the conduct of the people labeled. A policing algorithm decides where the police go, and so writes its own next batch of data. When an institution calls its intervention ``evidence-based,'' the intervention is also producing the evidence. It is making a new world and calling the result a measurement.

That is why consent matters outside the clinic. Whole populations can be enrolled, without signing anything, in experiments on their movement, work, schooling, speech, and credit. The technical language hides what is being done, which is a redistribution of risk. One group receives safety and another surveillance. A third loses its income, and a fourth is told to wait. The justification is statistical, but the distribution is political.

If citizens are to be experimental material, they are owed what any laboratory subject is owed: a stated objective, a public protocol, an independent monitor, a stopping rule, and an honest account of harms. What they are generally told is that the intervention is too complex to evaluate and too important to question. People who say this like to call it humility before complexity, when in fact it gives them immunity from accountability.

\section{The dignity of ``we do not know''}

It takes a certain courage to say ``we do not know'' to an audience that wants a command. The sentence solves nothing. It only stops an institution from passing off its appetite for control as knowledge.

Admittedly the sentence can be abused. It can mean the evidence really is incomplete. It can also mean the speaker has not looked, or is hiding an interest, or has learned to live off ignorance as a permanent excuse. So uncertainty should be itemized. What exactly is unknown? What is merely disputed? What would settle it, and how long would that take? What can safely be done in the meantime? And which actions can be undone?

That last question belongs on the front page of every policy paper. Where knowledge is uncertain, reversibility is a moral virtue. If we must act before we know enough, we should choose actions that leave our successors room to correct us. The opposite rule is to make the change irreversible while the evidence is still provisional, and it is the mark of a system that expects to own the future.

Institutions prefer certainty because certainty travels well. It fits a headline, a budget request, and a chain of command. Uncertainty calls for judgment, and judgment needs someone who can be held responsible and questioned. So the uncertainty is squeezed out before it reaches the public. The scientist gives an estimate. The administrator turns it into a threshold. The politician turns the threshold into a rule, and the rule becomes a moral test. The citizen on the wrong side of it has stopped being a person in a difficult situation and become a deviation from the model.

A forecast is not a promise, a probability is not a fate, and an error bar is not an embarrassment to be deleted before the minister speaks. Casti's book on what science can know about the future treats prediction as a problem with limits and makes no oracle of it \cite{casti:92a}. Feyerabend made the wider point that no single procedural formula captures how inquiry really works \cite{Feyerabend2001May}. His ruder question was how a society should defend itself against science once science is treated as beyond question, and it is the question of this chapter \cite{feyerabend-defense}.

The people who demand certainty are seldom the ones who will pay for the mistake. They want the authority of a prediction and the immunity of a forecast. When it fails, the fault lies with complexity, with human behavior, with the unforeseeable, or with the public's lack of obedience. The model is never wrong. Reality is merely difficult.

An honest institution would publish every forecast together with its conditions of failure and a named date for review. It would say who gains if the forecast is believed and who loses if it is wrong. It would also keep its records in a form that can be contested, including the assumptions, the dissenting reports, the abandoned alternatives, and each change of guidance. Without that, every generation inherits a polished myth in which the present was inevitable, and it concludes that the future must be inevitable too.

\section{The new priesthood}

The modern priesthood wears no robes. It wears lanyards. Its sacred objects are the dashboard, the simulation, the benchmark, and the confidence interval. Its liturgy is the briefing, and its incense is the glow of the projector.

I do not say this to insult expertise, which is real and valuable and painfully acquired. The danger starts when expertise hardens into a rank that nobody outside the guild may question. A technician who knows how a system works is not thereby entitled to decide what the system should do to a population.

The priesthood is seductive because it offers to take the politics out of political choices. Nobody has to argue about values, because one can ask what the model says. Nobody has to defend the distribution of costs, because the algorithm has ranked the options. And nobody has to admit that the future is open, because a large enough dataset is supposed to have closed it.

Artificial intelligence makes the temptation stronger. Language models produce fluent explanations, but fluency is not evidence and scale is not understanding. Even the research on their supposedly ``emergent'' abilities shows how far definitions and choices of measurement determine what appears to emerge \cite{Wei2022Jun}. Wrap a machine's output in enough technical glamor and people forget to ask the oldest question there is: how do we know?

The machine does not have to be conscious to hold power. It only has to be installed in hiring, lending, schooling, policing, or medicine. Its categories then become social facts. A person is refused a chance because a score was low, and no human being can say what the score means or how to appeal it. This is the priestly miracle in its modern form. A political decision becomes a technical output, and the technical output becomes fate.

A society does not decay because its people grow stupid. It decays when they are trained to be spectators of the decisions that govern them. They watch the experts argue, the numbers climb, and the policy arrive, and they find that their own part is to comply. The vocabulary changes first. ``Citizen'' becomes ``user,'' ``patient'' becomes ``case,'' and ``neighbor'' becomes ``risk.'' ``Question'' becomes ``misinformation.'' The change is more than cosmetic. Once a person is a data point, he can be optimized. Once dissent is contamination, the cleaning crew can be thanked.

And the more complicated the system grows, the more childish its public theater becomes. A civilization that can sequence genomes still offers its citizens two cartoon buttons: trust the experts, or trust nobody. The adult position is conditional trust, open evidence, reversible policy, and consequences for those who get it wrong. It does not fit on a banner, so the banner wins.

The spectator society is not quiet. It is noisy, furious, and permanently online. But noise is not participation. A million reactions can be collected, sorted, sold, and forgotten without one citizen altering one decision. That is the political genius of managed attention: it turns disagreement into content. The angry citizen becomes an engagement statistic and the scandal a temporary spike. The system can absorb any amount of criticism, so long as criticism never acquires a procedure or a hand on the wheel.

\section{Take back the question}

A science fit for a free society has to be more than accurate. It has to be open to inspection, contest, and correction. Nobody is asking citizens to repeat particle-physics experiments in the kitchen. The demand is that institutions make their authority legible.

That means a few plain things. Any recommendation with consequences should state its question, sponsor, assumptions, data, uncertainty, alternatives, winners and losers, and date of review. It should mark where the technical finding ends and the value judgment begins. Advisory committees should publish their minority reports and not treat disagreement as a public-relations accident. The public should know who was invited and who was not. Data and code should be open wherever privacy allows, and where they cannot be opened, an independent auditor should confirm the reason. Replication, maintenance, negative results, and corrections should be rewarded. It is no scandal that researchers respond to incentives. The scandal is that the incentives are discussed as if they were weather.

And every policy needs a brake: a stopping rule, a route of appeal, and an expiry date that only a public process can extend. The brake will be inconvenient. That is what a brake is for.

None of this dethrones science. There is no throne in the laboratory and there should be none in government. The purpose is to stop the word ``science'' from serving as a password that opens every door and shuts every mouth.

Nor does any of it promise a return to innocence. What has been built has been built, and it has left habits, dependencies, fortunes, and scars. The future cannot be reset. It can only be contested from inside the wreckage of the present, which is less satisfying than revolution and more useful than nostalgia.

The limits of knowledge are the best argument there is for dividing power. If nobody can know everything, nobody should govern everything on the strength of one model. If every measurement is partial, institutions should be plural. If the future is uncertain, the people who will live in it should keep the right to revise what was decided in their name.

Who owns the science? Whoever shapes its questions, its funding, its publication, its translation into rules, and its enforcement owns a piece of it. Who casts the scientists? Whoever controls attention, access, and the price of being wrong.

The first act of resistance is embarrassingly simple, and it is to refuse the script. Ask who wrote the question before you argue about the answer. Ask who gains from certainty. Ask what the experts were not allowed to say. Ask who gets to declare the experiment finished while the population is still living inside it.

By now it may be too late to stop the machine from existing. It is not too late to decide whether the machine may call itself wisdom.


% ===================== Chapter file ch04 =====================
\chapter{The Last Permissible Thought}
\label{chap:last-permissible-thought}

\begin{flushleft}
\large\itshape Truth, bullshit, tribal belief, and the moment when a society can no longer hear a warning in time
\end{flushleft}
\vspace{1.2em}

\noindent\textbf{Truth is elusive, not illusive.} What is elusive is hard to catch. What is illusive is a manufactured appearance. Reality is mostly the first kind of thing. Public reality is more and more the second.

Philosophers used to ask what is true and what exists. The question of our time is more nervous and far more revealing: \emph{under what conditions would I be willing to admit that something is true?} This is no seminar-room refinement. It is the central political fact of a society in decline. A bridge is unsound before the ministry accepts the report. A bank is insolvent before the regulator allows the word to be used. A policy is failing before the people who run it, live off it, report on it, or suffer from it can afford to say so. Reality arrives first. Recognition follows. Action limps in last.

Sometimes the delay does no harm, and sometimes it kills. A civilization becomes ``too late'' when the time a true description needs to become believable is longer than the time left to act on it. Truth is still there. It has only lost the race. That is the thesis of this chapter. Too late is no date on an apocalyptic calendar. It is a relation between three clocks:

\begin{enumerate}
  \item the clock of the underlying process;
  \item the clock of social recognition;
  \item the clock of institutional response.
\end{enumerate}

When the first outruns the other two, facts become posthumous. They are admitted once admitting them has become cheap. The dissenter is vindicated at the funeral. The inquiry is commissioned after the collapse, and its report is written in the passive voice so that nobody still on the payroll need feel accused.

A statement does not enter public life merely by being true. It needs a speaker with standing, words that are still allowed, an audience that can bear the conclusion, and an institution able to act on it without destroying itself. Take away any one of these and the fact remains a fact and changes nothing. Reality does not need a press office. Recognition does.

\section{Reality is not a vote, but belief is}

Public argument constantly mixes up four things:

\begin{enumerate}
  \item what is the case;
  \item what someone believes is the case;
  \item what someone can justify believing;
  \item what someone is prepared to say in public.
\end{enumerate}

A mountain on the far side of an unexplored valley does not begin to exist when it is mapped. A fraud is real before the court says so. But the map and the verdict decide when the rest of us may responsibly treat the matter as settled. The world supplies the facts. We supply the terms of admission.

This sounds elementary until incentives come into it. Then every level gets its doorkeepers. Scientists decide which anomalies deserve a look. Editors decide what is reportable. Courts decide what is admissible. Platforms decide which label a claim must wear. Employers decide which sentences are compatible with ``our values.'' And each citizen decides which facts would cost too much to believe. What the public calls ``truth'' is therefore usually a parcel: a fact, plus an approved method, plus an approved witness, plus an approved tone of voice.

The philosophers offer several tests of truth, and each fails when asked to do the whole job. Correspondence to reality, without humility, becomes the shout ``I see things as they are.'' Coherence without contact with the world is an elegant delusion. Pragmatism without limits turns ``useful'' into ``true,'' which is the favorite move of propagandists and managers. And social construction with nothing left over is a palace coup in which talk declares itself sovereign over gravity, infection, arithmetic, and death. The grown-up position is less glamorous. There is a reality independent of our descriptions, and our access to it is indirect, fallible, bent by interest, and open to correction.

Scientific institutions were built on that uncomfortable compromise. Popper insisted that evidence must be able to wound a theory \cite{popper-en}. Kuhn showed that no scientist inspects the world from nowhere; he works inside a paradigm that tells him what counts as a problem and what counts as a solution \cite{kuhn}. Each needs the other. Without Popper, paradigms turn into fortresses. Without Kuhn, philosophers pretend that evidence arrives unformatted.

Nor is being right the same as knowing. Gettier's famous cases showed that a belief can be justified and true and still true by luck \cite{Gettier1963}. The point travels well beyond philosophy. An investor predicts a crash for idiotic reasons and gets rich. A pundit repeats fifty rumors, hits once, and sells himself as a seer. A government picks the right policy on fraudulent evidence. The outcome does not vouch for the route.

And consensus settles nothing by its mere existence. Agreement is evidence when it comes out of independent inquiry, different methods, hostile testing, and a real reward for whoever refutes it. It is much weaker evidence when everybody shares the same sources, training, patrons, and penalties. Ten thermometers agree for a different reason than ten courtiers do. To book both as ``ten independent confirmations'' is epistemic money laundering.

So ``Is there a consensus?'' is the wrong question. Ask what produced the consensus, what could overturn it, and what happens to the person who tries. A consensus that no imaginable finding could defeat is not knowledge. It is jurisdiction.

\section{When would you believe it?}

Ask people whether they care about truth and they all say yes. The answer is worthless. The question that reveals something is conditional: \emph{what observation, from what source, by what procedure, would change your mind?} If there is no answer, the belief is not being held as a hypothesis. It is being held as property.

Belief has several thresholds, and they need not line up. What must I see? How much evidence is enough? Whose word counts? What would accepting this say about me and my people? How sure must I be before I pay a price?

A physician may privately suspect that a favored treatment does nothing. She needs more evidence before she changes her practice, and more again before she will criticize colleagues in public. An engineer sees the corrosion and has no authority to close the bridge. He knows the mayor will demand absolute certainty, because a closure is visible and a collapse is hypothetical. A citizen believes a policy is failing and goes on defending it, because to admit the failure would hand a victory to people he despises.

Now turn it round. We lower the bar for claims we like. A rumor about the other tribe is ``plausible enough to share.'' A documented scandal in our own requires a court verdict, three replications, and the signature of God. The same mind that lectures others about evidence becomes a brilliant defense lawyer on its own behalf.

Psychologists call this motivated reasoning. It is seldom a crude decision to lie to oneself. It is a selective recruitment of memories, standards, and counterarguments that goes on until the preferred conclusion feels earned \cite{Kunda1990}. We do not update a proposition only. We update a self, a set of loyalties, a career, and next Sunday's argument at the dinner table.

You can watch it anywhere. The manager calls rising numbers ``the trend'' and falling ones ``noise.'' The activist treats a vivid anecdote as proof when it helps and as an outlier when it hurts. The professor admires rigor in hostile papers and ``theoretical importance'' in friendly ones. The voter sees prudence when his own man compromises and corruption when the other man does. The believer credits a fulfilled prophecy to the doctrine and an unfulfilled one to faulty interpretation.

None of this means that all beliefs are equally baseless. That slogan is only another way of avoiding judgment. It means that the standards by which we sort strong beliefs from weak ones can be captured by the very beliefs they are supposed to judge. The customs officer is on the smuggler's payroll.

There is a simple discipline. Before you defend a belief that matters, write down the strongest observation that would count against it; a source outside your own camp whom you would accept; how confident you really are, as opposed to how confident you sound; what you gain, socially or morally, by going on believing it; and the date on which you will look again. If you cannot fill in the first two lines, you do not have a belief that responds to evidence. You have an identity arrangement.

Timing changes the standard of proof. A reservoir wall shows a crack. One inspection may be a false alarm, and it is sensible to wait for a second. But the water is rising, and every day of delay costs more. To demand courtroom certainty before taking a precaution that could be undone is no kind of rigor. It is the official's preference for being precisely blameless over being approximately alive.

Sometimes belief is part of the fact. In a bank run the depositor asks whether the bank is sound, and also whether other people think it is sound. A false rumor can bring about the failure it predicts. A true reassurance can fail because nobody expects anybody else to trust it. Institutions can die of public knowledge about public knowledge.

\section{Six laboratories of recognition}

The gap between truth and recognition is clearest where error has a price. Here are six cases. In each, somebody fails to believe what is already so, for a different reason.

\subsection{The diagnosis}

A patient gets an alarming test result. A careful physician asks about false positives, base rates, and the quality of the sample, and may order a second test. Hesitation at this stage is method, and nobody should call it denial.

Now suppose the second test agrees, the symptoms fit, and every week of waiting closes off a treatment. The patient wants a fifth opinion, because the first four were ``not holistic enough.'' The standard of proof is now moving in response to the conclusion. He has stopped asking how to tell disease from error. He is asking how to keep the old world alive another week.

The test for an institution is the same. Was the standard of proof fixed before the evidence came in? If it rises every time the unwelcome conclusion gets nearer, skepticism has turned into a drawbridge.

\subsection{The anomaly}

A mature research program meets a result it cannot explain. One anomaly should not overturn it; instruments fail, samples get contaminated, and chance produces monsters. Scientific conservatism protects knowledge from fashion. It can equally protect a paradigm from reality.

What matters is whether the anomaly can earn promotion, from nuisance to something worth replicating, and from there to a recognized problem that may force a revision. A living institution has a ladder for unwelcome data. A dead one has only a waste chute.

Resistance is fiercest where the evidence carries a moral charge. Semmelweis is remembered because hand disinfection implied something unbearable, that respectable physicians had been killing their own patients \cite{semmelweis2024}. The lesson is not that the lone genius is always right. Most lone claimants are wrong. The lesson is that a profession needs procedures that can tell a humiliating truth from a flattering error, and can do it before the funerals have run the experiment.

\subsection{The failing project}

A large public project misses a milestone. Leadership explains that ambitious programs meet delays, which may be true. Then the second milestone is missed, the budget grows, the scope shrinks, and a new dashboard reports ``percentage of workstreams initiated.'' When is the project failing?

No bell rings. Failure can be recognized only against a baseline and a rule laid down in advance. Without them every result can be rescued by narrative. Costs become investments, delay becomes learning, and cancellation becomes strategic renewal.

The use of euphemism is the lesser scandal. The greater one is that the organization may lose the very ability to describe failure in a form that travels upward. Each layer translates bad news into the dialect the next layer requires. By the time it reaches the board it has been washed, ironed, and folded into reassurance.

That is why ``speak truth to power'' is not enough. Power has to keep a receiver switched on. Where the bearer of bad news is promoted for loyalty, the news will be tailored.

\subsection{The intelligence estimate}

Before a war, analysts hold fragments. Some sources are reliable and incomplete, and others have reasons to exaggerate. The leaders face lopsided errors. To underrate a threat may be catastrophic, and to overrate it may license a catastrophe of their own making.

``Follow the evidence'' gives no help, because evidence does not say which mistake to fear more. Somebody has to state that openly and say which actions can be reversed. If the uncertainty is hidden in order to build support, the policy gains force for a season and the institution loses its credit for a generation.

Afterwards hindsight rearranges the file. Warnings are remembered as certainties, and the caveats vanish. An honest review reconstructs what was known at each moment and grades the reasoning against that. A political review asks whose career has to survive. The first studies error. The second hands out innocence.

\subsection{The slow variable}

Fast events get the attention, and slow variables decide the outcome. Maintenance backlogs, exhausted soil, an aging population, lost skills, public debt, sinking groundwater, and the silting-up of the law can all worsen for decades while daily life goes on as usual.

Slow decline is easy to argue away. Each year's change is small enough to dispute. The measurements get revised, and a good quarter is advertised as the turnaround. And since responsibility crosses electoral terms, every incumbent inherits the causes and bequeaths the consequences.

By the time the trend cannot be denied, the stock has changed. The missing engineers cannot be trained overnight. The children who were not born cannot be legislated into existence. The bridge now needs money that three other emergencies have already claimed.

This is the purest case of ``already too late.'' The numbers were public for years. What was lacking was a coalition willing to trade present comfort for a future benefit that nobody could see. There was no conspiracy to hide the figures. Everybody hid behind the calendar.

\subsection{The intimate fact}

Finally, a failing marriage. One partner has known for a long time that the trust is gone. The other senses something and explains away each sign separately. The friends say nothing, for fear of causing the break they are only observing. Everybody is considerate, and all that consideration is a machine for delay.

When it is finally said, the cry is ``Why did nobody tell me?'' But the warnings were there, in forms that could be dismissed. More evidence would not have helped. What was needed was the willingness to live in a more painful story.

I include the private case because societies are made of the same nervous systems. Nations postpone grief as well. Institutions mistake silence for loyalty, and citizens prefer a bad world they know to an alternative whose price is unknown. The scale is different and the ego is the same.

In all six laboratories the test is identical. Is the bar for belief set by the quality of the evidence, or does it move to protect the believer from the cost of the conclusion? The first is skepticism. The second is denial with stationery.

\section{Bullshit}

Harry Frankfurt drew the indispensable distinction \cite{Frankfurt-OnBullshit}. The liar and the honest man are opponents, but they play on the same field. Both care where the truth lies, the one to state it and the other to hide it. The bullshitter has left the field. Whether his words match reality is a matter of indifference to him. He cares about the impression.

That makes bullshit more corrosive than lying. A lie can be caught by holding it against a fact. Bullshit wears away the expectation that the comparison matters. It trains an audience to take language as mood and as a badge of membership. People stop asking whether a statement is accurate. They ask whose side it is on and whether repeating it will do them good.

Modern organizations are superb at producing it, because they demand words faster than reality supplies knowledge. A crisis requires a statement before the facts are in. A strategy retreat requires a vision before anyone has admitted the constraints. A grant application requires ``transformative impact'' before the experiment exists. A ministry requires measurable targets, so a purpose that cannot be measured is translated into numbers that can be reported and gamed. The resulting prose is not exactly false. It is indifferent to truth and dressed in truth's uniform, which is worse.

The industrial product is easy to recognize. Its abstractions cannot be proved wrong by anything. Moral adjectives stand where a causal mechanism should be. Statistics come without denominators. Claims are ``evidence-based,'' and no evidence is named. The passive voice removes whoever did the deed. Promises are judged by their intentions and failures are excused by complexity. And urgency suspends scrutiny until, later on, forgetfulness makes an audit impossible.

Bullshit is cheaper than lying, since you need not keep track of the truth. It is safer than candor, since it can always be denied. Challenge it and the speaker retreats from assertion to aspiration: ``That was not a factual claim. It expressed our commitment.'' The sentence was built with a fire exit.

Repetition does the rest. A statement heard often feels truer, simply because it has become easier to process \cite{HasherGoldsteinToppino1977}, and knowing better is not a reliable defense \cite{FazioEtAl2015}. Familiarity turns up in the stolen clothes of corroboration. A slogan met in ten places may have started in one press release. The nervous system registers fluency and does not ask about origin.

So let us stop flattering public communication with the name ``marketplace of ideas.'' Much of it is a laundromat for sentences. Claims go in stained with interest and come out smelling of consensus.

\section{The belief as mortgage}

Nobody meets a fact as naked data. We invest in what we take for granted. A belief acquired early becomes a fixture. A belief defended in public becomes a mortgage. A belief used to condemn other people becomes a hostage. The more moral and reputational capital is tied up in it, the harder it is to revise.

Festinger called the underlying pressure cognitive dissonance \cite{Festinger1957}. The pressure can be relieved without giving up the false belief. One can reinterpret the evidence, blacken the witness, extend the deadline, discover sabotage, or announce that success in the literal sense was never the aim. Reality supplies the discomfort, and imagination arranges the refinancing.

The investment has four layers.

\emph{Autobiography.} ``I believe this'' quietly means ``a sensible person with my history believes this.'' Contrary evidence therefore attacks a whole life story. The retired official asked to reconsider a policy is being asked to reinterpret thirty years of service. The scholar whose theory has failed must demote his own publications from achievements to curiosities. The parent told that a decision did harm hears an accusation against love itself. No wonder the mind hires lawyers.

\emph{Competence.} The higher the office, the more people have built their work on its occupant's premises. Error then becomes architecture. To correct the premise is to expose a whole chain of promotions, budgets, reports, and deferential silences along with the original misjudgment.

\emph{Morality.} This is the layer that matters most for this book. Many modern beliefs come welded to a virtue. To doubt the claim is to look indifferent to compassion, safety, equality, the nation, freedom, or science, whichever happens to be attached. The weld is lethal to thought. Values deserve advocacy, and causal claims deserve testing. When the two are fused, evidence that a policy does not work is heard as an attack on the good it was meant to serve. The nobler the goal, the less its failure can be mentioned.

\emph{Tribe.} A belief is also a password. What it says may matter less than what saying it announces. People learn which facts earn warmth and which produce that half-second of silence in which a room changes temperature. Then they absorb the boundary. The corridor never slams shut. It closes with a thousand career-minded coughs.

This explains why the uninformed are often more invested, and not less. Knowledge shows you the branching possibilities, the counterexamples, the problems of measurement, and the humiliating size of what is unknown. Ignorance sees a clean landscape and takes the absence of visible complications for proof that there are none. Kahneman's fast and slow thinking fit here, if we skip the cartoon version \cite{kahneman2011thinking}. Fast judgment supplies a coherent story. Slow judgment could check it, but it is lazy and expensive, and it is mostly hired to do public relations for a conclusion already reached. Mental shortcuts are no defect of the machine. They are how a finite mind operates at all \cite{Tversky1974}. The political question is which environments correct them and which turn them into weapons.

In a healthy culture, error can be digested. Forecasts are scored, policies are audited against a baseline, and papers are replicated. A leader can say ``I was wrong'' without handing the whole institution to its enemies. In a declining culture every admission is ammunition, so each faction buries its own mistakes and preserves its opponents' in amber. What you get is a museum of unforgiven errors and a landfill of unadmitted ones.

\section{Dunning--Kruger, without the cartoon}

The Dunning--Kruger effect is usually cited as an educated way of calling someone stupid. That use is itself tribal and lazy. The original studies were about calibration. People who did badly at logic, grammar, or humor greatly overrated how they had done, in part because the skill needed to perform well is the same skill needed to recognize a good performance \cite{Kruger1999-mv}. Incompetence damages the answer and the error detector together.

The meme has outgrown the evidence. Later work argues that some of the famous graphs are inflated by regression to the mean and by measurement error \cite{GignacZajenkowski2020}, and Dunning's own review is far more careful than the internet slogan \cite{Dunning2011-hf}. The defensible conclusion is not that the ignorant are uniquely arrogant. It is that everyone judges himself badly, and worst where feedback is weak and the standards of the field are invisible.

That more modest finding is the more alarming one, because public life online gives terrible feedback. Confidence shows at once, and calibration does not show at all. A forceful falsehood collects attention today, and the correction arrives next week, qualified and dull. The expert speaks in ranges because he has met the exceptions. The ignoramus speaks in absolutes because he has not. The platform measures engagement and promotes the cleaner signal. A man with a hundred unscored predictions feels bolder than a forecaster with a published record. One has self-esteem and the other has error bars.

Worshiping credentials is no cure. Experts can be bought, herded, vain, and wrong. The cure is \emph{scored exposure to reality}. Predictions should be made before the outcome is known, and procedures fixed before the results arrive. The failures must be kept on file, and critics allowed to read the file. Humiliation does little against arrogance. Accounting does a great deal.

\section{The vanity of the dissenter}

An attack on conformity has to include an attack on counterfeit dissent. Otherwise criticism of the mainstream becomes a new mainstream with worse footnotes. Readers who have been nodding along should take this section personally. So should the author.

A dissenter is not automatically an independent thinker. Contrarianism can be as tribal as conformity, and it has pleasures of its own. There is the thrill of seeing what the herd cannot see. There is the company of fellow initiates. And there is an explanation for one's lack of success that leaves one's genius intact. If the establishment rejects the claim, that proves the claim is dangerous. If it ignores the claim, that proves suppression. If it engages, that proves fear. The theory feeds on every possible response.

This is narcissism in rebel costume, and it copies the establishment's worst habit, which is to make identity the guarantee of truth. The official says ``trust me, I am certified.'' The contrarian says ``trust me, I am censored.'' Neither has offered any evidence.

Persecution is not validation. History has its punished truth-tellers. It has vastly more people who were punished, ignored, or laughed at and were simply wrong. Galileo does not endorse everyone who compares himself to Galileo. A scar proves that somebody hit you. It proves nothing about your cosmology.

Counterfeit skepticism gives itself away. It asks endless questions and would accept no answer. It treats a hole in the official account as proof of its own. It demands perfect evidence from institutions and takes anecdotes from friends. It reads the absence of evidence as evidence of a cover-up. It remembers the predictions that landed and deletes the rest. And it confuses a fearless manner with a courageous method.

Real skepticism is expensive, because it threatens every conclusion, the skeptic's included. It says what would reduce its suspicion. It keeps ``not proven'' apart from ``false,'' and ``possible'' apart from ``probable.'' It does not leap from one institutional lie to the claim that institutions do nothing but lie. Distrust of everything gives no independence. It is an auction in which belief goes to whoever tells the best story.

Nor should every unpleasant consequence be called censorship. If a speaker insults his audience and the audience walks out, nobody's freedom has been reduced. If a paper is rejected for flaws the referee can name, knowledge has been served. If a professional keeps asserting things without evidence in a field where clients depend on accuracy, discipline may be in order. The line is crossed when the standards depend on the viewpoint, when the penalties are out of all proportion, when the process is hidden, or when the institution will not say what evidence would reopen the question. Mere discomfort demonstrates nothing. A corridor is demonstrated by asymmetry that is built into the structure.

A wider corridor will let in more truth and more rubbish. That is the price of discovery. Nobody is owed equal amplification. The aim is to sort better without handing final authority to a sorter who cannot be challenged.

The dissenter's protection against his own vanity is to commit himself in advance. He should state the prediction and date it, name the evidence, and keep the failures. He should be able to summarize his opponent fairly, and to tell corruption from error and coordination from shared incentives. The purpose of dissent is better contact with what is there. Enjoying the view from outside is no purpose at all.

There is a last irony. Narrow corridors manufacture irresponsible opposition. When fair questions cannot be asked in respectable places, they move to places with lower standards. Careful critics learn to insinuate, because plain speech is punished. Audiences come to trust provocateurs, because the institutions sound evasive. And the center then points at the resulting madness as a reason for tighter control.

Repression and crankery are dance partners, and each makes the other look necessary.

\section{Tribe before truth}

We are not solitary reasoning machines. We are coalition animals. We keep track of who can protect us, punish us, hire us, marry us, cite us, invite us, and forgive us. Beliefs therefore have social prices.

Jonathan Haidt's work shows how much of our reasoning is done in the service of intuitions and group loyalties already in place \cite{Haidt-righteous-mind}. This does not make morality a sham. It makes moral eyesight selective. Every tribe has superb vision for the hypocrisy, cruelty, and motivated reasoning of the other. Its blind spot sits exactly where looking at itself would cost the most.

The pattern is symmetrical, even where power and evidence are not. One camp makes the nation sacred and calls inconvenient facts betrayal. The other makes emancipation sacred and calls inconvenient facts harm. Each treats official authority as a conspiracy when it stands in the way and as legitimacy when it helps. Each supposes that only the enemy has an ideology.

The exchange rates are well known. From an ally a claim is testimony, and from an opponent it is propaganda. Doubt that favors our side is called nuance, and doubt that favors theirs is evasion. When our extremist speaks he is an unrepresentative lapse, and when theirs does the mask has slipped. Censorship is safety when we do it and tyranny when they do. Expert consensus is decisive when it agrees with us and ``captured'' when it does not.

Much of this is not even hypocrisy. It is perception, and it happens before thought. By the time an argument reaches consciousness it has already been through customs.

Tribalism also changes what it means to change one's mind. In an inquiry, to revise a view shows that you are responsive. In a faction fight it looks like desertion. The convert may gain the truth and lose his friends, his audience, his income, and his sense of being a decent person. So people need far more evidence to cross a social border than to adjust a private estimate.

The most dangerous institutions are those in which factual disagreement can no longer be told apart from disloyalty. Once that happens the channels of information become ceremonial. Reports still circulate, hearings are held, and experts appear, but all for the purpose of acclamation. The institution has stopped asking reality what happened. It asks an approved reality to confirm the institution.

A tribe can survive a great many errors as long as it keeps some way of admitting them. It becomes terminally stupid when the correction of error is itself treated as treason. From then on better information makes things worse, because each side reads the other's evidence as proof of hostile coordination. The cure has become fuel.

\section{The window and the corridor}

The \emph{Overton window} is named after Joseph P. Overton. It is not the ``overtone window,'' tempting as the slip is. His model describes the range of policies a politician can support without unbearable electoral risk. Inside the window an idea is viable. Outside it an idea may be perfectly legal to utter and still fatal to champion. The window moves as culture, events, and advocacy move \cite{MackinacOvertonWindow}.

The concept is useful because it is no theory of truth. A policy can be inside the window and disastrous, or outside it and wise. The window is a weather map of permission. It is no compass.

The German \emph{Meinungskorridor}, ``opinion corridor,'' is a cousin. It came from the Swedish \emph{\r{a}siktskorridor}, a term associated with the political scientist Henrik Oscarsson's intervention of 2013. A linguistic study by Charlotta Seiler Brylla traces its use in quarrels over what may be said, and shows that the complaint ``the corridor is narrow'' can itself be a rhetorical move \cite{SeilerBrylla2019}. The caution is worth keeping. Sometimes a speaker really is punished for stating a relevant fact. Sometimes he has merely been contradicted and has promoted his bruises to martyrdom.

Three kinds of permission need to be kept apart. \emph{Legal}: may the sentence be said without the state stepping in? \emph{Social}: may it be said without expulsion from respectable company? \emph{Political}: may a man in office act on it and stay in office? Constitutional free speech concerns the first, the corridor the second, and the Overton window the third. Mixing them up produces endless bad arguments. ``You are not in prison'' does not show that the corridor is wide. ``People criticized me'' does not show that it is narrow. ``Most voters reject it'' does not show that it may not be discussed.

The corridor is a price list. In the middle are the statements that buy status. Near the walls are those that need a disclaimer, a credential, or a ritual declaration of good intentions. Beyond the walls are the statements that cast doubt on the speaker's character. Further out are those that bring dismissal, investigation, or the attention of a prosecutor. Each institution has its own list. A sentence cheered at a party congress will end a career at a university, and the reverse. There are many overlapping corridors, with different landlords.

Yet they line up. Editors come from the same schools as civil servants. Staff move from foundations into ministries. Professors become commentators, and platform policy borrows its language from the activists and the regulators. Nobody has to give an order. A shared class interest produces the same boundaries everywhere. Conspiracy is an optional extra, and sociology is enough.

\subsection{The spiral}

Elisabeth Noelle-Neumann described the ``spiral of silence.'' People watch the climate of opinion and keep quiet about views they expect will isolate them. Their silence makes the dominant view look still more dominant, which silences more people \cite{Noelle-Neumann}. Confidence grows in public on one side while doubt piles up in private on the other.

This is more than cowardice. Isolation is costly, and the cost runs from awkwardness to unemployment. Each silence makes sense for the person keeping it, and the sum of them falsifies the picture for everyone.

Timur Kuran added the economics. People state in public preferences they do not hold in private, because of what public statement costs and pays \cite{Kuran1995}. The published spread of opinion then no longer matches the private one, and what looks stable is brittle, because a great many people are waiting for a sign that they are not alone. Add \emph{pluralistic ignorance}, in which individuals privately reject a norm and wrongly suppose that everybody else accepts it. In the classic study, students overestimated how comfortable their fellow students were with heavy drinking, and so kept up a norm that few of them liked \cite{PrenticeMiller1993}. Public opinion can be a ventriloquist act performed by people who each think the dummy is everybody else.

Put these together and the result is disturbing. A society can have no majority that believes the official position and a large majority that repeats it. Everyone expects everyone else to repeat it, and the institutions take the repetition for belief. The official position is then socially real and sincerely held by hardly anyone. It governs appointments, headlines, and policy. And when some event reveals the hidden distribution, the consensus can collapse with astonishing speed. What had looked like conversion was only coordination.

None of this means that every exiled idea is a treasure. Most intellectual garbage is outside the corridor because it is garbage. The case for a wide corridor does not rest on equal respect for all claims. It rests on this: a society that cannot tell refutation from expulsion will sooner or later expel the man carrying the fire alarm.

\subsection{How the corridor moves}

The walls are moved by repetition, prestige, emergency, ridicule, and euphemism. An unthinkable proposal is discussed first as a curiosity. Then it is an unfortunate necessity, then a moral test, and at last common sense. The reverse journey is just as common. A former commonplace is found embarrassing, then harmful, and finally becomes evidence that the speaker is unfit for decent company.

Language goes ahead as the advance party. Call a restriction ``protection'' and you have imported a beneficiary. Call a benefit ``dependency'' and you have imported a disease. Call dissent ``misinformation'' and you have imported a judge. Call enforcement ``accountability'' and you have imported guilt. The label does not win the argument, but it decides which side starts in the dock. Bernays wrote openly about organizing opinion in a mass society \cite{bernais-propaganda}. Persuasion today is faster and has no visible author, and the lesson is unchanged. Whoever supplies the categories has usually won before the evidence is heard.

The result is a built-in delay. A new fact from outside the frame must fight first for a vocabulary, then for a hearing, then for official recognition, and only after all that for action. A falsehood inside the frame starts at the fourth step. That is how a civilization can be drowning in information and starved of thought.

\section{Mass formation without mysticism}

``Mass formation'' is a dramatic phrase. Used carelessly it suggests a hypnotized public moving as one body under a master's spell. That picture flatters the dissident, who becomes the only adult in a stadium of zombies.

The better account is less cinematic and needs no trance. Mass coordination grows out of ordinary things: uncertainty, imitation, fear for one's reputation, repetition, and signals from institutions. People need not believe the same thing for the same reason. They need only act as if enough others believe it.

Solomon Asch's conformity experiments are still instructive \cite{Asch1956}. Subjects had to make an easy visual judgment after hearing a unanimous group give an obviously wrong answer. The pressure raised their error rate. But many resisted, and a single ally in the room greatly weakened the effect. So the experiment does not show that people are sheep. It shows that feeling alone changes the price of trusting your own eyes.

Mark Granovetter's threshold model shows how very different private dispositions can add up to a sudden collective act \cite{Granovetter1978}. One man acts at once. A second acts when one other has, a third waits for two, and so on down the line. If the thresholds form a chain, a small spark runs through the whole crowd. Remove one link and an almost identical crowd stays put. And Cass Sunstein's work on group polarization shows that like-minded people who deliberate together move toward a more extreme version of what they already thought \cite{Sunstein2002}. A room does not average its occupants. It distills them.

Assemble the parts and the process has a familiar course. People sense disorder and have no explanation. A story arrives that names a cause, a victim, a culprit, and a remedy. Respected institutions signal that the story is sound. Media and colleagues repeat its terms until familiarity feels like evidence. Agreement becomes cheap, and asking questions begins to cost. Doubters imitate confidence so as not to stand alone. Every public declaration lowers the next person's threshold. The position becomes a test of character, so that revising it feels like corruption. Finally forms, budgets, and statutes preserve the whole thing long after the original emotion has cooled.

At no stage does everybody have to be deceived. The journalist finds the story useful if oversimplified. The official has his doubts and fears disorder. The activist believes every word. The corporation signs up as insurance for its brand. The citizen repeats the formula to avoid a row. Their motives differ and their signals agree. ``Everyone believes'' is therefore a category mistake. Some believe, some pretend, and some defer. Others are silent or have never thought about it. Institutions squash all these states into a tick in a box and call the total consensus.

Picture a ministry with a hundred analysts. Twenty strongly favor a new program and thirty mildly. Thirty doubt it but expect the leadership to back it. Twenty are opposed. The leadership announces that the program embodies the ministry's values and asks every division for an implementation plan. The thirty doubters conclude that resistance is pointless. Fifteen of the opponents fall silent and five resign. Six months later the ministry reports ``broad internal support.'' No poll was rigged, because no poll was taken. Consensus was manufactured by asking for implementation where judgment should have been asked for.

Or take a viral accusation. The first accounts mention it as a possibility. Influential voices treat hesitation as indifference to the victims. Employers issue statements. Journalists report the statements as ``mounting concern,'' and politicians cite the coverage. The evidence has not changed by a gram. The social fact has changed completely. The accusation is now believable because institutions are acting as if it were believed. Each participant sees only his local signal: the employee sees the memo, the employer the trending topic, the platform the engagement, and the journalist the statements. The loop has no sovereign author. It has a business model.

Mass formation, then, is \emph{coordinated believability}. It decides which statements get repeated and also which doubts can be understood at all. A criticism phrased in the story's own moral vocabulary may be tolerated. One that questions the vocabulary sounds mad or wicked. The formation controls the exits.

It can also dissolve. A visible dissenter breaks the isolation. A failed prediction raises the cost of repeating the formula. A defector from the elite signals permission, and the thresholds cascade the other way. But dissolution gives nobody back the lost years. The careers, budgets, laws, textbooks, and broken trust all remain. The crowd wakes up in a room it has already rearranged.

\section{Law as the bouncer}

A corridor can be enforced without law. Mockery, exclusion, hiring decisions, professional discipline, and a quiet word with the algorithm are usually enough. But law changes the physics. It turns a contested boundary into an official one, backed by investigation, expense, and force.

Every durable legal order restricts some speech. Threats, fraud, perjury, defamation, and direct incitement do not disappear when somebody recites ``freedom of expression.'' The hard question is whether a rule is aimed at a well-defined injury or whether it deputizes the authorities to police disputed reality.

Policing belief is not the property of one side. National security turns dissent into disloyalty. Public safety turns doubt into danger. Public morals are used against the heretic in sex, religion, and art. And the language of ``harm'' can be stretched until disagreement itself counts as an injury. Every tradition has its favorite emergency and its favorite censor.

The escalation climbs a ladder. A claim is said to be wrong. Then it is irresponsible, then harmful, then a sign of a dangerous disposition. Next it is a risk calling for administrative attention, and at the top it is conduct fit for punishment. Each rung sounds modest. Together they carry a sentence from the seminar table to the police file.

Hannah Arendt saw why factual truth is so exposed in politics \cite{Arendt-1967-truth}. Opinions can haggle, and facts are stubborn. A government can split the difference between two tax rates. It cannot split the difference between an event having happened and not having happened. Power puts up with rival values much more easily than with a verified record of what power did.

And the law need not proclaim an official truth in order to chill inquiry. The process is punishment enough. A complaint, an investigation, a seized laptop, a licensing review, and three years of litigation cost far more than any eventual fine. The acquittal then proves much less than civil libertarians hope, because the system has already said what it had to say through the ordeal. The same holds in private government. A platform label is no prison sentence, and an employer's inquiry is no criminal trial. But a society can keep free speech in its constitution while making speech that matters expensive everywhere else. Formal liberty and practical muteness get along very well.

Mill's defense of open discussion rested on our fallibility. The silenced opinion may be true, or partly true, or needed to keep an accepted truth from stiffening into prejudice \cite{Mill1859}. Mill was not naive about error. He was suspicious of the confidence it takes to give today's administrators the last word on tomorrow's knowledge. Orwell dramatized the final stage, in which power no longer merely forbids statements and claims authority over the categories in which statements can be formed \cite{Orwell-1984}. The most efficient censorship does not delete the answer. It removes the question from respectable grammar.

Not every moderation rule is totalitarian, and not every censored claim is true. Inflated comparisons are bullshit too. The line worth drawing runs between rules that deal with conduct and demonstrable victims and rules that give an institution open-ended authority over contested propositions. The first kind can protect inquiry, and the second annexes it.

Any restriction on speech should have to pass six tests:

\begin{enumerate}
  \item \textbf{Specificity:} Is the forbidden conduct defined narrowly enough to be predictable?
  \item \textbf{Evidence:} Must a concrete injury be shown, or is offense presumed?
  \item \textbf{Intent:} Does guilt require knowledge or recklessness?
  \item \textbf{Independence:} Can a separate body review the decision?
  \item \textbf{Symmetry:} Would its defenders accept the rule in their opponents' hands?
  \item \textbf{Reversibility:} Are there appeals, time limits, and corrections of the public record?
\end{enumerate}

The symmetry test is the acid. Every faction imagines its rule administered by its own saints. Institutions have to be designed for the day they are captured by the people you fear.

When law becomes the bouncer at the door of knowledge, truth needs an identity card issued by power. By the time the card arrives, the building may be on fire.

\section{Factories of believability}

No complex society can check everything from first principles. We depend on laboratories, courts, archives, statistical offices, newspapers, universities, and professional bodies, and lately on platforms and ranking systems as well. Trust cannot be avoided. The man who claims to think entirely for himself is generally relying on sources whose mediation he has stopped noticing.

The trouble begins when an institution confuses \emph{being trusted} with \emph{being trustworthy}. Trustworthiness is a practice. It consists of open methods, unwelcome evidence kept on file, disclosed conflicts, and a memory for one's own mistakes. Trust is only how the audience responds. Once that response becomes the target, managing the reputation pushes out correcting the errors.

Four factories work on this principle.

\emph{The credential factory.} A credential is useful shorthand. It says that someone has been trained and examined. But shorthand loses information, and a guild will defend its territory as readily as its standards. The right response is neither sneering at experts nor kneeling to them. It is to ask experts for the one thing that makes expertise valuable, which is disciplined contact with a reality that can prove them wrong. The expert to trust is the one whose uncertainty has a structure. He can say what is measured and what is inferred, and what would reverse his advice. Confidence without an audit trail is costume jewelry.

\emph{The media factory.} Selection is unavoidable in journalism. Trouble starts when every outlet selects alike and shared sourcing passes for independent confirmation. A report cites an official. A commentary cites the report, and a politician cites the commentary. The official then cites the growing consensus. Four institutions have spoken, and one proposition has gone round in a circle. Speed makes it worse. The first frame captures the audience and fixes the burden of proof. A dramatic falsehood can be round the world while its careful correction is still choosing a verb.

\emph{The bureaucratic factory.} Bureaucracies know the world through categories. Cases become fields, risk levels, and indicators. Once money and promotion hang on the indicators, staff learn to produce the look of success. Red cells on the dashboard call for explanations and green cells bring peace. Bad news travels upward with its author removed, and reaches the top as ``implementation headwinds.'' The top then broadcasts reassurance, which the lower ranks read as an order to find supporting numbers. You cannot point to the place where the institution lies. It goes false layer by layer.

\emph{The platform factory.} A platform need not believe a claim to make it believable. It only has to rank it, repeat it, and surround it with signs of attention. Users infer importance from visibility and visibility from engagement, and outrage and certainty beat calibration every time. A ten-word verdict travels further than a ten-page method. A screenshot outlives the correction attached to the original.

The four feed one another. A viral claim prompts coverage, coverage prompts an official response, the response becomes content, and the metrics justify more coverage. Somewhere on that circuit prominence is quietly converted into corroboration.

Distrusting everything is no remedy. Cynicism is gullibility wearing black. If all institutions lie, then the most charming liar wins. The remedy is \emph{traceable trust}. A claim should travel with its pedigree: sources, methods, conflicts, revisions, record of past predictions, and the disagreements along the way. For any claim with large consequences, ask how many truly independent sources stand behind it. Ask which part is observation and which is commentary on commentary. Ask whether the categories were fixed before the result, who pays if the claim is wrong, where the log of corrections is kept, and whether a critic could reproduce the analysis without the institution's leave.

The number of logos on a claim is not the number of independent observations behind it.

\section{Why correction comes too late}

If false beliefs could be changed like light bulbs, this chapter could end with a plea for media literacy. But beliefs are built into systems. They shape investments, careers, regulations, technologies, alliances, and vocabularies. Correction has to overcome an error and an infrastructure with it.

Economists know about technological lock-in. An early lead makes a path more attractive as investments pile up around it, even when a rival path would have been better \cite{Arthur1989}. Political institutions behave similarly. Timing and sequence matter, the benefits are concentrated, constituencies form, and the cost of reversal climbs \cite{Pierson2000}. Belief systems lock in the same way, and six latches hold them shut.

\emph{Reputation.} Those who vouched for a claim cannot revise it without revising their own standing. The louder the original assertion and the rougher the treatment of doubters, the dearer the retreat. An institution will defend a rotting position because an apology would bring its past conduct to light, and the state of the evidence has little to do with it.

\emph{Institution.} Programs create offices, and offices create procedures. Procedures create compliance industries, and all of them generate documents proving that they are needed. The original problem may shrink while the apparatus stays. No organization announces that its mission was accomplished or mistaken. It discovers new dimensions of the problem.

\emph{Money.} Jobs, contracts, property values, subsidies, and whole business models attach themselves to a narrative. A fact that threatens the narrative threatens cash flow, and it will be assessed by a network with a payroll and by no neutral mind. Interests do not have to buy false belief outright. It is enough to fund doubt about expensive truths and certainty about profitable ones.

\emph{Law.} Once a belief has been written into regulation, precedent, or liability, revising it becomes expensive. Officials follow the rule because it is the rule. Yesterday's judgment about the facts is today's working environment.

\emph{Network.} People choose their friends, feeds, and experts partly by shared belief. In time, contrary information arrives only by way of hostile messengers. Even an accurate correction is refused, because accepting it would mean trusting the other side. The network has immunized itself by controlling the source.

\emph{Vocabulary.} The deepest latch is in the language. Once the key terms have acquired moral and administrative force, a rival description can hardly be stated without accepting the frame. Nietzsche's early essay on truth and lying describes how worn-out metaphors harden into what passes for reality \cite{Nietzsche-WahrheitLuege}. A society can lose access to facts by losing neutral words for them.

Physicists have a word for the combined effect, \emph{hysteresis}. The way back is not the way in run backwards, and removing the original pressure does not restore the original state. A magnet keeps its alignment. The panic ends and the emergency powers stay. The theory fails, and the careers built on it go on reviewing grant applications. The slogan is retired, and its categories still organize the databases.

This is the exact sense in which it may already be too late. It is not too late for reality to win in the abstract. Reality always wins in the end, because bridges fall and ledgers stop balancing. It is too late for recognition to prevent the damage. It is too late to restore the trust that confident error used up, and too late to give back the silenced years to the people who paid for an institution's vanity.

The delay accumulates as debt does. A small falsehood spares somebody a small embarrassment. Then further statements are needed to protect the first one. Metrics are adjusted and critics classified, and archives are forgotten until the law codifies the workaround. The interest compounds. In the end the truth can be afforded only through default.

``Already too late'' should therefore be used with discipline. It does not mean that every outcome is fixed and effort is useless. It means that rescue will not be restoration. A late correction limits further harm, but it cannot bring back the branch of history that was not taken. The patient may survive without getting the organ back. The institution may reform without recovering its innocence.

\section{The terminal signs}

How can you tell that an order of belief is nearing the point where correction will come too late? The presence of error tells you nothing, since every order is full of it. Look at how the order treats its errors. Six signs are decisive.

\begin{enumerate}
  \item \textbf{No conditions of defeat.} Major claims are never paired with an observation that would count against them.
  \item \textbf{Lopsided skepticism.} Friendly claims get charity, and hostile ones face impossible standards.
  \item \textbf{Punished messengers.} Critics are judged first by identity, tone, and associations, and by their evidence later or never.
  \item \textbf{Institutional amnesia.} Failed predictions and reversed guidance vanish without any accounting.
  \item \textbf{Merged authority.} The body that declares what is good also claims the last word on what is true.
  \item \textbf{The demand for trust.} An institution that keeps asking to be trusted while resisting audit is announcing that reputation has taken the place of method.
\end{enumerate}

There are reforms worth trying. Protect internal dissent. Separate forecasting from advocacy. Publish prediction records, and reward the critic who turns out to be right. Fix the decisive tests in advance. Put expiry dates on emergency rules. Keep the original data and the history of corrections. Diversify sources, and not merely the faces in the photograph. Give adversarial reviewers status, time, and access, and not a ceremonial three minutes at the end.

But there is a cruel catch. The institutions that most need these reforms are the least able to see it. A system that treats dissent as sabotage will regard protected dissent as a subsidy for saboteurs. A bureaucracy that lives on vague metrics will call an audit simplistic. And a faction that identifies itself with virtue will call the symmetry test moral surrender.

That is why decline is mostly the work of sincere people. Villains are convenient because they can be removed. The durable failure is a mesh of decent motives, local incentives, cowardly silences, and procedures that nobody designed as a whole. Every participant can explain what he did, and none of them can justify the result.

The last chance of correction is seldom recognized as one. It looks like an awkward meeting, a disputed chart, an unfashionable paper, or a whistleblower with poor social skills. It may be a vote to delay, a judge who declines to stretch a definition, or an editor who prints a sentence that will anger the subscribers. Later generations call these moments turning points, after watching everybody turn the wrong way.

\section{After the corridor closes}

Truth does not die when a society refuses it. It becomes archaeology.

The files are opened decades later. The vocabulary relaxes. Officials concede that ``mistakes were made.'' A documentary finds the memo. Scholars explain that nobody at the time could have known, although some people did know and were punished for the inconvenience. Statues get explanatory plaques. Institutions perform contrition with money supplied by the people who inherited the consequences.

The declining society is not necessarily the one that holds the most false beliefs. It is the one whose beliefs cannot be revised without civil war inside the self, the tribe, or the institution. Its citizens mistake confidence for competence and repetition for corroboration. They take what is legal for what is true, and good intentions for good results. Its public language is bullshit in Frankfurt's exact sense, aimed at the impression, while the truth waits outside without an appointment.

The decisive question remains brutally simple: \emph{when would you believe it?}

Would you believe a fact if an enemy had measured it? If it damaged your cause? If it made your past self look a fool? If saying it cost you standing? If acting on it meant admitting that the cheap solution expired years ago?

If your answer is ``only when my tribe authorizes it,'' then membership has been placed above the world. If it is ``only when the law certifies it,'' then power has become the notary of reality. If it is ``only when everyone else says so,'' the spiral of silence has gained another turn.

The last permissible thought is no particular political opinion. It is the thought that one's own side may be wrong in a way that matters, that the approved words may hide the mechanism, and that recognition may arrive after the remedy has expired. A society can still be corrected only while that thought can be spoken aloud before the ruins make it obvious.

Reality has infinite patience.

We do not.


% ===================== Chapter file ch05 =====================
\chapter[Too Late to Change Course?]{Too Late to Change Course? Climate, Fervor, and the Cost of Freezing a Moving Planet}
\label{chap:moving-climate}

\section{The moving target}

The climate is not a machine with an off switch. It is a planet in motion: air, ocean, ice, rock, soil, forests, clouds, living things, and lately human industry, all trading energy and matter with one another. It has never held still, for a day or for a geological epoch. The Greeks did not live in the climate of the Romans. Rome expanded in one climate and shrank in another. The medieval North Atlantic was not the Europe of the Little Ice Age, and the glaciers that crept down Alpine valleys in the seventeenth century are not the glaciers now backing away from their moraines. Change is the ordinary condition of the climate system and always has been.

A society that has got used to a certain range of temperature, rainfall, coastline, and energy price would like the target to stop moving. People want the climate fastened to the moment at which they first noticed it, or to a sentimental picture of an original Earth. German has an exact word for the wish: the climate is to be \emph{festgezurrt}, lashed down like cargo. The wish is understandable. It cannot be granted.

It turns dangerous when it becomes a political program in which any departure from the chosen snapshot is a moral crime. Listen to how climate activism talks. Risk management has a dull vocabulary of probabilities, costs, and trade-offs, and that is not the vocabulary in use. What one hears is emergency, purification, confession, betrayal, and judgment to come. Whether a measure reduces a particular risk at a bearable cost has become a secondary question. The first question is whether supporting it shows that one stands on the right side of history. Virtue is no longer measured in tons of carbon avoided per euro, or in lives made safer per public dollar. It is measured in willingness to sacrifice.

I use the word \emph{sect} for this tendency, and I use it deliberately. I do not mean that every environmentalist is irrational or every warning false. I mean a recognizable social form, one that protects a doctrine by moral pressure and cares less about what the world will do than about what its members must be seen to believe.

An old religious pattern explains why it feels that way. German calls it \emph{Erbs\"unde}, inherited sin; the English term is original sin. In Christian teaching it is a condition one is born into, and no particular deed is needed to incur it. Climate rhetoric has produced a secular copy. Man is guilty for what he is, a creature that breathes, eats, heats its house, and travels, whatever he may or may not have done. The carbon footprint is the confession. Abstention is the penance. Net zero is salvation, the activist and the approved scientist interpret the law, and the doubter is a heretic. Nobody is reviving the catechism on purpose. But the emotional grammar is all there, with its universal guilt, ritual cleansing, and promised redemption. Chapter~\ref{chap:weaponized-original-sin} takes this inheritance apart piece by piece.

Let me be plain about what I am not saying. Climate change is not imaginary. The greenhouse effect is real. Atmospheric carbon dioxide has risen steeply since industrialization. The surface has warmed, the oceans have stored heat, the sea has risen, and most glaciers are losing mass. The Intergovernmental Panel on Climate Change concludes that human influence has unequivocally warmed atmosphere, ocean, and land \cite{IPCC2021HumanInfluence,IPCC2023Synthesis}. A cold day in May does not refute this, nor does a medieval vineyard.

But the physics does not dictate a politics. It does not tell us to dismantle industry, ration energy, make electricity unreliable, or keep poor countries poor, and it does not make prosperity a sin. It does not turn a model scenario into a prophecy or an activist's forecast into a scientific result. And it exempts no policy from the question of what it costs and what it achieves, merely because the word ``climate'' appears in the title.

The climate will change. The open question is whether we meet the change with intelligence or with ritual. We can build stronger cities, better water systems, tougher crops, and cleaner and more plentiful energy. Or we can try to punish ourselves into a motionless climate, while the same coal and oil are dug up and burned somewhere else and the people with the least lose the most.

The stakes are high because energy is no ornament of civilization. It pumps the water, moves the food, makes the fertilizer, heats the hospital, cools the house, smelts the steel, and keeps the vaccine cold. It lets a nurse, a farmhand, and a factory do in an hour what once took a day of muscle. A society deprived of cheap and dependable energy does not return to harmony with nature. It returns to disease, hunger, cold, and dependence on whoever controls the ration.

\section{When science becomes a creed}

Climate politics began with fair questions. Industrial societies have changed the composition of the atmosphere, and it was right to ask what that would do, who would bear the cost, and how to measure it honestly.

The trouble starts when concern hardens into a closed moral system. In such a system the line that matters separates the faithful from the unbelievers, and the difference between a measured effect and an uncertain one is lost. The faithful know that the future is already morally decided. Rapid decarbonization is good, growth built on abundant energy is suspect, anything touched by fossil fuel is unclean, and doubters are in somebody's pay. Once that map is drawn, evidence is read according to allegiance. The style has four marks.

\emph{It turns a range of futures into one apocalypse.} A climate projection is a conditional sentence: if emissions follow this path and technology that one, the model estimates a range of outcomes. Activist speech deletes the ``if'' and the range. A scenario becomes ``what will happen,'' and a long trend becomes a deadline.

\emph{It moralizes ordinary needs.} Heating, travel, concrete, meat, a house of one's own, and a factory job are treated as evidence of guilt. The word ``consumption'' wipes out the service rendered. A kilowatt-hour may be a refrigerated vaccine, a water pump, a night shift in a clinic, or a child's reading lamp. When energy is demonized, the poor are asked to give up what the rich already have.

\emph{It cannot criticize itself.} When an expensive measure fails to deliver, the cry is for more commitment, and nobody asks for a better evaluation. A failed prediction is recast as proof that too little was done. A technology that works but smells wrong is expelled from the moral circle, whether it is nuclear power, carbon capture, or a better gas turbine. The doctrine defends itself by narrowing what may count as evidence.

\emph{It uses children as instruments.} The frightened child accusing the adults is powerful because the image bypasses argument. Nobody can answer a child with a spreadsheet and not look a brute. But the force of a child's testimony establishes nothing about the size of a physical effect or the usefulness of a policy. Children deserve protection from climate risks. They also deserve protection from grown-ups who cultivate their fear and then march behind it.

The classroom offers a second kind of exhibit, in which adults perform science for children. During the 2026 election campaign in Baden-W\"urttemberg, a television crew from ARD followed the CDU's leading candidate, Manuel Hagel, into the Altenburg Community School in Stuttgart. Asked on the spot to explain the greenhouse effect to a seventh-grade class, he began: ``Between Earth and the Sun is the atmosphere. And when it gets thinner, the Sun gets hotter.'' He went on to blame exhaust gases and CO$_2$ for the thinning and the heating, and called the whole business the greenhouse effect \cite{NebelHagelClassroom2026}.

Nearly every link in that little chain is wrong. The atmosphere surrounds the Earth and is no screen hung between two heavenly bodies. Carbon dioxide does not thin it, and nothing we do warms the Sun. The Sun delivers short-wave radiation. The Earth sends energy back to space mostly as long-wave infrared. Greenhouse gases change how that outgoing radiation is absorbed and emitted, and so shift the planet's energy balance \cite{IPCC2021HumanInfluence}. No membrane is wearing thin.

Fairness requires the context. According to the school's head, nothing had been planned; Hagel had expected to help pupils with their work and was caught cold \cite{WittmerHagelSchool2026}. Anybody can botch an explanation without notice. A public that demands instant fluency from politicians and edits every hesitation into a humiliation will get bluffers.

But he had another answer available: ``I know the conclusion, but I could not explain the mechanism properly without looking it up.'' That sentence would have taught the class more than a correct diagram. What they got was a wrong explanation that no pupil was in a position to challenge, because it came from the candidate, in their own classroom, with a camera running.

Set this beside the consultant of Chapter~\ref{chap:who-owns-the-science}, whose carbon dioxide ``gets hotter than oxygen.'' The two stories contradict each other in every detail. They agree in one thing, which is that each takes the shortest imaginable road to the approved conclusion. One turns gases into metal rods, and the other turns the atmosphere into a worn-out parasol. Both are hard to challenge in polite company, because the destination is broadly right. Children should be taught the opposite lesson. A true conclusion does not make every story told in its support true, and ``I need to check'' is one of the first disciplines of science.

Science gains nothing by being worn as a badge. The scientist asks what was measured, what was inferred, what was assumed, and what would count as failure. The propagandist asks whose side you are on.

Two books are useful correctives here. Steven Koonin's \emph{Unsettled} argues that public summaries regularly claim more than the data and models support \cite{Koonin2021}. Vaclav Smil's \emph{Grand Transitions} puts energy transitions back into historical time. No society replaces its energy system by moral decree. It builds mines, ports, grids, refineries, and factories over decades, while demand keeps rising \cite{Smil2021}. Neither author is scripture. Together they establish a point the sect cannot admit: a physical problem can be real while the public story about it is exaggerated, incomplete, and politically dangerous.

\section{What we know, and what we do not}

The debate is poisoned by the demand to choose between total certainty and total ignorance. Climate science contains firm knowledge, strong conditional inference, large uncertainties, and open questions, and an adult society should be able to hold all four in its head together.

The firm part is firm. Greenhouse gases absorb infrared radiation. Adding carbon dioxide changes the balance between incoming sunlight and outgoing heat. The carbon dioxide concentration has risen sharply, and the extra carbon carries the isotopic signature of fossil fuel. The IPCC's best estimate of human-caused warming from 1850--1900 to 2010--2019 is about 1.07 degrees Celsius. Greenhouse gases contributed more than that, and man-made aerosols masked part of it \cite{IPCC2021HumanInfluence,IPCC2023Synthesis}. This is a strong statement about the past. It makes no claim that every storm is caused by carbon dioxide or that every future consequence is known equally well.

Weather is an event and climate is a statistic. A warming climate shifts the odds and the intensity of heat waves and downpours, and it dictates no single one of them. When a town floods, the causes include the rain, the terrain, the drains, the building permits, and the river management, and possibly a changed climate as well. To say that warming raised the odds of heavy rain is different from saying that warming flooded the town.

Then come the feedbacks. Melting snow uncovers dark ground. Water vapor amplifies. Oceans exchange heat and carbon with the air. Clouds cool or warm depending on their height and kind. Ice sheets flow, fracture, and float, and do not behave like thermometers. None of this casts doubt on the greenhouse effect. It explains why confidence varies from one projection to another. That more greenhouse gas warms the planet is one sort of statement. How the rain will fall on a given province in 2090 is quite another.

``Uncertainty'' is treated nowadays as a polite word for denial. It is information. An error bar does not confess ignorance. It states the range of outcomes consistent with the evidence, and to ignore it makes the danger no smaller and the policy more stupid. A society that spends without limit against every remote possibility has less to spend against certainties, among them poverty, unreliable power, weak hospitals, and bad water.

So ``Is climate change real?'' is the wrong question. The right one is this: which risks are large enough, likely enough, and tractable enough to justify which measures, and what do those measures do to people who are already vulnerable? It is less thrilling than a declaration of planetary emergency. It is the question on which lives depend.

\section{A short history of a restless climate}

Any honest account begins long before writing. The Earth has been a hothouse without polar ice and an icebox with glaciers over much of North America and northern Europe. Geologically we still live in an ice age, the Quaternary, which began some 2.6 million years ago. Our own epoch, the Holocene, is a warm interval within it.

The pacemaker of the glacial cycles is known. Slow variations in the Earth's orbit and tilt change where and when sunlight arrives, above all the strength of northern summers, while hardly altering the yearly total. The decisive matter is whether summer is warm enough to melt last winter's snow. When it is not, snow lasts, turns to ice, and reflects more sunlight. The oceans cool and take up carbon dioxide, and the cooling feeds itself. When northern summers strengthen, the sequence runs in reverse. A small push is amplified by feedbacks, and the first cause is not the biggest contributor.

There is also a hole in the theory. About 800,000 years ago the dominant glacial rhythm changed from roughly 41,000 years to roughly 100,000, and NASA says plainly that science has no clear answer why \cite{NASA2024Milankovitch}. This is more than a footnote. It shows that a theory can be powerful and incomplete together. The ice cores also record abrupt regional lurches, tied to ocean circulation and collapsing ice. The system has thresholds. That does not entitle anybody to announce a tipping point with a date on it.

One popular confusion needs clearing up. In the glacial cycles, carbon dioxide often rose after the warming began, as an amplifying feedback. That does not make it harmless today. A feedback in one setting can be a forcing in another, and fossil carbon warms the air whatever started the last deglaciation.

Written history adds a narrow and patchy window, and the patches follow whoever happened to write and whatever happened to survive. A drought in a Roman historian may be a local affair. A bad harvest may be due to war, taxation, or pests as much as to weather. Historians therefore read the chronicles alongside tree rings, pollen, cave deposits, ice cores, and moraines, and trust none of them as a perfect thermometer.

What the combined record shows is variability, and it shows societies that coped with it or failed to. The Greeks farmed a Mediterranean of sharp regional contrasts, and a failed harvest in Attica was not a new climate. Rome grew during what scholars call the Roman Climatic Optimum, a run of comparatively favorable conditions in parts of the Mediterranean and no uniform warm plateau \cite{McCormick2012,Neukom2019}. In the third century came civil war, fiscal strain, plague, pressure on the frontiers, and a less kindly climate. ``Climate caused the fall of Rome'' is as crude as ``climate never mattered'' \cite{McCormick2012,Harper2017}. In the sixth and seventh centuries a cluster of large volcanic eruptions veiled the sun and produced the Late Antique Little Ice Age, which should cure anyone of the idea that ``natural'' means ``harmless.'' The medieval warmth, roughly 950 to 1250, let the Norse farm in Greenland, but it was regional and far from synchronized around the globe. The Little Ice Age, which was sharpest in Europe in the seventeenth and early eighteenth centuries, froze rivers, pushed glaciers into villages, and ruined harvests. Volcanoes, the sun, and ocean circulation all had a hand in it.

Two opposite abuses of this history are common. One says that it was warm in the Middle Ages, so modern warming is nothing new. The other denies that the pre-industrial past had any warmth worth mentioning. Both are too simple. Regional climate varied a good deal before industry. The recent warming is unusually widespread and fast when set against the last two thousand years. Neither fact cancels the other.

Glaciers illustrate the same point. Many reached their greatest recent extent in the Little Ice Age and began to retreat in the nineteenth century, before modern emissions amounted to much. A photograph from 1850 therefore proves nothing about cause. But the pace of loss today is another matter. The monitoring services record net loss worldwide, with rates in the early twenty-first century that have no precedent in the observed global record \cite{Copernicus2026Glaciers,WGMSLatestGlacierData}. As to what a retreating glacier means for the valley beneath, that depends on snowpack, monsoon, reservoirs, and groundwater. Slogans are no help there, and the valley needs a watershed plan.

Through all of it the lesson of history stays the same. A society is not a thermometer. It stores grain, builds aqueducts, chooses its crops, and moves goods and people, or it fails to. The peasants of the Little Ice Age suffered because their buffers were thin. No malice on the part of the climate was needed.

\section{What is different now}

If the climate has always changed, what is special about the present? The answer is that we have added a strong and identifiable push to a system that was already variable, and we have done it while building a civilization of eight billion people who depend on one another's infrastructure.

The mechanism is not in dispute. Carbon that lay underground for millions of years has been returned to the air within two centuries. Greenhouse gases slow the loss of infrared energy to space, and the planet warms until the budget balances again. Orbital cycles cannot explain the present trend, because they are far too slow. The sun cannot explain it either, since its variation is too small and has the wrong pattern, and volcanoes cool. Natural variability still rules from year to year and decade to decade. It does not remove the human share of the long trend.

I state this flatly, because criticism of the activists must not become their mirror image. ``The climate changed before there were factories, so emissions are irrelevant'' is not a serious argument. Neither is the claim that because emissions matter, natural variability has ceased to. The past shows that the system changes for many reasons. The present shows that one of those reasons has become unusually strong and unusually quick.

There is a second half to the modern situation, and the sect never mentions it. No earlier society was as well equipped to protect people from weather. We forecast storms, cool buildings, store food, irrigate, move water, build sea walls, and breed crops for heat and drought. The same industrial system that produces the emissions produces the defenses.

This is the paradox at the center of the whole debate. Civilization is blamed for the risk, and civilization is also the instrument of resilience. A policy that weakens the industrial base in the name of the climate may shave one line of the emissions table while exposing people to heat, cold, hunger, and coercion. Nobody claims industry is pure. The moral question is whether its benefits and its harms are being weighed honestly.

\section{Observation, projection, prophecy}

Public discussion muddles three things. An observation says what has happened. A projection calculates what could happen under stated assumptions. A prophecy announces what will happen, with moral certainty and a date near enough to frighten. Science deals in the first two. Activism prefers the third, because prophecy mobilizes.

Climate models are not crystal balls. They are numerical sketches of physical processes, and they must simplify clouds, soils, ice, aerosols, and human behavior. Every run depends on its inputs. A high-emissions scenario does not promise that the world will follow it. A low-emissions scenario does not show that the technology it requires exists at scale.

Models are not worthless either. Older projections caught the broad course of warming rather well, once allowance is made for the difference between the forcings they assumed and those that occurred. Fair evaluation cuts both ways. A model has not failed because the world took a different emissions path. Nor has it succeeded because the curve looks roughly right when nobody checks the assumptions.

\subsection{Al Gore's calendar}

Al Gore's \emph{An Inconvenient Truth} made climate a mass subject in 2006 and gave it the vocabulary of moral emergency \cite{Gore2006}. Its core message, that emissions warm the planet with consequences for ice, oceans, and people, was sound. Scientists at the National Snow and Ice Data Center judged that he had the basics right, while correcting several details \cite{NSIDC2006Gore}.

Its fate also shows what happens when advocacy boils conditional science down to pictures and dates. The best-known case is the ``ice-free Arctic by 2013.'' The record is messier than the legend. In a speech in 2009 Gore cited a researcher and spoke of a 75 percent chance that the Arctic would be free of ice in some summer months within five to seven years. The researcher later said he would never have put it so precisely. The date came and went. The film itself had spoken of 50 to 70 years \cite{Putterman2021}.

None of this makes Arctic sea ice unimportant, and it has in fact declined substantially. The episode shows the life cycle of a prophecy. A hedged remark by a researcher becomes a date. The date passes and the underlying trend continues. The public is left to choose between cynicism and a demand for still louder alarm. Kilimanjaro is the same story. Its ice has shrunk, but for reasons that include humidity, snowfall, and sunlight as well as temperature, and the summit of one mountain is no thermometer for the globe.

\subsection{The Himalayan glaciers}

The most embarrassing case is in the IPCC's own 2007 report, where Working Group II stated that the Himalayan glaciers could be gone by 2035. No scientific literature supported the statement, and it was withdrawn in an erratum \cite{IPCC2007Errata}. The correction did not show that those glaciers are stable. Many of them are retreating. It showed that a great assessment body can pass on an unsupported claim when review fails.

One blunder in thousands of pages does not discredit the IPCC. But it should make us wary of anyone who answers such an error by saying that the mistake did not matter because the message was useful. Accuracy matters most where a report is used to justify expensive and coercive measures.

\subsection{Hansen in 1988}

James Hansen's Senate testimony of 1988 is often filed under failed prophecies, which is not quite fair. He presented three scenarios for different emission paths. Compared directly with the thermometer, his middle scenario ran about 54 percent too warm. Most of the miss came from assuming too much future methane and too many halocarbons. Fed with the forcings that really occurred, the relation between forcing and warming in his model agrees with observation \cite{Hausfather2020}. He was no oracle, and he was no fraud.

The lesson lies in what happened next. Politicians reach for the most alarming scenario because it is the most useful. The modeler says ``this is what a high-emissions path would do,'' and the audience hears ``this will happen unless you obey.'' By the time the date arrives, politics has moved on, and nobody goes back to check the assumptions.

\subsection{Deadlines}

Then come the countdowns: ten years to save the planet, twelve years to the end, a last window before the tipping point. Some of these refer to a carbon budget, some to a policy target, and some to nothing that can be traced to any assessment.

Some changes really are slow to reverse. Carbon dioxide lingers, the sea level lags, and ice sheets go on adjusting. But ``irreversible'' does not mean a cliff with an appointment. Risk rises with each increment of warming, and each increment avoided is worth something. The statement that every fraction of a degree matters is the opposite of the statement that after a missed deadline nothing matters.

The failed predictions that do most harm are the quiet ones, the implied promises that a given policy will solve the problem. A city bans a fuel, or a country shuts down reliable power stations, and emissions fall in the targeted column and rise in another. Production moves abroad. The local statistic improves and is celebrated. But the atmosphere has no way of registering virtue. It responds to molecules, and to the global total of them.

\subsection{The audit}

An honest scorecard is mixed. The Arctic was not ice-free by the early 2010s, and its ice has nonetheless declined. The Himalayan glaciers will not vanish by 2035, though many are shrinking. Kilimanjaro kept its ice past the advertised date and has lost much of it. Hansen's middle scenario ran warm, because his inputs were too high, and his physics held. The world did not end on any date announced in a speech. Meanwhile temperatures have risen, glaciers have thinned, and the oceans have warmed.

An audit is for learning, and nobody should run one to win a partisan round. A movement that shrugs off its failed predictions will lose the public's trust, and so will one that takes every observed change as confirmation of everything it ever said. Science lives by the readiness to say that this part was right, this part was wrong, this part is uncertain, and this part was never science in the first place.

\section{Who pays}

Activism speaks in the singular, of the planet, humanity, and our future. Risk comes in the plural. The same warming hurts one place and helps another, and the same policy cuts emissions in one country and raises poverty in the next.

The first divide is wealth. A heat wave in a city with air conditioning, shade, clinics, and a working grid is a different event from a heat wave in a settlement where water is carried by hand. The danger of warming cannot be separated from the danger of being poor. The second divide is by sector. Milder winters mean fewer deaths from cold, and hotter summers mean more from heat. Extra carbon dioxide feeds plants, and drought starves them. The balance sheet has two sides, and it comes out differently in different places. The third divide is by generation. Our descendants will inherit a warmer world, and they will also inherit whatever wealth and technology we leave them, which may include cheap desalination, advanced reactors, and hardier crops. A policy that makes the present poorer leaves the future less able to cope. What the generations owe one another includes capability, and the debt is not paid in atmospheric carbon alone.

This is where the language of sacrifice goes bad. The prosperous activist calls for lower energy use from a well-heated flat, and a poor family is told that affordable electricity must be given up for the planet. The campaigner denounces industrial farming and lives off refrigeration, synthetic fertilizer, and global logistics. The celebrity pleads for the oceans from the deck of a very large yacht. Such contradictions do not refute the underlying concern. They show the difference between a principle and a performance.

The fitting response to hypocrisy is a demand for policies that survive contact with ordinary life, and it would be foolish to wave away the ocean on account of the yacht. If a measure requires people to be poorer, colder, hungrier, or less mobile, its advocates owe an account. Who bears the cost? What benefit to the climate results, and is it global or merely local? What were the alternatives? Moral confidence is no substitute for a ledger.

\section{The price of doctrinal deindustrialization}

``Deindustrialization'' is spoken as though it meant the painless disappearance of dirty factories. In practice it means losing the capabilities that hold a complex society together: steel, cement, chemicals, machinery, fertilizer, refrigeration, construction, and the skilled people who know how these fit together. A country can cut its industrial output in only two ways. It can lower its material standard of living, or it can import the goods. In the second case the smokestack leaves the landscape and the factory moves across the border.

Energy is built into matter. Steel needs fierce heat and a reducing agent. Cement needs kilns. Fertilizer needs hydrogen, which today is made chiefly from natural gas or coal. Glass, aluminum, paper, plastics, medicines, and data centers need power that does not flicker. So do wind turbines, solar panels, batteries, and transmission lines, every one of which must be mined, smelted, manufactured, and shipped.

Nothing here argues for keeping every fuel for ever. It argues for treating the transition as an industrial project and dropping the idea of a moral cleansing. Shut capacity down before its replacement stands, and you get scarcity, which is not the same thing as decarbonization. People burn cheaper and dirtier fuel, postpone maintenance, move production to laxer jurisdictions, or simply go without.

I call the outcome \emph{industrial starvation}, and I mean the phrase to provoke. When energy prices jump or supply turns unreliable, the damage travels. Food costs more because fertilizer and freight cost more. Housing costs more because cement, steel, and glass cost more. Hospitals pay more to stay open, and small firms close. The burden falls unevenly. The comfortable household absorbs the bill or buys a battery, and the poor one chooses between heating and eating. The large firm relocates and the small one dies. What the economist describes as an efficient carbon price reaches the worker as a loss of dignity.

\subsection{Energy access is development}

Hundreds of millions of people still have no reliable electricity, and billions cook over smoky fires. For them energy consumption is the precondition of health, schooling, clean water, and paid work. Electric light lengthens the day. A refrigerator keeps food and medicine. A pump brings clean water, and a clean stove spares the lungs of women and children.

The debate becomes morally inverted when rich countries, which reached their own comfort on coal, oil, gas, dams, and reactors, advise poor ones to stay on a low-energy diet for the sake of a climatic baseline. A decent policy cannot ask Africa or South Asia to choose between development and atmospheric virtue. It has to make clean energy cheap and reliable. Making conventional energy dear is not the same achievement. Even the International Energy Agency treats access to electricity and clean cooking as central to development and to any transition worth calling just \cite{IEA2024Africa,IEAAccessAffordability}. ``Transition'' is an honest word only if it means more energy for those who have too little.

\subsection{How fast things really change}

Smil's history makes the rates of transition visible \cite{Smil2021}. Coal did not replace wood overnight, and oil did not replace coal overnight. Energy systems overlap for decades, because old plant stays useful, demand keeps growing, and the new system has to be built with the help of the old.

The same holds for materials. A low-carbon system requires mines, processing plants, ports, transformers, transmission corridors, skilled workers, and enormous quantities of concrete and steel. Wind and solar generation can grow quickly. To integrate them takes grids, storage, backup, and markets that function. Hydrogen is a carrier and no source of energy. Every technology has a supply chain behind it.

These are no arguments against the technologies. They are reasons to ask engineering questions. How much material is needed, and how fast can mines and factories expand? What keeps the grid stable in a windless January week? What becomes of the land, and of the waste? No slogan answers them.

\subsection{The oil that goes elsewhere}

Hans-Werner Sinn's ``green paradox'' adds the economist's warning \cite{Sinn2015}. If one region uses less oil, coal, or gas while the world market stays connected, the price falls and others buy more. The fuel one buyer refuses is not destroyed. It is sold to the next. And owners of reserves who expect stricter policy in future have a reason to pump faster now.

Put simply, if Europe burns less and Asia burns more because the oil has grown cheaper, Europe's figure drops and the world's figure drops much less, perhaps hardly at all. The atmosphere responds to the global quantity and takes no interest in national bookkeeping.

The paradox is no iron law. Reductions in demand do bite when they are large, lasting, and matched by real substitutes. Border adjustments, product standards, and international agreements can limit the leakage. But as a warning against unilateral moral theater the paradox is indispensable. A country can close an efficient power station and celebrate its national statistic, and then import the electricity, the steel, and the chemicals. The emissions have not gone away. They have gone abroad, and the consumer feels virtuous because he can no longer see the factory.

\subsection{Four ledgers}

Every climate policy keeps four sets of books. The atmospheric ledger records how much greenhouse gas is added or removed worldwide. The economic ledger records who pays and who profits. The human ledger records lives saved or lost and poverty relieved or deepened. The political ledger records whether people will put up with the policy, whether other countries will copy it, and whether a backlash will undo it.

A policy can look splendid in the first book and ruinous in the other three. A government's task is to keep all four, and it should not choose the one that feels best.

That is why the demonizing of consumption does so much damage. Lower consumption is no achievement in itself. A household that uses less heat because the house is well insulated has made progress. A household that uses less heat because it cannot pay the bill is deprived. A steelworks that needs less energy per ton has advanced. One that makes less steel while the country imports the difference has only improved a local number. Praise should go to policies that leave people safer, healthier, and better off while cutting emissions on a scale the atmosphere can notice. A policy that merely makes people consume less has earned none.

\section{Competence instead of penance}

The alternative to deindustrialization is technological seriousness, and nobody should mistake it for complacency. Cut emissions where cutting works. Adapt to what can no longer be avoided. Invest in whatever makes both cheaper.

\subsection{Adaptation is not surrender}

In activist circles adaptation is treated as capitulation. This is absurd. Every society adapts to its climate and always has. It builds for the local weather, stores water, tames rivers, and chooses crops to suit. The only question is whether adaptation keeps up with exposure.

Against heat there are shade, trees, reflective roofs, ventilation, warning systems, reliable power, and air conditioning. Air conditioning has become a symbol of decadence. It also saves lives. The sensible policy is to make it efficient and available, and to stop preaching against it. Against floods there are drains that work, channels kept clear, forecasting, zoning, and barriers where they repay their cost. On the coasts there are sea defenses, raised buildings, and planned retreat where defense is not worth the money. These are local decisions, and no global target can tell a town whether to build a wall.

For water there are reservoirs, leak repair, recycling, desalination, and prices that discourage waste. Scarcity is usually a failure of pipes and government before it is a failure of rain. In farming there are crops bred for heat, drought, and salt, shifted planting dates, and better irrigation. For health there are disease surveillance, vaccination, and hospitals that keep running in a heat wave. Wealth is itself a health technology. Rich societies lose fewer people to weather than poor ones do because their buildings, medicine, transport, and food supply are better. Their weather is no gentler.

\subsection{An open portfolio}

The world should release less carbon, and the route matters. The lasting route is to make clean energy better and cheaper than dirty energy. People adopt what is reliable, cheap, and convenient. They will never adopt anything on a global scale because a campaign has told them that consumption is sinful.

So the portfolio should be open. \emph{Nuclear power} gives firm, low-carbon electricity from very little land. Its difficulties with cost, construction time, waste, and public acceptance are problems of engineering and institutions, and they give no ground for excommunication. \emph{Hydropower} is excellent where geography and ecology allow. \emph{Wind and solar} are now major sources of new generation, but their variability has to be paid for in grids, storage, and backup. They should be judged by the reliability and emissions of the whole system and not by the nameplate capacity installed. \emph{Batteries} suit vehicles and short-term balancing, and cannot provide seasonal storage or furnace heat. \emph{Carbon capture} is costly and is no license to emit, yet it may be the only route for cement, steel, and chemicals, and a hard option should be improved before it is condemned for failing to fit a moral story. \emph{Hydrogen and synthetic fuels} belong where direct electrification fails, in some industrial heat, shipping, and aviation. Add better grids, geothermal energy, methane control, high-temperature heat pumps, and everything not yet invented.

A policy that pays for research, demonstration, and deployment lets this portfolio evolve. A policy that crowns one technology as morally pure locks society into a brittle system.

\subsection{The taboo on learning}

Deliberate intervention in the climate is controversial, and it should be. Brightening the atmosphere could shift rainfall and start quarrels between states. Carbon removal devours energy, land, and money. Neither is an escape hatch.

But research is not deployment. To refuse to study a technique because it makes us uneasy is to throw away knowledge that might be needed in a real emergency. A responsible program would examine feasibility, side effects, governance, liability, and what happens when the intervention stops, and it would do so in the open. The taboo is as irrational as the fantasy of a technical cure for everything.

The same goes for the adaptation measures that activists wave away as distractions. A child in a hot city gets no relief from a sermon against cooling. A coastal town is not protected by being told that the rising sea symbolizes adult guilt.

\section{Nine rules for a moving planet}

\begin{enumerate}
\item \emph{Measure outcomes and not intentions.} Count global emissions, reliability, health, displaced industry, and who pays. A national accounting gain is not yet a result for the climate.
\item \emph{Keep energy abundant while cleaning its sources.} The aim is more useful energy with fewer emissions. Less energy is no aim at all.
\item \emph{Adapt now.} Some warming is already committed. Protecting people from it is a duty and gives no excuse for emitting more.
\item \emph{Keep the portfolio broad.} Let technologies compete on performance and cost, and never on ideological pedigree.
\item \emph{Protect the poor from policy shocks.} If a policy raises the price of energy, compensate openly and do not pretend the cost is not there.
\item \emph{Coordinate across borders.} Rich countries must not be allowed to export their emissions while poor countries inherit the restrictions.
\item \emph{Audit predictions.} Record what was claimed, on what assumptions, for what date, and what would have counted as failure.
\item \emph{Protect scientific disagreement.} Consensus marks what is well established. It gives no license to silence criticism of models, costs, or institutional blunders. Facts constrain policy, and they do not choose an ideology.
\item \emph{Remember that civilization is a climate technology.} Buildings, energy, food storage, medicine, transport, and forecasting are what stand between people and the weather. To wreck the capacity to maintain them in the name of climate protection is a contradiction.
\end{enumerate}

\section{Against the frozen picture}

The Earth that activism asks us to save never existed. It is a picture put together from chosen memories: the weather of someone's childhood, a glacier in an old photograph, a favorite coastline, a pre-industrial landscape as painted by people who never had to live in one. The picture has real emotional force. It marks no stable physical state.

The climate has never been fixed. Ice ages came and went with the Earth's orbit. Volcanoes cooled the world within a year. Rome rose in one climate and fell in another. Medieval warmth was patchy, the Little Ice Age ruined harvests, and the glaciers have been retreating since the nineteenth century, lately at a much faster pace.

Human influence is part of that story now, and activists did not invent it. Greenhouse gases are warming the air, the sea, and the land. What follows from this is the whole question. The sectarian answer is to nail the climate to a chosen point, declare every deviation intolerable, and accept whatever human cost that requires. The civilized answer is to lower the avoidable risks while keeping the capacities that make people resilient.

Both sides care. They differ as a morality of sacrifice differs from a politics of competence. Sacrifice may be needed in a real emergency, but a permanent politics of sacrifice turns into a method of control. Competence asks for measurements, alternatives, reserves, and the chance to learn. It remembers that people need heat, cooling, food, transport, medicine, and work, and that clean technology has to be built by somebody and cannot be wished into being. It asks whether a policy works for the globe. Looking virtuous at home is not enough.

Two temptations are to be refused. The first is complacency, the notion that because the climate changed before, our emissions do not matter. The second is fanaticism, the notion that because our emissions matter, anything labeled climate action is justified. Both put identity where judgment should be.

The future will not be settled by whether the Earth returns to some particular number on a graph. It will be settled by whether societies can still produce energy, knowledge, food, shelter, health, and cooperation while the climate goes on changing. The humane course is to make energy cleaner, more reliable, and more plentiful, and to make civilization stronger so that people can live well on a planet that will never stand still. Impoverishing ourselves in the hope that the climate will settle down is the other course, and it will not work.

It is too late to freeze the planet. It is not too late to improve the choices open to the people who live on it.


% ===================== Chapter file ch06 =====================
\chapter[Demography, Migration, and Self-Correction]{Too Late to Change Course: Demography, Migration, and the Capacity for Self-Correction}
\label{chap:demography}

\epigraph{\enquote{That one must answer for the (foreseeable) consequences of one's actions.}}{Max Weber, \emph{Politik als Beruf}}

\section{A book of tables}

In 2010 a former Berlin finance senator and sitting member of the Bundesbank board published a book made largely of tables. Thilo Sarrazin's \emph{Deutschland schafft sich ab} set out birth rates, school results, employment figures, and welfare statistics, and it drew the conclusion that Germany was shrinking, aging, and importing its replacement workforce from the places least able to supply one. Before the book was in the shops, the chancellor let it be known through her spokesman that it was ``not helpful.'' The author lost his post, and in due course his party. A decade later he said it again, with newer tables, in \emph{Der Staat an seinen Grenzen}.

Some of the first book deserved its beating. Its excursions into heredity were amateur genetics, and they gave every critic the excuse he was looking for. But the tables on births, schooling, and welfare were never refuted by the outrage. Most of them were never discussed at all. For thirty years the standard reply to the arithmetic has been an adjective, and adjectives have no effect on arithmetic.

This chapter states the case as its defenders make it, without the customary anesthetic. Weber's distinction supplies the frame. The ethic of conviction asks whether my heart is pure. The ethic of responsibility asks what will happen. German migration policy has been conducted for two generations under the first, and the consequences belong to the second.

\section{Compound interest in reverse}

Population obeys compound interest. With a fertility rate around 1.4, each generation is roughly seventy percent the size of the one before. A hundred grandparents have seventy children, forty-nine grandchildren, and thirty-four great-grandchildren. That is no gentle decline. Within a century a birth cohort shrinks to a third.

Germany has been below replacement for half a century. Educated men and women postponed children, or skipped them, for reasons that each made sense at the time: the career, the rent, the feeling that there would be time later. Meanwhile life expectancy rose. The outcome is a country with more and more pensioners and fewer and fewer people to pay for them. A pay-as-you-go pension system is a promise made by people of working age to the old, and it assumes that there will be people of working age. Contributions go up and benefits go down, and the young are told that this is solidarity.

\section{The substitute}

Immigration was meant to fill the gap, and it does fill the gap in headcount. Whether it fills the gap in contributions depends entirely on who comes, with what schooling, and how soon they find work. A trained nurse recruited abroad and an asylum seeker with four years of school each count as one in the migration statistics. For the pension fund, the primary school, and the labor office they are very different events.

Germany chose not to choose. The guest workers of the 1950s and 1960s were recruited for the assembly line on the assumption that they would go home. ``We asked for workers, and people came,'' as Max Frisch put it. They stayed, and their families followed, and nobody had thought to prepare a policy for that. In 2015 the error was repeated on a larger scale and with greater self-satisfaction. The asylum system buckled, deportations of rejected applicants remained rare, and family reunification multiplied the numbers. The chancellor's ``Wir schaffen das,'' we can do this, was pure ethic of conviction. It named no cost and no limit, and it did not say who ``we'' were.

The demographers have an honest name for the arrangement, \emph{replacement migration}. The term was coined at the United Nations and has nothing to do with the political fringe. A population kept up in this way is not stabilized. One population is substituted for another, cohort by cohort, and the country is committed to ever larger inflows merely to stand still. In the large cities, children from migrant families already form the majority in the youngest age groups. Whatever one thinks of that, it is a transformation of the country within a single lifetime, and no electorate was ever asked to approve it.

\section{What institutions can carry}

Public systems are not infinitely elastic. Germany's were built for a population that was stable, largely homogeneous, well schooled, and trusting. When large numbers arrive without the language, without much schooling, and from places with very different rules of life, the systems give way in predictable ways.

\emph{Schools.} In many urban classrooms most pupils do not speak German at home. The teacher teaches the language and keeps order, and the curriculum gets what time is left. Standards are lowered quietly, because high failure rates would be embarrassing. Two groups of children lose by this. The migrant children do not catch up, and the native children in the same classrooms get a thinner education than their parents had. The international comparisons have shown large gaps for years, and they persist into the second generation. The causes are language, the habits of the home, and what a child has learned before his first day at school. They do not cure themselves.

\emph{Health and welfare.} Waiting times lengthen. Doctors work through interpreters. Unemployment and dependence on benefits remain far higher in some immigrant groups than among natives, even after many years in the country. A welfare state designed for people who pay in for most of their lives is being asked to support people who may never do so.

\emph{Housing.} New arrivals and low-paid natives compete for the same cheap flats in the same cities, and the well-to-do who applaud immigration do not live there.

\emph{Security.} The police statistics show certain groups of young men heavily overrepresented among suspects in violent and sexual offenses. A statistic about a group convicts no individual. But the official habit of publishing such figures late, in aggregate, or not at all has done more for public distrust than any agitator could. A state that hides its numbers passes sentence on itself.

What results is a Potemkin village on a national scale. The speeches celebrate diversity and the reports confirm that the model is working. The schools still hand out certificates and the hospitals still treat patients, and the welfare offices still pay. But the performance of the old Federal Republic has gone from behind the facade. And every wave arrives before the last one has been absorbed, so that there is never a pause in which to consolidate.

\section{The culture nobody is allowed to mention}

Integration on this scale has not worked as promised. Many adults never learn German properly, and in some districts a child can grow up, shop, pray, and marry without needing it. Attitudes that are normal in the countries of origin do not dissolve at the border. They concern the position of women, hostility to Jews, the limits of free speech, and the precedence of state law over religious law. Teachers, social workers, and police officers have been reporting this for decades. To repeat it accuses no individual. It describes what happens when a society imports people faster than it can pass on its own rules, having first lost the confidence to say what those rules are.

The sentimental reply is that culture is enrichment, and that anyone who notices friction has a defect of character. This is how a real problem was handed to the extremes. Fair questions that could not be asked in respectable company were asked elsewhere, by worse people, and the established parties then pointed at those people as proof that the questions were indecent.

\section{Who was at the station}

In September 2015 crowds stood in Munich's central station and applauded the arriving trains. It was the ethic of conviction at its most photogenic. The people who applauded then went home. The bill went to people who had not been at the station: the primary-school teacher, the clerk in the housing office, the nurse on the night shift, and the low-paid tenant looking for a flat. Compassion of this kind is cheap for the giver, because somebody else pays. It also asks for no accounts. Nobody who clapped has ever been asked how it turned out.

A state exists to secure the continuity of its people and its institutions. That includes deciding who may enter, selecting those who can contribute, and insisting that those who stay accept the constitution as the rule for everybody. To fail at this is negligence, and calling it kindness does not improve it. The weaker residents pay for the negligence first.

\section{The point of no return}

Germany cannot go back to the society of 1990. The natives are too few and too old. The newcomers are too many for institutions that are already overloaded. Standards, once lowered, are hard to raise, and the political class still prefers denial, because admission would be a confession.

What remains is management. Pensions will require higher contributions or deeper cuts, and probably both. Schools will struggle to hold what standards are left. The cities will sort themselves into quarters that have little to say to one another. The welfare state will cost more and deliver less. Each further year of unselective immigration makes all of this worse and postpones the painful consolidation that would at least stop the slide.

That is what ``too late to change course'' means here. Nothing is ever fixed absolutely. But the cheap options expired years ago, and the remaining ones are all expensive and unpleasant, and they will be chosen by people who had no say in the decisions that made them necessary. To admit this is to admit defeat. It is also the only honest ground left on which to build anything.

That is the case at full strength, as its defenders make it. It is a verdict, and a verdict should be checked against measurements. Section~\ref{sec:germany-demographic-stress-test} of the last chapter does that, figure by figure, and the reader may judge how much of the verdict survives.


% ===================== Chapter file ch07 =====================
\chapter[The Inheritance Paradox]{The Inheritance Paradox: Devoted Parents, Equal Children, and Dirac's Conundrum}
\label{chap:inheritance-paradox}

\begin{quote}
\itshape
``all children should be given an equal chance, that is, the same
opportunities, for development.''
\end{quote}

\section{A small provocation}

The sentence comes from Paul Dirac, in a short piece called ``The Futility of War'' that he contributed to the 1982 Erice seminar on nuclear war \cite[pp.~173--174]{dirac-81}. Dirac was not a man who wasted words, and his argument is almost too short to be called a theory. It opens a very large door all the same. What follows is my development of it. Dirac supplied the spark and should not be blamed for the rest.

He noticed that two principles which nearly everyone accepts pull in opposite directions. We expect parents to care for their children, to make sacrifices for them, and to improve their prospects. We also think it unjust that a child's prospects should depend on the family he happens to be born into. If the parents succeed, the children no longer start level. If public authority prevents the parents from succeeding, the sacrifice has lost its point. ``These two principles therefore tend to exclude each other,'' Dirac wrote, and he called the conflict ``the fundamental problem of human relationships''~\cite[p.~174]{dirac-81}.

I bring this up in a book about sentimentality because ``equal opportunity'' is the most popular promise in modern politics, and it cannot be kept. Every party makes it, and every parent who votes for it goes home and breaks it by helping with the homework. The promise survives only because nobody asks what it would cost to fulfill.

Nobody disputes that parents should love their children, or supposes that children must be treated identically in every respect. The difficulty is to honor the most intimate obligation we have without letting birth become a political destiny. It reaches into property, schools, neighborhoods, health, friendship, status, and the meaning of merit. Each generation is the environment of the next, and the conundrum follows from that.

\section{Two promises that cannot both be kept}

Call the first the \emph{promise to parents}. Parents may do for their own children what they would not do for strangers. They feed them first, comfort them at night, pay for their schooling, move house for a better school, introduce them to useful people, and save for their future. They are allowed to take pleasure in the fact that their children are better placed because of it. There is nothing embarrassing or selfish about this that a higher morality will someday remove. It is most of what being a parent means.

Call the second the \emph{promise to children}. No child should be shut out of a decent life by the accident of birth. None should go without a doctor, a safe home, a competent school, or a real chance to find out what he is good at, merely because the adults who raised him were poor or powerless.

Neither promise demands equal outcomes. Parents want their children to flourish. They do not want every child to earn the same examination marks and salary, and nobody else wants that either.

The promises collide the moment parental effort has effects outside the family, which is at once. One family buys a quiet room, a tutor, a safer street, a summer of unpaid exploring, or an introduction to an employer. The gain does not stay inside that family. A place at a selective school, a scarce internship, or a coveted job is a \emph{relative} good. When one child moves up, the standard against which the others are judged moves with him.

Take two children. Both are fed, loved, and competently taught. The first has parents who can pay for coaching, who know which applications matter, and who can carry her through a year of failure. The second has parents who love her just as much but work long shifts, know less about the institutions she must get through, and cannot finance a delay. Both have the legal right to apply. Nobody would say that their opportunities are equal.

So the conflict does not begin with the mansion and the corrupt appointment. It is present in small and blameless acts, and it becomes visible when those acts are added up and repeated over generations. The family says of one child that he is its own to protect. Justice says that another child, a stranger, counts for just as much. Both statements are true, and they cannot both be acted on to the full.

\section{Equal what?}

``Equal opportunity'' is a popular phrase because it means four different things, and each listener picks the one he likes.

The weakest meaning is \emph{legal}. No law reserves the position for a race, a sex, a religion, or a family. This is a real achievement, but an unlocked door is not yet an opportunity.

The second is \emph{fair competition}. The examination questions are not for sale, and the job is not kept for the nephew. This is the meritocrat's meaning. It is still incomplete, since a clean race between a trained runner and an untrained one has a predictable result.

The third is \emph{substantive}. People of similar ability and similar willingness should have similar prospects, whatever their origin. That was Rawls's point in asking for more than formal openness~\cite{Rawls1999}. A door is not really open to a child without shoes, time, bus fare, or confidence, who knows nobody able to explain how the building works.

The strongest meaning would \emph{equalize the conditions themselves}. Family wealth, health, neighborhood, vocabulary, and parental attention would all be neutralized, and perhaps natural talent as well. Here the ideal starts to devour the thing it was meant to serve. If every unequal influence has to be cancelled, parents are no longer simply bringing up their children. They are distributors of advantage, and they will have to be monitored.

A plainer vocabulary is more useful, and it has three terms. The \emph{floor} is the minimum below which no child may fall: food, safety, medicine, schooling, and a reasonable prospect of standing on his own feet. The \emph{field} is the range of difference that remains above the floor. The \emph{ceiling} is the point at which private advantage turns into public domination. Above it a family is buying access, exemptions, influence, and the power to write the rules for everybody else, and no longer merely comfort.

Most quarrels about equality become clearer once you ask which of the three is at issue. A person can defend a wide field and a high floor together. He can accept modest inheritance and oppose dynasties. What he cannot honestly do is say ``equal opportunity'' and leave the object unstated. Equal access to what, at what age, compared with whom? Until that is answered, a society can congratulate itself on formal equality while inheritance quietly settles the results. Or it can go to war against every difference and overlook the plain deprivation that really does wreck a child's chances.

\section{What a family is for}

Economists speak of the family as a device for moving resources from adults to minors, because money is the part that can be seen. But a small child does not need a bundle of goods. He needs somebody who notices that the bundle is missing, understands what is wrong, and is still there the next day. Brighouse and Swift have worked out in detail what goods of this kind are, and how much inequality a society ought to accept in order to keep them~\cite{BrighouseSwift2014}.

This is why the proposal to replace parents with institutions is harder than it looks. An institution can provide meals, medicine, and teaching, and it may be more reliable than a given family. It cannot provide the one person who feels answerable for this particular child. A child who knows that someone will still be there after the mistake, the illness, or the defeat has an asset no budget line supplies.

Parental partiality does three jobs. It \emph{protects}, by shielding the child from dangers he cannot yet judge. It \emph{develops}, by cultivating abilities and habits. And it \emph{represents}, in that parents speak for children who cannot yet speak for themselves. The third job is the least noticed and may be the most valuable. Parents interpret institutions, make the telephone call, write the letter, and ask for the second chance. A child whose mother knows how to contest a school's decision has an advantage that appears in no income table. An office can be neutral on paper and still be legible only to certain families.

It follows that a state serious about fairness must do more than move money. It must make its own institutions navigable for people who did not inherit the knowledge of how they work. That means clear procedures, independent advocates, and schools that can be trusted.

There is also a positive case for leaving parents a good deal of room. Imagine a society in which every act of parental care is examined as a possible breach of equality. Does the bedtime story confer an unfair cultural advantage? Must the family holiday be declared? Such a society would do more than limit inequality. It would turn affection into an administrative problem, and it would produce the fearful and performing parents such a regime deserves.

None of this places the household beyond justice. Families can be sites of violence, neglect, and petty tyranny, and children need public protection because private life is not always kind. The line that matters runs between \emph{preventing harm} and \emph{equalizing benefit}. A child must be protected from abuse. It does not follow that every difference in parental warmth must be evened out. Parents must not be able to buy their child a public office. It does not follow that they must be stopped from helping him prepare for a hard examination. The question to ask of any family advantage is how far its effects reach into the common world.

\section{Five inheritances}

The word ``inheritance'' suggests a will and a house. That is the smallest part of what passes from one generation to the next.

\emph{Material.} Money, land, a flat, a firm. The great service of wealth is that it lets a family wait, and this matters more than anything it lets them consume. It pays the rent during training, carries an unpaid internship, and makes a failed venture survivable. Of two equally able and willing young people, one can spend six months learning a skill without wages, and the other must take the first steady job he is offered. An observer who sees only the outcome will put it down to ambition. For the child with reserves a failure is information. For the child without, the same failure ends his schooling.

\emph{Bodily.} Health, nutrition, sleep, and treatment shape what a child can do long before any school assesses him. A child who misses lessons because of an untreated illness has made no bad choice. No society can make all bodies alike, but it can stop avoidable differences from becoming life sentences.

\emph{Cultural.} Every family hands on a culture. The trouble is that powerful institutions reward one culture in particular. Picture an admissions interview. One candidate has been trained since infancy to present curiosity as a virtue and to turn her experiences into a polished story. Another is just as curious and has been brought up to think that talking about oneself is boasting. The interviewer sincerely believes she is measuring motivation. What she is measuring is a manner of speaking acquired at the dinner table. An inherited confidence is converted into a credential. The credential brings a position, and the position becomes the dinner table of the next generation.

\emph{Social.} Names, contacts, and the experience of being expected in certain rooms. A network is more than a stock of favors. It is an information service that tells its members which trade is growing, which doctor to ring, and which mistake can be quietly put right. A society can live with unequal friendships. It cannot live with a connection that turns a public office or a licensed profession into a family estate.

\emph{Moral.} Families also teach what to expect of the world. A child learns from them whether institutions help or humiliate, whether effort is rewarded, and whether there is any sense in planning ten years ahead. This is no simple tale of confident rich children and cowed poor ones. Poor families can pass on extraordinary solidarity and resourcefulness, and rich ones can pass on anxiety, dependence, and the conviction that the world owes them a hearing. But the world pays better for some dispositions than for others.

These five matter because a policy aimed at only one will miss. A state could tax estates heavily and still let elite schools, segregated districts, unpaid internships, and closed professions reproduce the same hierarchy untouched. Inheritance has no single lever. It is an ecology.

\section{How a difference becomes a destiny}

The worst feature of inherited advantage is that it breeds. A small lead in the first years raises the odds of later success. Success brings resources for the next stage, and at the end the result is read as proof that the winner was simply more able.

The mechanism can be written down. Let $A_t$ be a family's effective advantage at stage $t$, taking money, health, knowledge, confidence, and connections together. Let $I_t$ be what the family invests, $R_t$ the return on that investment, and $E_t$ the effect of the surroundings. Then
\begin{equation}
  A_{t+1} = A_t + R_t(I_t) + E_t .
  \label{eq:advantage-recurrence}
\end{equation}
If the return grows with what a family already has, two families with a modest initial gap diverge quickly. For money alone the familiar version is
\begin{equation}
  W_{t+1} = (1+r_t)W_t + Y_t - C_t ,
  \label{eq:wealth-recurrence}
\end{equation}
with wealth $W_t$, return $r_t$, earned income $Y_t$, and consumption $C_t$. The equations predict nothing about any particular society. They are a reminder that a head start does not sit still. It earns.

The version that has nothing to do with money is the more important one. A child given early help with language likes school a little better. Liking school, she reads more, and reading builds vocabulary. With a better vocabulary the next lesson is easier, and good marks bring encouragement. The loop runs just as well backwards. Poor health makes attendance patchy, and a child who is often absent comes to think school pointless and stops trying. Low effort narrows the jobs open to him, and a precarious job is bad for health and housing. What looks at the end like a personal failure may be the last stage of a sequence that began before the person concerned could make any meaningful choice.

The compounding is social as well. Those chosen for prestigious posts become the teachers, managers, judges, and gatekeepers who decide what promise looks like in the next round. They reproduce themselves in complete sincerity.

One caution is needed. A tendency is not a fate. Rich children fail and poor children rise, and no statistic strips an individual of his agency. The lesson is to distrust the meritocratic fairy tale that starts from the result and reasons backward to the character of the winner.

The practical conclusion concerns timing. A grant at twenty-five cannot undo an unequal childhood, and a public library cannot give back the sleep lost in an unsafe home. Whoever wants to widen opportunity has to act early, before advantages have hardened into certificates and disadvantages into a child's description of himself.

\section{The child as person and as competitor}

A child's situation can be looked at in two ways. \emph{Absolutely}: does he have enough food, safety, schooling, health, affection, and room to grow? \emph{Relatively}: where does he stand among the others who compete for the same places?

The absolute question comes first. A society can be unequal while nearly all its children live securely. It can also be admirably equal because everybody is poor, sick, and at the mercy of officials. Equality by itself tells you nothing about whether anyone is doing well.

But relative position cannot be dismissed as envy, because standing decides what a person is allowed to do. Three kinds of goods should be told apart. \emph{Consumption goods}, such as the bigger house and the better holiday, are comparatively harmless. \emph{Positional goods}, such as the place at the selective school, are serious, because one child's gain is another's loss. \emph{Power goods} confer the ability to shape the rules by which places are handed out, and they threaten the equality of citizens as such.

It follows that the response to inheritance should not be uniform. Nobody can defend a system in which a family buys a court ruling or a clause in a statute. Nobody can defend forbidding a parent to give a child a violin and a quiet room. The work of judgment lies between these two, with private schools, expensive tutoring, family firms, and gifts of property.

Scale changes the nature of the thing. A modest inheritance gives a young adult stability. A vast one can buy newspapers, political access, and the competition. So ``the parents earned it'' cannot be the end of the discussion. One must also ask what the transfer will do to the world the next generation has to live in.

\section{Merit, with less arrogance}

Meritocracy is offered as the reconciliation. Parents may help at home, provided that public positions go to talent and effort. Everyone is free to compete, the best wins, and the result is legitimate.

There is truth in this, and those who sneer at merit should recall the alternatives. Work, skill, discipline, and courage are real. An institution that refuses to recognize them soon loses its competence, and then the public's trust.

But meritocracy easily becomes a way of hiding how merit is manufactured. An examination can be honestly marked and free of corruption, and still reward ten years of preparation, a certain way with language, and coolness under pressure. That does not invalidate the examination. It shows that clean selection and equal preparation are different things. The competition inspects finished products and is then described as a pure measure of effort, although effort itself was helped or hindered long before. The literature on equality of resources and on luck makes the same point~\cite{Dworkin2000,Anderson1999}. Effort is real enough. The conditions under which it shows and pays are inherited unequally.

The remedy is to keep merit and take the arrogance out of it. Merit is a good reason for giving someone a job. It is no proof that he deserves everything the job then brings within his reach. A winner can have earned his place and still be in debt to teachers, to public roads, to his mother, and to luck.

Institutions can act on this. Selection procedures can look for ability in the various forms in which different backgrounds express it. There can be several routes into a profession, apprenticeships and second chances among them, and not just one expensive credential. Above all, it must be possible to lose without being ruined. Where the winner takes everything and the loser is marked for life, a rich family's ability to cushion failure becomes overwhelming. A secure floor is therefore a friend of competition. It lets people attempt hard things.

\section{The state is already in the room}

The debate is usually staged as family against state, as though the state stood outside and stepped in only when private conduct went too far. In fact the state is inside the process from the beginning. It defines property, enforces contracts, certifies schools, zones neighborhoods, writes inheritance law, and decides which diploma opens which door. A house has a market price because a legal order recognizes ownership. A degree is worth something because employers and licensing boards say it is. There is no natural distribution lying underneath the institutions. The state shapes inheritance whether it likes it or not, and the only question is whether it does so openly and for reasons it can defend.

This cuts against two opposite illusions. The first is that inherited advantage is a private matter with which public institutions have nothing to do. The second is that public control comes free of cost and brings no new forms of power with it. A tax office or a school board can flatten one hierarchy and build another. Officials are human, and they have nephews too.

So fairness means constraints on public power as well as private. Rules should be general, visible, and open to challenge. A child should need neither a powerful parent nor a powerful bureaucrat in order to be heard.

Universal services matter here. A service reserved for the poor tends to become a poor service, underfunded and stigmatized, while the comfortable buy their way out. A good service used by everybody creates a common floor and a common experience. Schools, clinics, libraries, and public transport narrow the distance between families without an inspector in every living room.

The state can also change the game, and not merely compensate the losers. It can break monopolies, open the professions, publish how institutions perform, and limit the conversion of money into political voice. None of this makes families equal. It stops the most fortunate families from closing the public routes behind them.

The test for intervention cannot be whether something benefits a child, or whether it creates a difference. All good parenting does both. Better questions are these. Does the advantage deprive others of something basic? Does it buy a scarce public position? Is the resulting hierarchy open to outsiders? Does it grow from one generation to the next? And can it be regulated without wrecking a valuable relationship or creating a worse concentration of power?

\section{Four tempting solutions}

\subsection{Abolish inherited advantage}

The purest proposal is to wipe out every advantage of family background. All children would get the same resources and the same schools, and no private coaching, contacts, or property would be allowed.

Its appeal is plain, and it exposes an awkward fact. Most people who are enthusiastic about equal opportunity mean the opportunity for their own children to compete from a position that is already favorable.

But total equalization would have to invade the family. A book, a conversation, or a parent's free afternoon would become a suspect transfer. It would need an authority able to inspect and compare the contents of households, and such an authority would soon be more arbitrary than the inequality it was set up to cure. The proposal mistakes the impossibility of perfect neutrality for the uselessness of partial remedies.

\subsection{Leave families alone}

The opposite proposal protects family and property almost without limit. What was honestly acquired may be freely given, and if the children of the fortunate do better, that is the price of freedom. This is the logic of the entitlement theory of property~\cite{Nozick1974}.

It is right that people have projects, attachments, and a claim to what they have built. It is right, too, that an equal distribution maintained by surveillance and force is no improvement.

It fails at the point where private advantage becomes public power. A parent leaving a modest house is in a different moral position from a family handing down control of a dominant platform, a media empire, or a regulated monopoly. That a transfer is legal does not tell us whether it leaves other people citizens or makes them dependents. The view also treats the market's starting point as self-evident, when the value of most assets rests on roads, courts, and collective decisions that the owner did not make.

\subsection{Redistribute afterwards}

A third proposal lets families create differences and then compensates with taxes and transfers. This has the great merit of staying out of the household. It is also on firm philosophical ground, since the rules that make ownership and bequest possible are themselves public arrangements and no facts of nature~\cite{MurphyNagel2002}.

Redistribution is indispensable, and it comes late. A transfer at twenty cannot buy back early security, early language, or the confidence of having been expected to succeed. A check does not open a closed network. There is a political catch as well. If advantage is allowed to pile up and is corrected only at the end, its beneficiaries may by then be strong enough to dismantle the correction.

\subsection{Make losing bearable}

The fourth proposal keeps the differences and lowers the stakes. If medicine, schooling, housing, and a dignified life are not prizes reserved for winners, inherited advantage matters less. A child can lose a contest and keep his future.

This is the most promising of the four, because it deals with the relative problem without requiring identical households. It does not make opportunity equal. It makes inequality less destructive.

But some goods are positional by nature. A research institute has only so many seats, and a good street only so many houses. Somebody must still decide who gets them, and how much inherited influence a republic can stand.

Each proposal holds one piece of the truth. Equalization guards against hierarchy of birth. Family liberty guards intimacy. Redistribution protects those who start behind, and lower stakes protect those who stumble. A workable order needs all four in unequal amounts.

\section{The pie is not fixed}

Arguments about inheritance tend to assume a fixed quantity to be divided. For scarce positions that is correct. For an economy as a whole it is false.

Families and institutions help to produce goods as well as divide them. A stable legal order makes long-term investment possible. A school creates abilities that did not exist. Research yields knowledge that benefits millions who never paid for it.

Growth does not cure injustice automatically, since a bigger pie can be eaten by fewer people. Nor is every inequality excused by the claim that it promotes growth. But policy has to be judged on both counts, production and distribution.

Incentives are part of the moral landscape, because the total available to everyone depends on them. If people expect that every extra effort will be confiscated, some will work less, invest less, or leave. If they expect that nothing can be passed to their children, some kinds of saving and long-term care will wither. People do work for meaning, duty, and pride as well as money, but they are not saints.

The opposite expectation is just as corrosive. Where inherited wealth guarantees the children's position whatever their competence, the result is complacency, rent-seeking, and political defense of obsolete arrangements. A society that wants invention has to protect the right to create, and deny the right to rule for ever on the proceeds of a past creation.

History adds a warning against romantic shortcuts. The great reductions of inequality have mostly followed war, revolution, state collapse, and plague, the ``great levelers'' of comparative history~\cite{Scheidel-2017}. Careful moral design has seldom achieved as much. Destruction does level a landscape. That shows how stubborn inherited hierarchy is, and it recommends nobody's catastrophe. Reform should therefore be patient. It should raise the floor, widen the entrances, remove the rents, and check concentrations of power before they set hard. A fairer distribution that keeps productive capacity intact is worth far more than a brief equality bought with ruin.

\section{What can actually be built}

One more fact complicates the picture. Parents produce a public good. They raise the future workers, nurses, taxpayers, inventors, and neighbors, bearing most of the cost while society collects much of the return. That is the standing argument for public support of families through schools, clinics, parental leave, and housing. It is cost-sharing for a common investment, and nobody should mistake it for charity. Where good schools and clinics are open to all, parents can give love, attention, and guidance without having to buy a whole parallel infrastructure. The family remains partial. The background against which it is partial becomes less unequal. Shared institutions also give children a common world, in which future citizens at least catch sight of one another.

On this basis a structure can be sketched. It has six load-bearing parts.

\emph{A high and unconditional floor.} Food, medicine, safe housing, protection from violence, and good schooling, for every child. ``Unconditional'' is the important word. A child should not have to prove that his parents chose wisely before he is given the means to make choices of his own.

\emph{Early help without surveillance.} Because advantage compounds, the first years are the cheapest time to act. Maternal care, good nurseries, and early language teaching should be offered as ordinary services for everybody, so that accepting them is no badge of failure. The state must come down hard on abuse and neglect, and it must not treat every unconventional household as a defective one.

\emph{Several ways in.} Transparent recruitment, paid apprenticeships beside costly degrees, and no credentials whose only function is to protect those already inside. Standards need not fall. They should measure what the job requires, and stop measuring inherited familiarity with the way to the job.

\emph{A ceiling on conversion.} This is the most important of the six. Money buys education. Education yields credentials, credentials lead to office, and office makes the rules that protect the money. Every link may be legal, and the chain produces a hereditary ruling class. Public action should aim at the links, which are monopoly, regulatory capture, opaque admissions, purchased political access, and conflicts of interest. Private life need not be equal. Private advantage must be denied a public veto.

\emph{Inheritance as a social dividend.} No fortune is made in a vacuum. Roads, courts, science, stable money, and public health all contribute to the value of private assets. To return a share to the common fund is to acknowledge that co-authorship, and it insults no family. Inheritance tax is the obvious instrument, and the case for it grows where returns on capital outrun ordinary incomes~\cite{Piketty2014}. But it needs distinctions. A small bequest that keeps a young family out of debt is not a conglomerate. A tax that ignores liquidity forces the sale of the farm, and a tax riddled with exemptions spares only the largest fortunes. It should also be remembered that taxing the estate does not undo what the money did earlier, when it paid for tutors, addresses, and networks. And the legacy to worry about is control more than comfort. To inherit a productive asset is one thing. To inherit a bottleneck through which other people must pass is dangerous, and competition law and campaign finance rules may matter more here than the estate tax.

\emph{A culture of bounded success.} Rules will not hold in a society that worships dynasties. A parent may want his child to do well, and the child is not thereby entitled to win every public contest. A family may keep its business without owning the market. A rich man may speak in public, and his money should not buy him a louder vote.

\section{The remainder, and a test of honesty}

When the floor is high, the routes are open, and the concentrations are checked, inequality will still remain. Some parents are more patient, better informed, or luckier. Some children are born with temperaments and bodies that suit the institutions of their time. No public order can reduce all this to a common measure without destroying the variety of human life.

One may call the remainder unjust, on the ground that no child earned his start. But the word conceals a choice. To remove every unearned difference, how much privacy must be given up, and how much parental freedom? What new authority would decide which advantages count? A moral complaint is not yet a policy. The opposite error is to call every surviving difference ``deserved'' because it passed through a market.

The right attitude is disciplined discomfort. We should be able to say that a difference is morally arbitrary, that it is permitted for a reason, and that the reason has limits. An inequality can be allowed and not celebrated.

The problem also works as a lie detector. People defend their own children's advantages as love and call other people's corruption. They praise competition when they are ahead and ask for protection when they are behind. They invoke merit when the result suits them and structural injustice when it does not. Loudest of all are those who preach equal opportunity from platforms and buy positional advantage in private. They are sincere in both, and they never set the two side by side.

No decree abolishes this bias, but institutions can make it harder to act on. Good universal services stop the rich from turning every public good into a private purchase. Open admissions leave less room for favoritism. Independent review gives the outsider a means of fighting back, and conflict-of-interest rules keep family loyalty out of public office.

An honest society can still be one in which some people receive more from their parents than others. Its honesty lies in stating the reasons for permitting this, showing the cost to others, and making sure the advantage never hardens into exemption from the common rules. Such a society will be unequal. It will be less likely to mistake a hierarchy of families for a natural hierarchy of persons.

\section{Living with an unsolved problem}

The conundrum starts from two ordinary thoughts. Parents should help their children. Children should not be sentenced by their birth. The first leads to intimacy, sacrifice, and the freedom of families. The second leads to equal citizenship and the refusal to let private fortune become public fate.

The conflict is structural. It survives every change of regime and vocabulary, because children are formed in particular relationships and justice addresses them as equal members of one world. The same gift is love inside the family and a handicap for someone outside it. The same network that finds one child a vocation closes a profession to another. Dirac saw that two good principles exclude each other~\cite{dirac-81}. The task of institutions is to decide where the collision can be softened and where it must simply be borne. They must also decide where private advantage has to be stopped before it becomes public domination.

Pushed to the limit, either ideal destroys something valuable. Total equalization needs an intrusive power that would flatten family life and with it the motives for care and creation. Unlimited inheritance lets advantage accumulate until citizenship itself is unequal.

What can be defended is management. Set a high floor early enough to interrupt the compounding. Make the public world navigable for those who did not inherit the map. Keep the routes open. Tax and regulate transfers that confer command, and leave alone those that confer ordinary security. Preserve a wide sphere in which parents may love, teach, and provide. And treat what inequality remains as a cost that has to be justified, never as a natural order to be admired.

No such settlement is final. A school reform opens one door and the housing market shuts another. A tax reduces concentration and a professional network quietly restores it. Each generation has to check again whether the floor, the field, and the ceiling still hold.

Dirac's small provocation is valuable because it forbids the comfortable pretense that one principle can simply defeat the other. That is the sentimental answer on both sides. The egalitarian has it when he promises equal chances without mentioning the price, and the conservative when he calls every inheritance earned. Resignation would be the wrong response to a problem that cannot be solved. The better one is to build institutions that keep the contradiction from turning cruel, and to watch that our own partiality does not pass itself off as justice.


% ===================== Chapter file ch08 =====================
\chapter{Government Between Power and Liberty}
\label{chap:government-and-liberty}

\begin{quote}
\itshape
The central political problem is not how to make government strong enough to
do everything, nor how to make it weak enough to do nothing.  It is how to
make public power strong enough to protect a private sphere in which no
government is master.
\end{quote}

\section{The old choice and the modern temptation}

Every political order is an arrangement for using force. It settles who may make rules, who may tax, who may punish, who may command the army, and who may change the rules themselves. A government can be kindly, elected, and hedged about with limits. It is still an institution that can make you do what you would otherwise not have done. Nobody can talk sensibly about government without talking about liberty.

The usual picture is a tug of war. On one side stands the strong state, which is orderly, decisive, and competent. On the other stands the free individual, who speaks, worships, trades, moves, refuses, and lives by his own lights. Each exposes the weakness of the other. Liberty without public power is the liberty of the strong to push the weak around. Public power without liberty turns protection into tutelage and order into fear.

The matter is harder than a tug of war, because government has two jobs that pull against each other. It must restrain private violence and fraud, and it must restrain itself. The police have to be strong enough to stop a gang from ruling a neighborhood, and not so strong that they become a gang in uniform.

Democracy comes into this as a method of handing out public power and taking it back. It is not the same thing as liberty. A majority can vote to censor a minority, seize its property, or ban its religion. Nor is it the same as good government. Voters can choose incompetents, reward corruption, and send the bill to people not yet born. Democracy decides who governs and how they can be removed. Liberal institutions decide what even a government with a majority may not do.

There is a sentimental view of democracy, and this chapter rejects it. On that view the people are wise, their will is one, and an election reveals it. None of this is true, and the first part of it can be proved false by mathematics. Yet I shall argue that for anyone who puts individual liberty first, constitutional democracy is still the best form of government on offer for large, mixed, modern societies. The reason is unflattering to the voters. Democracy is best because power in it is dispersed and can be replaced, and because it can be contested in public. The citizen keeps a standing right to say no.

This is a defense by comparison. I do not claim that democracy decides better than an enlightened monarch or a council of experts. Either may now and then decide more intelligently. But no such regime offers a reliable way of keeping the ruler enlightened, benevolent, and willing to step down. A government that the governed can correct is less dangerous, even when it is noisy, slow, and often wrong.

The defects of democracy are visible and can be fought over. Those of concentrated power stay hidden until it is too late to correct them. Democracy does not manufacture a general will. It is machinery for keeping any temporary will from becoming a master nobody can question.

\section{Government is organized coercion}

A tax is not a subscription, and a court judgment is not advice. Behind every criminal statute stands a policeman. To say so is no argument for anarchy. It is a matter of honest description. A government that calls everything it does ``help'' is hiding the compulsion. A critic who calls every rule oppression is hiding the private violence that rules prevent. The same law is an unwelcome interference to one man and a shield for his neighbor.

Take property. The state that protects my house from a thief also forbids me to turn it into a tannery. There is no neutral point at which government has simply gone away. The legal order is always giving powers to some people and denying them to others.

``Strong government'' therefore needs to be taken apart. A state can be strong in capacity and weak in arbitrary power. It collects taxes efficiently, runs honest courts, and repairs the roads, and it cannot silence an opponent or take a house without judicial review. Another state is feeble in capacity and strong in arbitrary power. It cannot deliver the mail, and it can arrest anyone it likes. The line that matters separates general power exercised under rules from discretionary power, where an official decides case by case who is protected and who is exposed.

The rule of law is the device that converts the first kind of strength into something tolerable. Rules must be public, general, and reasonably stable, and they must bind the rulers. That does not make every law just. It makes law something that can be criticized and revised, and it takes away the ruler's personal will as the sole source of obligation. Nobody can plan a life on an official's good mood. A limited government is by no means an idle one. Its powers are bounded by institutions that citizens can understand and challenge.

\section{Berlin's two liberties}

Isaiah Berlin supplied the vocabulary~\cite{berlin-liberty}. \emph{Negative} liberty is the area within which I can act without deliberate obstruction by others. Its question is how many doors stand open and who is blocking the rest. \emph{Positive} liberty asks who governs me, and what I need in order to be my own master.

Berlin did not simply cheer for one against the other. A man who is starving, illiterate, or terrorized by the local strongman has few doors open, however little the state interferes with him. Negative liberty presupposes a public order that protects people from assault, fraud, and private tyranny.

His warning was aimed elsewhere, and it bears directly on the theme of this book. Positive liberty can be taken over by an authority that claims to know what people would choose if only they were rational, educated, and properly liberated. Dissent then becomes a symptom. The government is not coercing the dissident, it is educating her. It does not censor the newspaper, it protects the true interests of the people. The official becomes the interpreter of your freedom, and obedience is presented as a higher form of self-rule. Every paternalism of the last century has spoken this language, and the benevolent ones speak it most fluently.

The liberal answer is not to abandon common purposes. It is to deny that a common purpose entitles the state to occupy the whole of life. The private sphere covers conscience, belief, inquiry, family, art, peaceful speech, and movement, and it includes the freedom to be eccentric and wrong. People become individuals by trying things no official would have designed. A government that protects only the approved forms of autonomy is not protecting autonomy.

Democracy matters here because it keeps the boundary disputes open. It does not settle in advance where one man's freedom ends and another's protection begins. It provides a public arena for the quarrel, lets the loser stay in the game, and offers ways to change a rule without first overthrowing the state.

\section{Why the strong hand tempts}

The wish for strong government is not a primitive mistake. It grows out of real failures of uncoordinated action. No city gets a transport network if every landowner has a veto. No country defends itself if conscription is optional. Epidemics, floods, bank runs, and invasions call for decisions faster than a debating chamber can produce them.

There is the further attraction of decisiveness. When parliament is deadlocked and public debate has sunk to repetitive insult, a leader who promises simply to act has a strong pull. People will give up a good deal of freedom for the promise of order.

But a decisive ruler is not thereby a correct one. A strong government can suppress the information that would have corrected its policy and punish the messenger who reports the failure. It can stretch a temporary emergency into a permanent mandate. The capacity to act is an asset only where something exists that can stop an action that has turned dangerous.

Then there is the question of the successor. Suppose the present ruler is wise. He will die. He may fall ill, or appoint his nephew, or be replaced by a clique that understands repression better than administration. A constitution is an attempt to solve the problem of the unknown successor. It gives up some freedom of action for today's government in order to preserve the freedom of citizens under tomorrow's.

The same holds for elected governments. A party with a majority is tempted to rewrite the electoral law, tame the courts, lean on the press, and fill the independent offices with loyalists. It explains that it must be allowed to carry out its mandate. But a mandate to govern is no mandate to make sure that one never loses again.

``Government of the people'' can mean that the government acts for the people's good. It can also mean that the people can dismiss the government. The second meaning flatters rulers less and matters far more. A government can sincerely believe that it serves the people and still become their master.

\section{The oldest objections}

Democracy has been criticized for as long as it has existed, and the classical critics were no mere snobs.

Plato's objection in the \emph{Republic} is the ship~\cite{plato-republic}. Nobody would let the passengers vote on navigation. The question is who knows how to sail. Democratic competition rewards the man who can charm the crowd, and he need not understand the good of the city. Plato feared, too, that democracy would turn freedom into license. When every restraint is an insult to equality, authority dissolves, and a demagogue arrives to protect the people from the elites and gathers all power in the people's name. Too much freedom prepares the ground for tyranny.

Aristotle was more empirical. He classified regimes by whether they serve the common interest or the rulers' own. Rule by the many can be a legitimate constitution, and it can also degenerate into the poor majority plundering everyone else. His remedy was mixture, an order in which no group can dominate for good~\cite{barnes-aristotle}.

Two modern insights are already present. Political equality does not imply equal competence. And a regime can be popular in form and oppressive in substance.

The democrat's reply to Plato does not deny that knowledge matters. It asks who is to identify the knowers, who is to check them, and how they are to be corrected when they are wrong. Expertise is indispensable in medicine, engineering, and law. It does not follow that experts should have the last word on the purposes their skills are to serve.

Madison supplied the modern version of Aristotle's mixture~\cite{Madison1787}. Faction is permanent, since it could be abolished only by abolishing liberty. So its effects must be controlled. The republic is enlarged so that a local majority becomes a national minority, and its institutions are divided so that a coalition must win in more than one place. Delay is a defect when speed is needed. It is a virtue when it gives minorities time to organize and the public time to read the bill.

Representation, which scale makes necessary, creates agents with interests of their own. A representative may come to know more than his voters. He may also come to depend more on his donors, or to be more loyal to his party than to his district. Elections discipline him somewhat. So do public records, open hearings, a free press, auditors, and courts. An election held in the absence of all these is a ceremony. Liberal democracy is a competitive ecology, and it amounts to much more than an event called election day.

\section{The mathematics of disappointment}

The first difficulty arises before anyone has said a word about stupid voters. Let three voters rank three options:

\begin{center}
\begin{tabular}{c|c}
Voter & Ranking \\
\hline
1 & $A \mathrel{>} B \mathrel{>} C$ \\
2 & $B \mathrel{>} C \mathrel{>} A$ \\
3 & $C \mathrel{>} A \mathrel{>} B$
\end{tabular}
\end{center}

Each voter is perfectly consistent. But put the options to the vote in pairs. $A$ beats $B$ by two to one. $B$ beats $C$ by two to one. And $C$ beats $A$ by two to one. The group prefers $A$ to $B$, $B$ to $C$, and $C$ to $A$. This is Condorcet's paradox~\cite{Condorcet1785}. The group is not short of information. Its majorities contradict one another, and no single ranking fits them.

The political consequence is that the order of voting decides the result. Start with $A$ against $B$ and you get one winner, and start elsewhere and you get another. Whoever sets the agenda holds power without commanding a single extra vote, be he the committee chairman, the party whip, or the government that controls the parliamentary calendar. Procedure is no technicality.

Kenneth Arrow generalized the trouble~\cite{Arrow1963}. Suppose we want a rule that turns individual rankings of three or more options into a social ranking, and we ask for four modest things:

\begin{enumerate}
\item \textbf{Unrestricted domain.} The rule works for any set of individual rankings.
\item \textbf{Pareto efficiency.} If everyone prefers $A$ to $B$, so does society.
\item \textbf{Independence of irrelevant alternatives.} Society's ranking of $A$ against $B$ depends only on how individuals rank those two.
\item \textbf{Non-dictatorship.} No single person's ranking always prevails.
\end{enumerate}

Arrow proved that no rule satisfies all four. The popular summary, that democracy is impossible, is too strong. The theorem concerns one class of rules and one demanding set of conditions, and it does not say that elections are irrational. What it does say is sobering enough. No neutral device exists that preserves every reasonable notion of fairness and also yields a coherent collective ranking. Every system gives something up. And no voting rule can be defended by the claim that it perfectly represents the people. Representation is always representation under rules, and the rules shape the result.

There is a bitter joke in the theorem. Dictatorship does well by its formal standard, since the dictator's ranking is always coherent. But a coherent will has no claim on us for being coherent. It can be consistent and also cruel, ignorant, and impossible to remove. Democracies accept some incoherence as the price of giving nobody the last word.

A third result completes the set. Gibbard and Satterthwaite showed that any reasonable rule for choosing one winner among three or more candidates can be manipulated~\cite{Gibbard1973,Satterthwaite1975}. Under every such rule a voter will sometimes do better by misstating his preferences. The man who likes $A$ best and votes for $B$ to keep $C$ out is no liar. He is responding to the rules. It follows that no ballot simply reports ``the true will of the voters.'' Plurality voting, ranked ballots, approval voting, and the Borda count each transform preferences in a way of their own. To choose among them is to choose which distortion one can best live with. It is a political choice.

So the general will is not hidden somewhere, waiting for a better voting system to uncover it. In general it does not exist.

\section{The ignorant voter}

Anthony Downs began with a puzzle~\cite{Downs1957}. In a large election the chance that your vote decides the outcome is next to nothing. Voting costs time and effort. A strictly self-interested person would stay at home. Yet people do vote, out of duty, loyalty, and the wish to be counted. An election is a ritual of membership as well as a decision procedure. A citizen votes to say ``I count.''

The companion problem is \emph{rational ignorance}. If my vote almost certainly changes nothing, I have no reason to spend three hundred hours mastering tax law, energy systems, and monetary policy. So I use shortcuts: the party label, a trusted group, the look of the candidate, a story about how the world works. Shortcuts are sensible, since nobody can verify everything. They are also easy to exploit. Parties offer identity in place of evidence, media outlets live on outrage, and platforms spread a vivid falsehood faster than a dull correction.

Bryan Caplan goes further. Voters are worse than uninformed, he argues, because they are biased in systematic ways, above all about economics. Being wrong costs the individual voter nothing, and believing what feels good pays at once~\cite{Caplan2007}. Jason Brennan draws the conclusion that democracy lets people wield power over others without knowing what they are doing~\cite{Brennan2016}.

These arguments are largely correct, and a defender of democracy who denies them is being sentimental. The right response is to ask, compared with whom? Experts have guild interests, professional fashions, and reputations to protect. The useful question is which arrangement brings error to light and allows it to be corrected. Which group is free of error is not worth asking, since none is.

On that question democracy does respectably, provided we stop treating elections as referendums on policy detail. A citizen may not understand how a budget is built. He can still tell whether the hospital works, whether officials lie, whether opponents are being locked up, and whether power changes hands peacefully. An election is a crude instrument of dismissal, and that is the use it should be put to.

The argument that crowds are wise holds only under conditions. People must have partly independent information, and their errors must not all point the same way. Propaganda and tribal loyalty destroy both conditions. When everyone drinks from the same poisoned well, a large crowd is merely wrong with great confidence. So a democracy needs knowledge to enter politics by more than one door: independent statistics, an investigative press, a professional civil service, opposition parties, courts, and protected whistleblowers. A parliament can be ignorant while the system as a whole contains knowledge that the parliament cannot erase. Hayek made the parallel point about markets. No planner possesses the knowledge that is scattered among millions of buyers and sellers, and prices coordinate some of it~\cite{Hayek-45}. Markets have their own failures, and politics has to set their rules.

The danger is greatest when one election hands over concentrated power. If a single victory delivers the courts, the media, the police, and the electoral commission, ignorant voting stops being a tolerable weakness in a dispersed system and becomes the gate to permanent capture. The more power an election transfers, the stronger the walls around it have to be.

\section{The majority and the individual}

Majority rule treats citizens as equals. That is its charm, and it tells us nothing about whether the majority is right. A majority can restrict the speech, property, movement, and worship of a minority and load it with burdens it has no power to resist.

The tyranny of the majority need not be dramatic. It can be the daily pressure to conform, the decision that certain opinions are not respectable and that certain people must keep proving their loyalty. Tocqueville described a democratic \emph{soft despotism} in which citizens are not terrorized and are gradually reduced to dependence, conformity, and supervised comfort~\cite{Tocqueville1835}. He wrote this nearly two centuries ago, and it reads like a description of the present.

The liberal remedy is to take some matters out of the reach of ordinary majorities. These are speech, conscience, fair trial, association, and equal standing before the law. Courts, second chambers, federal structures, and supermajority rules exist for the purpose.

The remedy has its own paradox. Who authorized the judges to overrule the majority? Judicial review can decay into government by lawyers, with ``rights'' stretched to cover whatever policy the bench prefers. An honest position admits both dangers. Not every majority decision is valid, for then my liberty would depend on my popularity. Judges are not infallible either. Rights should be clear, courts should give their reasons, and legislatures should be able to answer by amendment. Liberal democracy is a mixed regime. It trusts citizens enough to choose governments and not enough to give any majority unlimited power. It trusts judges enough to guard the private sphere and not enough to make them the authors of the political order.

\section{How representation is bought}

\emph{Pork barreling} means steering public money to a particular district or interest in exchange for support. It is often called corruption, and the truth is subtler and more depressing. Voters do not elect a national average. They elect someone who is expected to come home with a road, a hospital wing, or a subsidy. When every representative behaves in this way, parliament passes a pile of locally popular projects that is wasteful taken as a whole, and the taxpayers who finance it cannot trace their bill to any single vote.

This is the logic of concentrated benefits and diffuse costs~\cite{Olson1965,BuchananTullock1962}. The beneficiaries are few, organized, and grateful. The payers are many and scattered, and each loses too little to notice. The benefit is photographed at the ribbon-cutting. The cost disappears into taxes, debt, inflation, and the things that were not built. A project can make excellent political sense and be a waste for the country. With logrolling, in which I vote for your bridge and you vote for mine, the waste turns into a system.

Banning earmarks does not remove the spending. It moves into ministries and procurement offices, where it is harder to see. Disclosure, audit, and published costs are more useful. They do not abolish self-interest, which no regime can do, but they force it to justify itself in daylight.

\emph{Rent-seeking} is the same thing on a larger scale. Political power is used to obtain an income that open competition would not provide. The instrument may be a subsidy, a licensing barrier, a tariff, a tax exemption, a regulation that happens to cripple a competitor, or a contract awarded to a friend. It is always defended in the name of jobs, safety, or the national interest. Sometimes the defense is even true. A firm that can concentrate on one parliamentary committee outweighs a million consumers who would each save a few euros. In the end, success in the market depends on controlling the rulebook more than on serving customers.

The end state is oligarchy. Political, economic, bureaucratic, and media elites form a lasting network that runs public institutions for its own benefit. No conspiracy is required. People with the same schooling, vocabulary, and interests select one another quite naturally. The parties can go on competing while the range of permitted policy stays inside limits the network has set. The countermeasures are familiar: visible lobbying, contestable contracts, limits on the revolving door, a press that follows the money, and counter-organization by those who pay. The countervailing groups develop rent-seeking habits of their own, of course.

Elections add a problem of time. Politicians are rewarded for benefits that arrive before the next vote and punished for costs that are visible now. Maintenance, pension reform, and research lose out. Future citizens have no vote. Autocrats are no better, since they build monuments and conceal failures. But democracy needs deliberate counterweights: audit offices, fiscal rules, and special caution toward anything irreversible. One caveat applies. ``Future generations'' are also the favorite clients of guardians who would like to escape present accountability. The future gives experts no blank check. The best rule is that a policy which closes off options needs a stronger justification than one which can be repealed.

\section{How democracies die politely}

Parties simplify politics. They bundle issues, recruit candidates, and give the voter a rough map. They also turn every question into a test of loyalty. Polarization differs from disagreement. It begins when the other party is seen as an illegitimate enemy whose victory would justify dropping the usual restraints. Citizens then accept censorship, gerrymandering, and selective prosecution, so long as the other side is the target. The party has become more important than the rule that allows a party to lose.

Democracy needs an opposition that can win. It needs a government that does not treat criticism as treason, and an opposition that does not treat governing as usurpation. Constitutional loyalty means accepting that one's side may lose and remain free to try again.

Democracies seldom die in a coup nowadays. They are hollowed out by people who were elected~\cite{Diamond2019,HaggardKaufman2021}. The electoral commission is captured, the journalists are intimidated, friendly businessmen are rewarded, the courts are packed, and the districts are redrawn. Each step is lawful, popular, or ``necessary to stop a dangerous opponent.'' Elections continue. Losing them has become impossible.

Backsliding usually starts from real grievances over corruption, insecurity, immigration, or official arrogance. A leader promises to break the system and describes every restraint as the work of an unrepresentative elite. Some of the restraints really are unjust. But the cure on offer replaces dispersed unaccountable power with concentrated unaccountable power. And established parties that have answered real grievances with contempt are in no position to be surprised when the demagogue arrives.

Emergencies are the other route. A state may have to restrict movement, requisition goods, and collect data in a hurry. The danger is that the emergency resets the baseline. The temporary measure leaves behind surveillance systems, administrative habits, and commercial interests that resist removal, and the government learns that a frightened public will accept what a calm one would refuse. The principle is that emergency powers must be necessary, proportionate, public, limited in time, and open to review, and they should be written down before the emergency. A state that says ``trust us indefinitely'' is asking citizens to give up the faculty by which they could judge whether trust is still deserved.

\section{The lot}

Election is not the only democratic device. Ancient Athens filled most of its offices, its council, and its huge juries by \emph{sortition}, random selection among citizens. Athens was no model of inclusion, since women, slaves, and resident foreigners were shut out. The principle is what interests us~\cite{dow_aristotlekleroteria_1939}.

The Greeks regarded election as the aristocratic method. It picks out the prominent, the rich and eloquent and well connected. The lot gives every citizen the same chance of governing, and in that sense it is the more democratic of the two.

Its advantages are real. It needs no campaign money and builds no career. It leaves no patronage network behind. It brings into the room people who would never stand for office. Its weaknesses are just as real. The lot does not select for competence or courage. An assembly of amateurs can be steered by whoever briefs it, and its members go home before the consequences arrive.

Serious republics therefore mixed the methods. Athens elected its generals and drew lots for nearly everything else. Venice chose its Doge through an elaborate alternation of lot and vote. Florence drew names from bags. Modern citizens' assemblies revive the idea~\cite{Sintomer2016,Bouricius2013}. A random sample of citizens studies a question, hears the rival experts, deliberates, and reports to parliament or the public. An assembly of this kind has no constituency to bribe with pork and no donors to please. Its use lies in advising, auditing, and embarrassing the professionals, and it should not replace the government. The real virtue of the lot is that it breaks into the political career, which otherwise perpetuates itself, and reopens the path from ordinary citizenship to public authority.

The lesson extends to all electoral reform. Plurality voting gives clear governments and wastes votes. Proportional representation reflects opinion and blurs responsibility. Ranked ballots help with spoilers and do not escape the impossibility theorems. Term limits break up incumbency and strengthen the unelected staff. Every reform exchanges one fault for another. The honest reformer names the exchange and does not promise that a technical change will finally reveal what the people want.

\section{The constitution as a machine against concentration}

A constitution is usually described as a statement of values. It is better understood as a technology for preventing the concentration of power. Separation of powers obliges institutions to cooperate. Federalism lets localities resist and experiment. Second chambers, judicial review, and fixed elections each add a point at which power can be stopped.

Every one of these has a price: gridlock, protected local abuses, overweighted provinces, political judges, short horizons. Constitutional design is not a choice between good institutions and bad ones. It is a choice of which failures to fear most.

The liberal priority is \emph{reversibility}. It must be possible to repeal a bad policy and remove a bad official. A censored opinion must be able to resurface, and a minority must be able to organize until it is a majority. The system must not depend on the virtue of whoever currently rules.

Constitutionalism also makes it safe to lose. Where defeat at the polls means prison, expropriation, or exile from public life, every election is a fight for survival, and parties will cheat sooner than lose. Stable democracy requires that losing office can be survived.

Besides voice there is \emph{exit}. A person can change town, employer, church, or newspaper. Exit is unevenly available, and competition among jurisdictions can become a race to the bottom. But liberty is safest when protections overlap, when a person can vote, speak, sue, move, publish, and appeal. No single institution then needs to be perfect. A state that monopolizes employment, media, schooling, and housing leaves no exit at all. Democracies allow independent parties, unions, churches, universities, firms, and newspapers. Each of these can turn into a small oligarchy, but together they ensure that the state is not the only source of security, status, and information.

A word on equality, since democracy is constantly asked to deliver it. Equality before the law is a condition of liberty, and a social minimum widens real freedom by taking hunger and illness out of every decision. But equal outcomes are not equal liberty. A government that equalizes results has to decide which differences are permitted, and will end by forbidding free choices because of the unequal results they might have. The useful distinction separates \emph{difference} from \emph{domination}. A rich man with more options is no political problem. A rich man who can buy offices, laws, newspapers, and immunity is one. Policy should go after the points of domination, which are monopoly, corruption, capture, and violence. Differences of comfort and taste can be left alone.

\section{The paradox of tolerance, read properly}

Popper's paradox of tolerance is nowadays quoted as if he had recommended suppressing any opinion hostile to liberalism. He recommended nothing of the kind. His worry was that an open society can be destroyed by movements which use its freedoms while refusing the reciprocity of open debate. A movement that would silence all opposition once in power can turn tolerance into the instrument of its own abolition~\cite{Popper-TOSAIE}.

The order of Popper's argument matters. The first answer to an intolerant doctrine is rational criticism and public argument. An unpopular opinion gives no license to censor. Neither does an attack on democracy itself, or an illiberal proposal. Political freedom includes the freedom to challenge the existing order. Things change only when a movement refuses reason as such. It forbids its followers to listen to argument, declines to answer argument, and turns to force and persecution. For that case, and for that case only, Popper wrote: ``We should therefore claim, in the name of tolerance, the right not to tolerate the intolerant.''

The right he claims is not a general power to punish bad opinions. It is the right of an open society to defend the conditions under which disagreement remains possible. It applies against incitement to persecution, against armed bands, against intimidation that keeps people from the polls, and against preparations to abolish peaceful opposition. To be radical, offensive, or wrong does not reach the threshold. The threshold is the substitution of coercion for argument.

This puts obligations on democracy's defenders, and they are regularly forgotten. Argument must be kept open, speech must be distinguished from violence, and the grounds for any coercion must be made public and subject to review. Otherwise ``defending democracy'' becomes a polite name for suppressing inconvenient competitors. The man who cites Popper in order to ban an opinion he dislikes has turned the paradox upside down. It is he who has given up arguing.

Hard cases exist. A movement may talk lawfully and prepare to make a change of government impossible. Such signs justify scrutiny of what it actually does, and they do not justify automatic prohibition. A democracy should answer intolerant arguments with arguments for as long as argument is possible, and defend itself against force when force is put in the place of argument. On coercion it takes a side. On opinion it holds back.

\section{What democracy must not promise}

The strongest case for democracy is comparative. Monarchy offers continuity and makes the quality of government a matter of heredity. Aristocracy cultivates service and turns birth into authority. Technocracy knows a great deal about narrow questions and cannot say who sets the goals or disciplines the experts. Autocracy is quick, and can be corrected only by the autocrat's leave. Revolutionary regimes promise to end domination and concentrate it in the party that speaks for the revolution.

Democracy does not solve these problems. It spreads them out. It makes authority temporary, visible, and open to challenge, and it removes rulers without civil war. Churchill's well-worn remark, that democracy is the worst form of government except for all the others that have been tried, is a warning against romance~\cite{churchill-democracy}. Democracy is good because it builds into the constitution the possibility that the present rulers are wrong. A perfect government would have no need of that provision, but real ones do.

So democracy should promise little. It cannot promise to satisfy every preference, as Arrow showed. It cannot promise that the majority will be right, for which history is sufficient evidence, or that elections will select able administrators, since campaigning and governing are different trades. It cannot promise the end of conflict, because conflict is what liberty looks like in practice. It cannot even promise small government, since protecting liberty takes capacity. Least of all can it promise moral innocence. A policy does not become right because a majority voted for it, and an election is not a rite of purification.

What it can offer is modest repair. Guard the constitutional core of speech, conscience, association, due process, and replaceable governments. Limit the conversion of money into power by disclosure, audit, and competition. Make pork visible by naming and costing it. Bring ordinary citizens in through assemblies and juries chosen by lot, as advisers and auditors. Improve public information without setting up an official truth. Keep defeat survivable. Write the emergency rules in calm weather. None of this abolishes voting cycles, ignorance, faction, or ambition. It keeps an error from becoming a reason to close the system against correction.

Rules are not enough either. A country can have a constitution, elections, and independent courts and still be a place where everyone treats every institution as a weapon. Liberal democracy rests on voluntary obedience. I accept a lawful decision I dislike because I keep the right to contest it. Fear does not make me accept it, and I need not think the rulers right. Obedience of this kind depends on reciprocity. People bear public burdens when they believe that the rich and the well-connected bear them too. A tax law enforced selectively teaches that law is a service one buys.

It also takes a certain civic competence. A citizen should be able to tell a policy failure from a constitutional failure, and corruption from a result he merely dislikes. He should know the difference between losing a vote and losing the possibility of ever winning one. Schools can teach how a bill becomes law, how a budget conceals its costs, how a court reasons, and how democracy has been used to dismantle itself. A citizen who knows only that democracy is good will defend it badly. And there is a place for civil disobedience, provided it is public, principled, proportionate, and aimed at an injustice that can be repaired. It marks the gap between what is legal and what is legitimate.

Democracy needs citizens who can lose without feeling humiliated and win without becoming intoxicated. The loser must not conclude that the public is illegitimate, and the winner must not conclude that his opponents are enemies. Political maturity consists in telling power from truth. A victory confers authority to act within a mandate. It does not reveal that the victor's view of the world is correct.

\section{Governing without becoming master}

The whole argument can now be put briefly. Liberty needs a protected sphere and a public order that restrains both private and official coercion. Only government can secure that order, and government is itself the most concentrated coercive power there is. The task is to make it both capable and correctable.

Democracy supplies the correctability, and it does so imperfectly. Majorities oppress, elections are manipulated, and parties are captured. Voters are ignorant, voting rules cycle, and representatives trade public money for support. Bureaucracies protect themselves. Courts overreach, and the media inflame. But in a constitutional democracy these faults are distributed among many centers of power. A citizen can resist by speech, association, vote, lawsuit, newspaper, town hall, or simply by leaving. None of these guarantees success. Together they make it hard for any one institution to close every door.

The protections of liberty usually look like inefficiency. A court delays the executive. A newspaper prints the secret, parliament refuses, the opposition obstructs, and a city declines to comply. All this is maddening because it stops power from moving in a straight line. It is valuable for the same reason, since power that moves in a straight line is very hard to stop.

The conclusion is constitutional modesty. Use elections, and do not mistake them for justice. Use expertise, and do not give it authority over purposes. Use majority rule, and protect the minority's right to become the next majority. Use strong government where only public power can prevent domination, and bind it by law, review, and expiry dates. Use the lot to interrupt professional politics, and keep responsibility clear.

Some arrangement of power is unavoidable, and the only choice is between ways of arranging it. Liberal democracy is the best of them for a people that puts individual liberty first, because it treats the ruler as a temporary officeholder and never as an owner. It does not guarantee that citizens will be wise. At its best it guarantees that they may go on challenging whoever is wrong.


% ===================== Chapter file ch09 =====================
\chapter[The Small Uses of Wickedness]{The Small Uses of Wickedness: Brecht and Two Supermarket Queues}
\label{chap:small-wickedness}

\section{When goodness has no door}

Bertolt Brecht's \emph{The Good Person of Szechwan} is usually described as a play about how hard it is to stay good in a bad world. That is true and much too mild. The play treats goodness as a practical problem of boundaries. What becomes of a person so good that she cannot shut her door, refuse a request, collect a debt, or tell an emergency from an appetite? She does not become a saint. She becomes a resource, and a resource open to everybody is soon used up.

Three gods come to Szechwan looking for one good human being. The comedy starts there, since the world is supposed to run on their moral order and they can find nobody able to live by it. Wang the water seller tries to get them a bed for the night, and one door after another closes. Everybody has a reason. At last Shen Te takes them in. She is a prostitute who can barely keep herself. The person with the least room is the one who makes room.

The gods pay her, and she buys a small tobacco shop. She hopes that a modest income will allow her to live decently and help others. For a moment the old problem looks solved, and it seems that security will make goodness affordable. In fact the shop becomes the address to which every unmet claim in the neighborhood is delivered. The former owner wants rice. An unemployed man wants cigarettes. A family that once sheltered her moves in, followed by a widening circle of relatives who eat the stock and take her inability to say no for permission. A carpenter demands payment for shelves she never ordered. The landlady wants references and rent in advance. Each demand sounds reasonable when taken singly, and together they ruin her.

Shen Te's property, time, and attention have lost their inside and outside. She feels the pain behind every appeal and cannot rank the appeals, or keep back what she needs to survive. She cannot bring herself to say that one person's need creates no unlimited claim on another. Her shop is a lifeboat. It floats, so everyone climbs aboard, and if she lets every hand hold on, it will go down with all of them.

So she invents a cousin. She puts on men's clothes, changes her voice, and comes back as Shui Ta, the hard-headed relative who is supposed to run the business. Shui Ta does what Shen Te cannot do under her own name. He throws out the uninvited family, bargains with the creditors, stands up to the landlady, and knows a gift from an entitlement. The neighbors who praise Shen Te's goodness deplore his harshness. They do not know that the harsh man is the condition on which the generous woman can go on existing.

At first Shui Ta appears only in emergencies. Then the emergency becomes permanent. Shen Te falls in love with Yang Sun, an unemployed pilot whom she saves from hanging himself and who proves quite willing to turn her devotion into cash. The wedding falls through and leaves her pregnant. Now she has a child to protect, and Shui Ta takes over for longer and longer. The shop grows into a tobacco factory. The people who once ate Shen Te's rice now work under Shui Ta's discipline. She has survived, by reproducing the hardness from which she wanted to rescue everyone. In defending the material basis of her goodness she has become an exploiter.

When Shen Te has been missing too long, Shui Ta is accused of doing away with her. The gods return as judges. With the courtroom cleared, she takes off the disguise and explains the trap. To be good to others and stay alive, she had to split herself in two. The gods, faced with the failure of their commandment, decline to amend it. They are delighted that their one good person still exists. They tell her to carry on being good, and they float away. The epilogue hands the unsolved problem to the audience~\cite{Brecht-GMvS}.

Brecht is not recommending that we become Shui Ta. Shui Ta grows from necessary firmness into domination, and that is part of the indictment. But the play also forbids the comfortable conclusion that warmth is enough. Good intentions without limits pay for the very conduct they deplore. A person who never says no will end by needing someone else to say it for her, and that someone will charge for the service.

\section{Institutions that refuse to disappear}

What Brecht inflicts on one shopkeeper happens to whole moralities and whole countries. A teaching can praise renunciation, forgiveness, hospitality, and tolerance. Any institution that means to carry the teaching through time must still decide whom to admit, what to permit, which claims to honor, and when to resist. If it refuses to decide, its openness becomes the means by which it is emptied or captured.

\subsection{Nietzsche and the slave revolt}

Nietzsche gave the most aggressive version of the argument. In the first essay of \emph{On the Genealogy of Morality} he describes a ``slave revolt in morality.'' An older, aristocratic scale had called strength, health, rank, and success \emph{good}, and had kept \emph{bad} for the common and the weak. Those who were denied any outlet in action turned the scale upside down. Their \emph{ressentiment} became creative. The powerful enemy was renamed \emph{evil}, and poverty, humility, patience, and suffering were renamed \emph{good}. Nietzsche places the decisive reversal in Jewish priestly valuation and treats Christianity as its victorious, universal form~\cite[First Essay, \S\S~7--10 and 16]{Nietzsche-GM}.

This is often flattened into the claim that Christianity was a religion designed for Rome's slaves. That takes the provocation for a census. Nietzsche's ``slave'' is a psychological type, anyone who cannot act openly against a superior and takes revenge by changing the moral meaning of superiority. Early Christianity did speak to slaves and the poor. It also crossed every social line.

Under the polemic there is a real problem. Turn the other cheek. Love your enemies. Forgive seventy times seven. Give to whoever asks. Read with prudence, these commands discipline pride and vengeance. Read without any limit, they ask the believer to give up the conditions under which such values can be borne at all. A person who ranks every demand above her own survival is Shen Te. A community that honors every attack on its way of life and counts self-defense as betrayal will preserve the commandment to yield by yielding away everyone who still believes in it. A morality of self-denial pushed to the end denies its own future. In his third essay Nietzsche goes further and says that a will turned against life would sooner will nothingness than not will at all~\cite[Third Essay, \S~28]{Nietzsche-GM}.

\subsection{The Church as its own Shui Ta}

Historical Christendom never tried unlimited meekness. The Roman Church preached charity and forgiveness, and it built property law, canon law, courts, taxes, diplomacy, and procedures for disciplining members and shutting out opponents. The Holy See could negotiate, refuse, condemn, and excommunicate. At times it went a great deal further, and used or blessed coercion that sits oddly beside the Sermon on the Mount.

One can call this hypocrisy, and often it was. It was also an institutional answer to Brecht's problem. The Church did not stake its survival on the hope that every claimant and enemy would return its charity. Its administration served as a corporate Shui Ta, a stern guardian of the material body in which the language of mercy could go on being spoken. The arrangement has kept Christianity alive for two thousand years. Survival does not prove, of course, that every means of survival was Christian.

The Reformation did not abolish the hardness. It moved it. Some Protestant traditions, the Reformed and the Puritan above all, brought discipline out of a distant papal office into the congregation, the household, and the conscience, where it could be more immediate and more severe. Others, among them the Mennonites and the Quakers, developed demanding forms of nonviolence. To reject Rome did not remove the need for a boundary. It produced several new ways of drawing one.

Nietzsche would say that these institutions survived by betraying the morality they proclaimed. A friendlier account is that an ethic made for one person's conscience cannot serve unaltered as the constitution of an organization. The tension is there on either account. Mercy needs an institution, and institutions develop interests that mercy cannot govern.

\subsection{Popper's boundary}

The political form of the same structure is Popper's paradox of tolerance, which I discussed in Chapter~\ref{chap:government-and-liberty}. Unlimited tolerance can destroy tolerance, when an intolerant movement uses an open society's freedoms to silence its opponents and abolish future disagreement. Tolerance then flows out exactly as Shen Te's rice did. It extends its protection to the organized destruction of the thing protected.

Popper's qualification is essential, and it is the part people leave out. He did not recommend suppressing a doctrine for being offensive or illiberal. As long as it can be met by argument, suppression is unwise. The right of self-defense arises when a movement refuses argument, forbids its followers to listen, and answers reasons with ``fists or pistols''~\cite[vol.~I, ch.~7, note~4]{Popper-TOSAIE}. Anyone can call his opponent intolerant and then claim the privilege of gagging him. Popper wrote no blank check for that maneuver. What counts is conduct, meaning intimidation, persecution, and violence.

So Popper corrects Nietzsche and Brecht together. A value does not survive by unlimited surrender, and it does not survive by unlimited counterattack. It needs a boundary that can tell criticism from capture and an awkward neighbor from an enemy. A principle that prohibits its own defense is an instruction for its own disappearance.

\section{The first supermarket: two items and a refusal}

Now two smaller scenes from a world without gods. Both took place in supermarkets, where strangers have to share out space, time, and attention among themselves with no script.

One day in Vienna, on my way into a supermarket, I noticed two women talking by the entrance. I took them for Roma, which was an inference and may have been wrong. One stayed at the door and asked the passers-by for money. The other went inside. Begging is a familiar sight in Vienna, in the underground as well, and it seems to come in waves. The seasons account for some of that, and so do the police, who move people from one pitch to the next. Informal arrangements about who works which corner play a part. So does the observer, who notices the busy weeks and forgets the quiet ones.

At the checkout I was holding two items. The woman who had gone in stood behind me with one. She asked whether I would let her go ahead, since she had only the one.

I said no.

She looked surprised, and then annoyed. ``Why?'' Her manner suggested that she was in a great hurry.

``Because I am a bad person,'' I said.

She shook her head. I think she was astonished that a local would refuse so small and so customary a favor. She may only have been irritated by an answer built to end the conversation. In either case it worked. The phrase offered her no principle to dispute and no excuse to knock down. I had not said that my time was worth more, or that she had picked the wrong queue, or that some urgent business prevented me. I accepted the unattractive part and kept my place.

I paid and put my two items in my rucksack. On the way out I saw her again. She was standing at the entrance, deep in conversation with the woman who was begging. Her hurry had evidently found time for a chat.

That last glimpse gives the story the shape of a fable, and I should not pretend to know more than I saw. Perhaps the talk at the door was urgent in its own way. What the scene does establish is narrower and sufficient. Her single item was a reason for asking and gave her no right to pass. I had been invited to make a small present of my time. I was free to make it and equally free not to. By calling myself bad I exposed the odd moral inflation by which declining a favor is made to sound like an offense.

\section{The second supermarket: who is next?}

The second scene took place in a village supermarket in Upper Austria. I had come to collect two sides of salmon that I had ordered at the meat and cheese counter. Behind the counter a friendly saleswoman with rather halting German was toiling to satisfy an elderly customer with a long list of special wishes. I took my place behind a heavy man who was already waiting.

More customers arrived. One of them, a stout old farmer's wife, did not join the line. She planted herself at the far end of the counter, nearer the saleswoman. From behind the counter it was now impossible to tell who had come when. There were two queues, or one queue and one rival claim.

After five minutes or so I noticed a second employee fiddling with something at the bread shelves. ``Could you please put another person on the counter?'' I called. Nothing happened. Then she began, unhurriedly, to slice cheese, apparently for a customer who was not there.

Several more minutes went by. Both clusters grew. The one saleswoman who was serving called into the back room, ``Come out, please, we have customers.'' Nobody came. Only when she went in herself did a plump young colleague emerge into the daylight. She had dug herself in back there, doing heaven knows what. She tied on her apron with a sulky air and looked around.

``Who's next?''

I pointed to the heavy man in front of me. ``He is.''

He was served. Meanwhile the friendly saleswoman had finished with the old lady. She turned to the farmer's wife, who had arrived after me and was standing closer.

``I was here first,'' I said. ``The two sides of salmon I ordered, please.''

The farmer's wife gave me a look of disgust and began to order anyway. When I pointed out that she had come after me, she did not appear to hear.

And then the heavy man, whose turn I had just defended, stabbed me in the back. He spoke in broad local dialect, which surprised me, because I had lazily taken him for a southerner. ``Now, now, let the lady go first. We take things easy round here. It doesn't matter who's first and who's later.''

The farmer's wife agreed warmly. Taking things easy suited her very well, so long as she was served first. I insisted, and got two heads shaking in unison. How could anyone be so impatient?

The whole quarrel was about three minutes of my time, and it taught me a good deal for that reason. A queue is a tiny institution. It turns the accident of who arrived when into a rule that strangers can apply without knowing one another. Nobody has to prove that his need is greater, his rank higher, his roots in the village deeper, or his errand more important. First come, first served is a crude rule. It also keeps every purchase from turning into a contest of nerve and elbows.

Once the line has split, politeness becomes a cover for the question of who pays for everyone's being relaxed. The heavy man could afford to be generous with my place. The farmer's wife could praise the easygoing local ways because at that moment they moved her to the front. If I had given in, harmony would have been restored and the entire cost charged to me. My insistence on the order of arrival made the conflict visible, and so the insistence was treated as its cause.

The honest version has another side. The staff were busy, and the old lady was entitled to proper service. Work in the back room does not vanish because the customers are glaring. And notice what my irritation did. It took an inventory of bodies and accents: the fat man, the stout farmer's wife, the saleswoman with the foreign accent, the lazy girl in the storeroom. Those details are in my memory because irritation is observant. They do not prove what I felt they proved at the time. I report them because a truthful anecdote has to keep the narrator's uncharitable eye. The reader need not mistake that eye for a camera.

\section{A modest license to be bad}

Neither scene comes anywhere near Shen Te's ruin. Nobody lost a livelihood, no child was in danger, and I did not have to invent a male cousin to get through the checkout. The scale is comic. But the structure is hers. A request is described as trifling, because the person asked will be the one to bear the trifling cost. A rule is declared unimportant by those who profit from suspending it. And the person who resists looks meaner than the person who imposed.

There are good reasons to let somebody with one item go first. There are good reasons to give money to a stranger and to be patient with an overworked shop assistant. Courtesy makes life among strangers bearable. But courtesy stops being courtesy once it can be demanded and enforced by reproach. A gift is voluntary only if the giver was able to refuse.

Here wickedness has its small and limited use. ``I am a bad person'' is a comic shield against the demand that every boundary be justified to the satisfaction of whoever wants to cross it. It shuts the little courtroom in which the petitioner appoints himself judge and the person asked has to prove his innocence. Sometimes the only way to keep your place is to stop pleading for permission to keep it. The sentimental conscience finds this unbearable. It wants a reason for every refusal that it can approve, and so it never refuses anything.

The danger is that the joke turns into an alibi. A refusal protects a boundary. Contempt enjoys the refusal and enlarges it into a theory about whole classes of people. Shen Te needs Shui Ta's ability to say no, but Shui Ta does not remain a doorkeeper. He starts arranging the world in his own hard image. The art, if there is one, would be to borrow his firmness and not catch his appetite for mastery.

No ringing commandment follows from this. Help when you have freely chosen to help. Keep the queue when the queue is the only protection against the boldest elbow. Say no in good time, before stored-up resentment starts talking in stereotypes. And when ``being bad'' feels especially delicious, consider that the pleasure may have taken over from the boundary.


% ===================== Chapter file ch10 =====================
\chapter[The Weaponization of the Old Testament]{The Weaponization of the Old Testament: Original Sin Put to Modern Use}
\label{chap:weaponized-original-sin}

\section{A doctrine whose name is absent}

The phrase \emph{original sin} does not occur in the Bible. Nor does the Hebrew Bible contain a doctrine that every person is born guilty of Adam's offense. What Christians later taught was put together over centuries and across three languages. Genesis supplied a story of primal disobedience. Paul gave the story a universal meaning. Later theologians, Augustine above all, made a system of it. To call the result an ``Old Testament doctrine,'' as the title of this chapter mischievously does, is to mistake the root for the tree \cite{NRSVue2021,TimpeChristianSin2021,BeatriceTransmissionSin2013}.

Genesis~3 is short. A man and a woman live in a garden where everything is allowed except the fruit of one tree. The serpent disputes the threatened penalty, and they eat. What follows first is shame, and the text offers no theory of hereditary guilt. They cover themselves, hide, and blame each other. Then come hard soil, painful labor, dangerous childbirth, death, and expulsion. The story explains estrangement, from God, between the sexes, and from the ground. The word ``sin'' does not appear in the chapter. It turns up in the next one, where sin is ``lurking at the door'' for Cain, and Cain is told that he must master it \cite[Gen.~3; 4:7]{NRSVue2021}.

Other passages pull in other directions. Psalm~51:5 (51:7 in German and Hebrew numbering) laments, ``Indeed, I was born guilty, a sinner when my mother conceived me.'' Christians have long read that as proof of inborn sinfulness, though a penitential poem is hardly a theory of biological transmission. Exodus~20:5 speaks of consequences that reach later generations. But Deuteronomy~24:16 forbids punishing children for their parents, and Ezekiel~18:20 insists that the child does not bear the parent's guilt. The texts contain both inherited consequence and individual responsibility, and they keep the two apart \cite{NRSVue2021}. Our contemporaries have lost the distinction.

\section{Paul builds the bridge}

The decisive passage is Romans~5:12--19. Paul sets Adam and Christ in parallel. Through one man came sin, condemnation, and death, and through one man come obedience, justification, and life. First Corinthians has the same pattern: ``all die in Adam,'' and Christ is the ``last Adam'' \cite[Rom.~5:12--19; 1~Cor.~15:21--22, 45]{NRSVue2021}.

Adam now stands for the condition of the whole race. But what exactly is inherited? It might be mortality, or a damaged nature, or an inclination to sin, or condemnation, or the guilt of another man's act. And how can a responsibility that precedes any choice be just? A detail of translation shaped the Western answer. The Greek of Romans~5:12 is usually rendered ``because all have sinned.'' The Latin text that Augustine used could be read ``in whom all sinned,'' as though the entire human race had somehow acted in Adam. The translation did not create the doctrine, which has older roots, but it helped fix its form \cite{BeatriceTransmissionSin2013}.

The finished doctrine is an ambitious piece of thinking. It tries to explain why wrongdoing is universal without making God the author of evil, and why death and disordered desire are in place before anyone's first deliberate act. And in Christian terms Adam is only half of it. The Fall is paired with grace. The debt is stated in order that it may be paid. Whoever keeps the guilt and drops the redemption has already turned the doctrine into a weapon.

\section{From Augustine to an inherited condition}

The classic Latin formulation came out of Augustine's quarrel with Pelagius in the early fifth century. Augustine did not invent every part of it. But he put the inherited condition at the center of his account of grace, baptism, sexuality, and the crippled will, and the councils of Carthage in 418 and Orange in 529 made his view the Western inheritance \cite{BeatriceTransmissionSin2013,TimpeChristianSin2021}.

Even among Christians the doctrine has never meant one thing. The Eastern churches stress inherited mortality and corruption more than inherited blame. The present Catholic catechism calls original sin a state ``contracted and not committed,'' and says expressly that it is not a personal fault in Adam's descendants \cite[paras.~404--406]{CatechismOriginalSin}. The Protestant confessions often put the corruption of fallen man more starkly. So the slogan ``you are guilty because you were born'' is a polemical sharpening that none of the churches teaches in that form.

Judaism went another way. It does not consist of ``the Old Testament'' alone. It has the Tanakh, whose canon and order differ from the Christian arrangement, and beside it a vast body of rabbinic law and argument. The rabbis never adopted Augustine's system. They grant that Adam's act changed the human situation and that man contends with an evil inclination. They also insist that each person is able, and obliged, to choose, to repent, and to return \cite{HaymanRabbinicFall1984}. Man is not innocent in this view. He is answerable for what he does, which is another thing entirely.

\section{The secular heir: born with a footprint}

The old structure has had a remarkable second career. A child is born and at once begins consuming. He breathes out carbon dioxide and drinks water. He must be kept warm, fed from farmland, clothed, housed, and driven about. Before he has formed a single intention he has a carbon footprint, a water footprint, and an ecological footprint. The accountancy is then turned into a morality tale. To arrive is to incur a debt. Ordinary metabolism is evidence for the prosecution. Cutting down is penance, and a certified payment buys remission.

The comparison with original sin is exact in one respect and dangerously inexact in another.

It is exact wherever a measurable inherited \emph{condition} is quietly converted into personal \emph{blame}. Everyone is born into systems of energy, transport, food, and housing that he did not choose. In rich countries those systems carry heavy environmental costs. But to take part in a cause is not yet to be guilty. Responsibility depends on knowledge, freedom, available alternatives, power, and the difference one's own decision makes. Wipe out those distinctions and an emissions inventory becomes a doctrine of stain.

It is inexact if it is used to deny the physical problem. Human activity has unequivocally warmed the atmosphere, the oceans, and the land, and fossil carbon accumulates whether or not anyone feels guilty about it \cite{IPCC2021HumanInfluence,IPCC2023Synthesis}. A physical constraint does not go away because its publicists have borrowed the manner of the pulpit. I have nothing against measurement, conservation, or prices. My objection is to the shortcut by which a description of resource use becomes a verdict on the worth of the user.

Three examples show what the shortcut costs in simple accuracy.

\begin{enumerate}
\item \emph{Breathing.} The carbon dioxide we exhale belongs to the short biological cycle. Plants took the carbon in our food out of the air a few months or years ago. Breathing adds no ancient geological carbon to the atmosphere, as burning coal, oil, and gas does. Food production has emissions of its own, certainly. But ``you pollute by breathing'' is simply bad carbon accounting \cite{MITBreathingCO2}.

\item \emph{Drinking.} Fresh water is scarce, dirty, or badly distributed in many places. Thirst is no misdemeanor. Agriculture takes about 72 percent of the fresh water withdrawn worldwide, while billions of people have no safe drinking water. Prices, pipes, choice of crops, leakage, pollution, and allocation matter far more than any general indictment of people for drinking \cite{UNESCOWWDR2024}.

\item \emph{Energy.} A kilowatt-hour from lignite, from uranium, from falling water, from wind, and from sunlight are five different things. They differ in emissions, waste, use of land, and reliability. To treat energy as sinful in itself hides the engineering choices by which people can live well while particular harms are reduced.
\end{enumerate}

At bottom the error is one of grammar. ``This activity has a cost'' turns into ``this person is a cost.'' Once that step is taken, the human being is no longer the subject for whose sake institutions exist. He is an unwanted load on a planetary balance sheet. Environmental responsibility has given way to an environmental anthropology, which holds that a person is damage, where the older view held only that people sometimes do damage.

\section{Allowances, offsets, and indulgences}

Carbon payments are often compared with indulgences, and the comparison is effective because both keep a ledger of offense and compensation. It is also easily abused, so the history should be got right.

In the early sixteenth century, indulgence campaigns were tied to the rebuilding of St~Peter's in Rome. A Mainz indulgence of 1515, issued at the instance of Johann Tetzel, survives as a document of that fund-raising drive \cite{EmoryIndulgence1515}. Technically an indulgence gave no permission to sin and did not buy forgiveness. In Catholic doctrine it remitted the temporal punishment remaining after guilt had been forgiven. The scandal lay in the sales machinery, and in the ease with which a subtle doctrine was heard as a transaction in which money came in and liability went out.

Carbon instruments need the same care. An \emph{allowance} in a cap-and-trade scheme is a permit under a legal ceiling that is lowered year by year. The European Union's Emissions Trading System obliges emitters to surrender allowances and lets them trade. The cap, the monitoring, and the enforcement do the work, and nobody's soul is involved \cite{EUETS}. A voluntary \emph{offset} is another matter. The buyer pays for a reduction somewhere else and sets it against his own continuing emissions. Offsets can finance real projects. Their honesty depends on whether the reduction would have happened anyway, whether it can be measured, whether it lasts, and whether it is counted twice. Current guidance accordingly puts direct reductions first and reserves durable removals for whatever cannot be eliminated \cite{OxfordOffsettingPrinciples2024}.

The indulgence comparison is fair only where the instrument is given a moral function that its physical design cannot carry, where the buyer reasons that he has paid, need change nothing, and has a clean conscience. It is unfair when it pretends that every price on emissions is a fraudulent absolution. A well-made cap reduces total pollution even if every participant remains exactly as sinful as before. Public policy should prefer that humdrum result to a theater of purified souls.

There is a political danger all the same, in any regime of permanent confession. Whoever designs the categories decides which necessities count as excess and whose consumption must be cut. He also decides who may buy release. The rich turn money into certified virtue, and the poor meet the same ledger as higher prices and fewer choices. The old hierarchy is back in secular dress, with guilt laid on everyone and dispensations reserved for those who can pay.

\section{The neglected commission}

Against the picture of man as an intruder stands a sentence from the first chapter of Genesis that one seldom hears quoted. God blesses humanity and commands it to ``fill the earth and subdue it; and have dominion'' over the other creatures \cite[Gen.~1:28]{NRSVue2021}. The Hebrew verbs are strong ones. The verse cannot honestly be rewritten as an instruction to admire untouched wilderness from a respectful distance. It authorizes reproduction, settlement, cultivation, and rule, and the transformation of a world that was given as good and was not delivered as a finished home.

The commission is no blank check. Genesis~2:15 puts man in the garden ``to till it and keep it.'' Biblical law goes on to limit ownership, protect animals and the poor, command rest, and remind the owner that the land is not absolutely his. Pope Francis reads the two verses together in \emph{Laudato si'}. Dominion cannot justify unrestrained exploitation, he says, because it has to be read with tilling and keeping \cite[para.~67]{FrancisLaudatoSi2015}.

Each text corrects a temptation. Dominion without keeping is plunder. Keeping without dominion turns into hostility toward every human mark on the landscape. The first forgets that power must answer for itself. The second forgets that human intelligence, labor, and building belong to the creation the text describes. A city, a field, a reservoir, a power station, or a hospital is not justified merely because people built it. It is not condemned merely because it changed something either.

The pairing also shows what is lost when the Fall is cut loose from the creation. In Genesis, humanity is blessed before it transgresses. Fruitfulness and work come before the curse. The ground resists, but tilling it was never the crime. A rhetoric that remembers only the inherited stain and forgets the earlier blessing has made a tragic story about misused freedom into an accusation against life.

\section{What may legitimately be inherited}

We do inherit consequences. A newborn receives an atmosphere altered by earlier fires, cities built around earlier decisions, the public debt, institutions, knowledge, machines, landscapes, and inequalities. It would be as silly to deny this as to deny that Adam's children live outside Eden. But an inherited \emph{situation} and inherited \emph{guilt} are different things. The first can create duties for the future, and for that purpose nobody need invent culpability for the past.

A defensible environmental ethic asks four prosaic questions. What physical harm is being done? Who can change it at a bearable cost? Which measure does more good than harm? And how should the burden be shared among people whose power, history, and alternatives differ enormously? These questions leave room for taxes, standards, invention, conservation, and even sacrifice where the evidence supports it. They require no child to apologize for breathing, and no citizen to regard his ordinary existence as a moral defect.

Original sin has been turned into a weapon wherever the verdict is known before any act has been examined, and humanity is held guilty for being present. The usefulness of this is obvious. A guilt that attaches to existence can never be paid off. It gives administrators, campaigners, companies, and markets a permanent claim on everybody's conduct. That usefulness is also the reason to resist. A morality fit for free people has to distinguish being from doing, necessity from waste, and consequence from guilt. It has to know real repair from innocence bought at a counter. Without those distinctions the old drama of fall and redemption shrinks to a one-sided technique of compliance. The Fall is kept and grace is left out, the ledger never closes, and life itself is entered on the debit side.


% ===================== Chapter file ch11 =====================
\chapter[Two Plus Two Is Five]{Two Plus Two Is Five: The Politics of Administrative Truth}
\label{chap:two-plus-two-five}

\section{The minimum freedom}

The most ambitious kind of political conditioning does not aim at making people approve of a policy. It aims at making them accept that the difference between true and false is itself a matter for the authorities. George Orwell stated the minimum resistance in \emph{Nineteen Eighty-Four}: ``Freedom is the freedom to say that two plus two make four.'' If that is granted, he adds, all else follows \cite[Part~I, chap.~7]{Orwell-1984}.

The sentence has little to do with arithmetic. Two plus two is Orwell's deliberately boring proposition. It needs no party program, no expert commission, and no moral enthusiasm. Its value is that it does not depend on what anyone wants. In the torture chamber O'Brien is not satisfied when Winston \emph{says} that four fingers are five. A lie told under threat leaves a witness alive inside. O'Brien wants Winston to \emph{see} five when the Party requires five. The victory is complete when the victim has surrendered the faculty by which he could have noticed the contradiction, and the coercion no longer feels like coercion \cite[Part~III, chap.~2]{Orwell-1984}.

Modern administrations seldom ask for anything so crude. Nobody announces that four is five. The unit is changed, or the boundary of the system, or the denominator, or the time horizon, and the result is then presented as though nothing had been altered. A measurement at the exhaust pipe becomes ``zero emissions.'' An uncertain scenario becomes ``the forecast.'' A political wish becomes ``the science,'' and a gain in the aggregate becomes a benefit to everybody. Each of these may rest on a defensible calculation. The deception begins at the moment when the conditions under which the calculation holds drop out of the sentence.

So this chapter is no catalogue of conspiracies. Most public misdescription is more ordinary. A useful metric is promoted to the whole picture. A slogan outlives its qualifications. Institutions copy one another's vocabulary. And critics who notice the omission answer with an equal and opposite exaggeration. The task is to defend arithmetic without pretending that every disputed question of policy is arithmetic.

\section{Four ways of making five}

What people loosely call ``lies'' falls into at least four classes.

\begin{enumerate}
\item A \emph{plain factual contradiction} asserts that an observable event did not happen when the evidence shows that it did. This is the nearest thing to two plus two is five.

\item A \emph{boundary substitution} reports a correct figure for part of a system and gives it the name of the whole. ``Zero CO$_2$ at the tailpipe'' becomes ``zero-emission transport.'' The second does not follow from the first.

\item A \emph{modal substitution} turns a target, a scenario, or a conditional projection into a fact about the future. ``This pathway could limit warming'' becomes ``the climate will be held at this temperature.''

\item A \emph{normative substitution} presents a choice among competing goods as though science had selected the goal as well as the means. Evidence can estimate consequences. Without an added judgment of value it cannot say how costs, risks, and liberties ought to be distributed. The step from \emph{is} to \emph{ought} is still a leap \cite[Book~III, Part~I, sec.~1]{Hume-Treatise}.
\end{enumerate}

The distinctions matter because the remedies differ. A contradiction should be withdrawn. A bounded measurement should have its boundary put back. A forecast needs its probabilities and assumptions. A political judgment needs an honest argument about values. To shout ``misinformation'' at all four is as useless as to call all four harmless.

\section{A fuel consumption that no car has to achieve}

The European Union's carbon-dioxide standards for new cars are an unusually clean specimen. The official fleet-wide values can be converted into liters of fuel, at roughly 2.31 kilograms of CO$_2$ per liter of petrol and 2.68 per liter of diesel:

\begin{center}
\begin{tabular}{lccc}
\hline
Period & Fleet target & Petrol equivalent & Diesel equivalent \\
       & (g CO$_2$/km) & (l/100 km) & (l/100 km) \\
\hline
2025--2029 & 93.6 & 4.1 & 3.5 \\
2030--2034 & 49.5 & 2.1 & 1.8 \\
From 2035 & 0 & 0.0 & 0.0 \\
\hline
\end{tabular}
\end{center}

The arithmetic is correct, and the table invites a misreading. No individual car has to meet these figures. They are reference values for the \emph{average} certified tailpipe CO$_2$ of each manufacturer's newly registered fleet. The value of 93.6 applies from 2025 to 2029. An amendment of 2025 allows compliance for 2025, 2026, and 2027 to be judged on the three-year average, and sets no different target for those years \cite{EUCarCO2Standards,EUCarFlexibility2025}.

A battery-electric car counts as zero at the tailpipe because it has no tailpipe. The generation of its electricity, the manufacture of its battery, the raw materials, the charging losses, and the scrapping all fall outside this particular metric. A manufacturer can therefore bring down its regulated average by selling electric cars, although each of them consumes energy and embodies a good deal of mined material. The electricity does come ``out of the socket.'' But a socket is a point of delivery and generates nothing.

It does not follow that electric cars merely move the same emissions somewhere else. The life-cycle comparison is an empirical matter, and it depends on the electricity mix, the size of the battery, the mileage, and the combustion car chosen for comparison. The European Environment Agency concludes that a typical electric car in Europe causes less greenhouse gas over its lifetime than a comparable petrol or diesel car, while its manufacture causes more \cite{EEAElectricVehiclesLifecycle}. ``Zero-emission car'' is thus a misleading name. ``Electric cars do nothing for the climate'' is, for Europe, generally false as well.

The correct sentence is long: \emph{for the purposes of one regulation, battery-electric vehicles are assigned zero certified tailpipe CO$_2$, and under present European conditions their life-cycle emissions are usually lower than those of combustion cars, but not zero}. Public language first shortens this to ``zero-emission vehicle'' and then to ``zero emissions.'' Nobody falsified a measurement. Two plus two was made into five by peeling off the label.

The line for 2035 needs a date stamp of its own. After the 2025 amendment the law still prescribed zero grams per kilometer for new cars and vans from 2035. By 2026 the Commission had proposed replacing that with a 90 percent reduction and making up the rest by specified measures \cite{EUCarCO2Standards}. The ``end of the combustion engine'' is therefore the current state of a piece of legislation, and it may yet change. It has never been a physical fact.

\section{Migration: ``manageable'' is not a datum}

The statement that a large wave of migration is ``manageable'' has the grammar of a finding and no unit of measurement. Manageable for which town, which budget, which school, which housing market, which hospital, and which police district, and over what period? And by what standard? The standard might be that the state has not collapsed, or that the budget balances, or that rents are stable, or that the newcomers have jobs and speak the language, or that crime is low, or that the public consents. Until those terms are supplied, ``manageable'' is a declaration of resolve. It may be brave or it may be reckless, but nobody can yet test it.

Migration has no single effect either. Outcomes depend on age, education, legal status, reason for coming, family, language, selection policy, the speed and scale of arrival, and the spare capacity of the receiving country. A skilled worker hired into a vacancy, a foreign student, a citizen's spouse, an asylum seeker forbidden to work, and a refugee family set off five different fiscal and institutional histories.

Aggregates conceal who gets what. Total GDP rises when the population rises, even if GDP per head stands still. Employers and consumers may gain while workers in certain trades and tenants in tight housing markets lose. A reassuring national average is compatible with acute pressure on a few cities and their schools. ``Migration benefits the economy'' is therefore incomplete for the same reason as ``migration ruins the economy.'' Neither says which migrants are meant, which institutions, over what time, compared with what, or to whose benefit.

The evidence is conditional as well. The OECD's \emph{International Migration Outlook 2024} reports that record inflows strained reception systems, housing, and public services, while countries went on looking for foreign workers to fill real shortages \cite{OECDMigrationOutlook2024}. A simulation by the EU's Joint Research Centre finds that larger inflows yield only small fiscal gains unless the barriers to employment are removed, and much larger gains if they are \cite{JRCFiscalMigration2020}. Neither result supports ``nothing important will change.'' Neither supports ``integration can never work.'' Capacity is no incantation. It is a ratio between a flow of people and the institutions that are supposed to absorb it.

Political speech avoids that ratio. When difficulties appear, ``manageable'' is redefined after the event to mean that the state is still standing. When benefits appear, they are all credited to migration, and nobody asks what a different admissions policy would have produced. A claim that no observation could ever defeat is not a claim about facts. It is a badge of allegiance.

\section{The climate is not a thermostat}

Another popular abbreviation holds that policy can ``keep the climate'' at a chosen point. Taken literally this is false. Climate is not one number, and the global mean temperature is no thermostat setting. Regional temperatures, rainfall, storms, ocean currents, and ice go on varying. Volcanoes, the sun, and the internal sloshing of ocean and atmosphere do not consult the carbon budget.

The opposite slogan, that emissions policy cannot affect future temperature, is false too. The IPCC finds a nearly linear relation between cumulative man-made CO$_2$ emissions and the warming they cause. To stabilize that human contribution at any level, global net emissions of CO$_2$ must fall to about zero. The size of the remaining budget is uncertain, and other gases and natural variability still matter \cite{IPCCWGI2021TechnicalSummary}.

A target of 1.5 or 2 degrees is therefore no promise that the climate will stand still. It is no piece of numerology either. It is a boundary of risk drawn on a modeled global average, with stated probabilities, baselines, and uncertainties. ``Limit warming'' is accurate, and ``set the climate'' is not. Net zero would stabilize the human-induced part of the warming. It would not stabilize the weather.

The example shows a polemical trick that both sides rely on. Advocates advertise more control than the models justify. Critics then demolish a silly version in which ministers operate a planetary thermostat. The public is left to choose between omnipotence and impotence, and the evidence supports neither.

\section{Other respectable sums that come to five}

\subsection{``Net zero'' means zero}

It does not, and the work is done by the word \emph{net}. In IPCC usage, net-zero CO$_2$ is reached when man-made emissions are balanced by man-made removals over a stated period. Net zero for all greenhouse gases requires a further convention for comparing different gases \cite{IPCCWGI2021TechnicalSummary}. A net-zero plan may thus keep residual emissions and offset them with removals.

That is legitimate bookkeeping and no fraud in itself. The misrepresentation begins when a headline lets the audience hear \emph{gross} zero. It becomes serious when the promised removals are short-lived or counted twice, when they would have occurred anyway, or when they exist only in a model of the year 2050. Conversely, to point out the residual emissions is not to show that the plan does nothing. Everything depends on the boundary and on the quality of the balancing item.

\subsection{``Renewables are the cheapest electricity''}

On one well-defined comparison this is true. IRENA reports that most new utility-scale renewable plants commissioned in 2024 generate more cheaply, by levelized cost of electricity, than the cheapest new fossil alternative \cite{IRENARenewableCosts2024}. Levelized cost is a sound measure of what it costs one plant to produce a kilowatt-hour over its lifetime.

But a generating plant is not an electricity system. Customers need power at the moment they switch on, which need not be the moment the wind blows. The cost of a system includes transmission, balancing, storage, reserve capacity, curtailment, and whatever else reliability requires. Every technology imposes system costs, and they differ in kind and in size. The OECD's Nuclear Energy Agency accordingly warns against taking the cost at the plant for the cost of a reliable low-carbon system \cite{OECDNEASystemCosts2019}.

The honest statement is that new solar and wind plants can generate very cheaply, while the cheapest \emph{dependable} system depends on location, the share of variable sources, the grid, storage, and the available backup. Solar and wind are not secretly ruinous, and renewables are not always cheapest. To delete ``by levelized cost'' is the same operation as deleting ``at the tailpipe.''

\subsection{``One hundred percent recyclable''}

The label generally refers to technical design, or to what selected components could in theory undergo. It does not say that the item will be collected, sorted, processed without contamination, turned into material of equal quality, and used again. Collection losses, composite materials, additives, degradation, energy, cost, and the demand for the recycled product all lie between what is possible and what happens.

The gap is visible on the scale of a continent. According to the European Environment Agency, recycled material met only 11.8 percent of the EU's material use in 2023 \cite{EEACircularMaterialUse2024}. Recycling is not thereby shown to be useless. The figure shows why a circular economy also needs repair, reuse, longer-lived products, and less material to begin with. ``Could in principle enter a recycling process'' does not mean ``will come back for ever without loss.''

\subsection{``The science says we must''}

Science can say that a measure will probably reduce a specified risk at a specified cost. It can compare scenarios and expose impossible promises. Without political premises it cannot decide whose risk counts most, which rights may be restricted, how uncertainty should be valued, or how the present should be weighed against the future. To say so is not to attack science. It marks the place where empirical premises need moral and constitutional ones beside them.

As casual shorthand ``the science says'' is harmless. It becomes an abuse when a dispute about means, proportion, distribution, or rights is restated as a dispute about a laboratory result. It also tempts the opponents to attack the science when their real objection is to the value judgment. Administrative truth flourishes when both sides misname what they are arguing about.

\section{Why the short version survives}

It is tempting to blame ``state-funded media,'' and the explanation is too coarse. Public broadcasters, subsidized newspapers, commercial networks, and independent publishers work under different rules. Funding does not prove obedience, and repetition does not prove a central command. Convergence needs no conspiracy.

Governments and international bodies have permanent press offices, accredited experts, scheduled announcements, standard statistics, and the authority to define the regulatory categories. Their language arrives ready to print. Journalists work to deadlines and need sources that can be quoted and held responsible. Media researchers call the result ``indexing.'' Mainstream news tends to reproduce the range of disagreement among recognized officials, and seldom goes beyond it. The thesis was developed for the United States, and the incentive behind it travels well \cite{BennettPressState1990}.

Short claims also win the competition. ``Zero-emission car'' fits a headline. A paragraph on tailpipe boundaries, electricity mix, vehicle lifetime, and fleet averaging does not. ``We can manage'' conveys confidence, and a table of fiscal, housing, and schooling conditions conveys doubt. ``Renewables are cheapest'' can be remembered. The difference between plant cost and system cost cannot.

The distortion need not start in the newsroom. Ministries, campaigners, companies, researchers, and critics all learn that a clean statement travels further than a qualified one, and the media amplify statements that have already been selected for repeatability. Public money may bring an outlet close to power. Advertising, ownership, audience tribes, and dependence on access exert pressures of their own. The trouble is an ecology of incentives. No single profession is the villain.

\section{An audit for public arithmetic}

Before accepting or rejecting a polished public claim, seven questions are more useful than asking which camp the speaker belongs to.

\begin{enumerate}
\item \emph{Where is the system boundary?} Tailpipe, power station, national territory, supply chain, or whole life cycle?

\item \emph{What is the denominator?} Total GDP or GDP per head? Total emissions or emissions per unit? The national average or the worst-affected district?

\item \emph{Is the number physical or administrative?} A measured quantity, a certified test value, a model output, an accounting credit, and a legal target are five different things.

\item \emph{Compared with what?} The existing fleet, a new combustion car, the bus, or no journey at all? A different admissions policy? A different energy system?

\item \emph{Over what time?} Construction may be expensive and operation cheap. Migration may cost now and pay later. An offset may promise storage for less time than the carbon stays in the air.

\item \emph{Where are the uncertainty and the distribution?} An expected gain in the aggregate is compatible with heavy losses for particular groups, and a model average can hide a wide range.

\item \emph{What observation would prove the claim wrong?} If every outcome can be redescribed as success, the statement is propaganda or a badge, and it asserts nothing testable.
\end{enumerate}

The questions apply to official optimism and to oppositional gloom alike. The critic who ``corrects'' zero tailpipe emissions by declaring that electric cars have no advantage has not restored four. He has picked a different five. The government that points to national GDP and ignores GDP per head and local capacity makes the same mistake as the opponent for whom every migrant is an identical fiscal liability.

\section{The right to keep the remainder}

Two plus two is four because the operation, the units, and the objects are held constant. Political arithmetic turns dangerous when one of them changes without notice. The best defense is therefore to insist on labels, which serves better than reflexive distrust. Say tailpipe and mean tailpipe, never the total. Keep the fleet apart from the single vehicle, net apart from gross, and a scenario apart from a prophecy. An aggregate is not a benefit to all, and an aspiration is not an observation.

Orwell's freedom is modest and demanding. It protects an outside world that does not change when a ministry, a broadcaster, a corporation, a movement, or a majority changes its vocabulary. It also disciplines the dissenter. The right to say that four is four gives no license to call every inconvenient estimate five. It carries a duty to keep definitions fixed, admit uncertainty, and correct one's own side. The pleasure of exaggeration has to be given up as well.

Even with that discipline, not everything follows. Policy still needs judgment, compromise, and the acceptance of real costs. But without it nothing follows by reason at all. There is only the approved answer, and the number of fingers one has been told to see.


% ===================== Chapter file ch12 =====================
\chapter[Puppy Eyes and Wolf Teeth]{Puppy Eyes and Wolf Teeth: Sentimentality, Reason, and the Unequal Price of Compassion}
\label{chap:sentimentality-rationality}

\section{The argument hidden in a face}

A bear cub does not have to argue. The round head, the short muzzle, the big eyes, and the clumsy walk do the arguing. A puppy needs no theory of animal ethics, since it only has to look up. Human attention is not handed out according to a neutral census of interests. Faces capture it, and so do voices, helplessness, and resemblance to ourselves.

There is nothing contemptible in this. A species whose young need years of care would not last long without a fast and powerful pull toward the vulnerable. In experiments that varied the infantile features of human faces, a stronger ``baby schema'' made the faces seem cuter and increased the reported wish to look after them \cite{GlockerBabySchema2009}. The study was about human babies and says nothing directly about wolf pups or kittens. Still, it is no mystery that young animals enlist the same response. The response comes first, and any calculation comes afterwards.

So the useful distinction does not set feeling against reason, as though one were a stain and the other the bleach. Feeling supplies most of the ends that reason works for. Hume put it as provocatively as he could: reason ``is, and ought only to be the slave of the passions'' \cite[Book~II, Part~III, sec.~3]{Hume-Treatise}. Facts alone cannot tell us whether the survival of wolves, the livelihood of a shepherd, or the life of a sheep matters. Something must already care.

\emph{Sentimentality}, as I use the word, is narrower than sentiment. It is an emotion protected from the information that would complicate it. It frames the appealing creature tightly and leaves the costs outside the frame. It also takes the strength of a feeling as evidence for a policy. Rationality does not abolish compassion. It asks compassion to count, to compare, to follow the chain of causes, and to notice who is paying for the beautiful conclusion.

\section{Two dogs, four cats, and a missing denominator}

``Two dogs and four cats'' is a caricature and no statistical household. Like any good caricature it enlarges a real feature. Pets have moved into the emotional category of family members. An industry survey published in 2026 counts 140 million European households with at least one pet, which is 49 percent of households in the countries covered. Between them they keep 306 million animals and spend 29.4 billion euros a year on pet food \cite{FEDIAFStatistics2026}. The figures come from the pet-food industry and should be read as a description of a market. They are no census of motives.

They do not show that pets are replacing children. People keep animals for company, affection, routine, security, and sport, and nobody can read demography off a shopping trolley. But the imagined menagerie does capture a moral style. Concentrated tenderness is directed at a few named animals, a large commercial apparatus supports it, and the owner experiences it as proof that he loves nature.

The proof fails. A man can spend freely on medicine for an old dog and never give a thought to how pigs are kept. A woman can adore cats and regard insects, toads, fungi, carrion, and mud as regrettable flaws in an otherwise pretty ecosystem. These people need not be hypocrites. Emotional closeness is simply selective. The named animal in the kitchen and the nameless animal in the flock live in different psychological worlds, though both can suffer.

Affection for a pet can widen moral concern, and it can also domesticate one's whole picture of nature, so that every good animal is an innocent personality waiting to be appreciated. But wild nature is not a big pet shop. Predation, starvation, parasites, and the killing of young are no policy failures caused by a shortage of kindness. The trouble starts when the intimacy that suits a companion animal becomes the model for governing an ecosystem.

\section{Charisma is not biodiversity}

Biodiversity is an abstraction, and its public relations are poor. It consists of genetic variation, populations, species, habitats, and processes. Most of it is small, nocturnal, slimy, prickly, poisonous, or invisible, and it does not look into a camera. Conservation therefore appoints ambassadors: the panda, the tiger, the elephant, the bear, the wolf. A charismatic species turns a technical problem of habitat connectivity into a story that attracts attention and money.

The bias has been measured. In a study of more than ten thousand participants in a zoo's adoption program, people chose charismatic species far more often, and being endangered did not by itself make a species more likely to be picked \cite{ColleonyCharisma2017}. This does not make the donors fools, or the money wasted. It shows that public preference and conservation priority are two different variables.

The difference cuts two ways. ``Biodiversity'' can be a respectable translation of affection, in which one wants the magnificent animal first and supplies the ecological vocabulary afterwards. But charisma can also be put to work. To protect a wide-ranging predator, large connected habitats may have to be protected, and those shelter a great many unphotogenic species. The poster animal pays for the marsh. It would be as careless to dismiss every large-carnivore program as sentimental as to pronounce every one of them an ecological necessity.

The test is whether the claimed benefit to biodiversity survives when the charismatic picture is taken away. Which populations, habitats, or functions improve? Compared with what other use of the same money? At what density of predators? With what effects on grazing, forestry, and other species? A word that answers all of these before they have been asked is no scientific finding. It is a benediction.

\section{Return is not reintroduction}

Accuracy matters down to the verbs. People speak of ``reintroducing wolves and bears to the Alps'' as though both had arrived in crates sent by environmental officials. According to the European Commission the wolves came back by themselves and spread naturally across the Union. Nobody reintroduced them. Bears are a different case. They have been deliberately released in certain places, among them Trentino and the French Pyrenees \cite{ECLargeCarnivoreFacts2026}. Between 1999 and 2002 the Life Ursus project set ten wild bears from Slovenia free in Trentino, with the declared aim of founding a viable population in the central Alps \cite{TrentinoLifeUrsus}.

The distinction does not settle what should be done next. A natural arrival is no command to accept any number of animals, and human help does not make a population worthless. But it marks two kinds of responsibility. A government faced with a wandering wolf is managing a species that returned. A government that releases bears has acted deliberately, and it is directly answerable for the consequences it could foresee.

The return is substantial. The Commission estimates about 23,000 wolves in Europe in 2023, which is 35 percent more than in 2016 \cite{ECLargeCarnivoreFacts2026}. In a 2025 survey of 10,807 residents of EU countries that have large carnivores, support for the recovery of wolves, bears, and lynx generally outweighed opposition. Most respondents nevertheless opposed further growth of the populations and opposed hunting them too. And country people and city people differed much less than the usual culture-war story supposes \cite{ChapronCarnivoreOpinion2026}. Public opinion, in short, contains an unresolved wish. The animals should be back, and they should neither multiply nor be shot.

\section{What predators actually do}

Large predators are not ornaments. They kill wild deer and boar, leave carrion, compete with other predators, and change the behavior of their prey. Bears spread seeds. The restoration of a native species may also be a legitimate goal in itself, even where no dramatic chain of ecological effects can be shown. The biodiversity argument has a real biological core.

It has limits as well. The effects depend on prey density, terrain, hunting, forestry, roads, and settlement. A review of large carnivores in landscapes dominated by people warns that the popular picture of predators as automatic ecosystem engineers rests on a handful of comparatively wild places. Human activity can weaken, replace, or redirect the effects, so that outcomes vary greatly from one setting to another \cite{KuijperCarnivores2016}.

And the Alps are no empty wilderness awaiting completion. They are a worked landscape of managed forests, roads, ski lifts, villages, tourism, hunting, and summer pastures. Humans already control the deer, cut up the habitat, and decide which land stays open. In such a place ``wolves increase biodiversity'' is a hypothesis. It needs a location, a mechanism, a baseline, and a measurement. It may hold in one respect and be false or trivial in another. To count the returned wolf as one more native species is fair. To claim every downstream benefit without measuring any is not.

The opposite slogan is no more rational. That a predator kills livestock does not show it to be ecologically useless. That a sheep is domesticated does not make it ecologically irrelevant. Two complete fairy tales are on offer, the healing wolf and the worthless killer. The rational course is to decline both and ask what happens in this valley under these conditions.

\section{The sheep outside the frame}

The publicity photograph shows a wolf, alert, in the snow. It does not show the flock at dawn. It leaves out the animal still alive that has to be put down, the sheep that are simply gone, and the hours of searching a mountainside. It leaves out the farmer. These omissions are morally convenient. If the wolf is a sentient creature and no mere unit of population, the sheep cannot become a mere unit at the moment it is torn open.

The scale can be stated honestly. The European Commission warns that national reporting is not harmonized, and says that wolves kill at least 65,500 head of livestock a year in the EU, about 73 percent of them sheep and goats. It adds that the reported losses come to roughly 0.07 percent of the Union's sheep and goats \cite{ECLargeCarnivoreFacts2026}. Both figures are relevant, and neither cancels the other.

The percentage answers a continental question: wolves are not going to wipe out European sheep farming by arithmetic. The absolute number and its distribution answer a local one: for particular farms and districts the losses can be severe. An average spreads the dead animals over millions of sheep that were never at risk and over citizens who lost nothing. This is the missing-denominator problem that runs through all public policy. A cost that is small for society can be very large for the few who bear all of it.

Austria gives a closer view. For 2025 the Austrian Centre for Bear, Wolf, Lynx records 454 sheep or goats and 26 cattle killed by wolves. A further 654 sheep or goats were recorded as missing in connection with attacks. Compensation rules differ from one federal state to the next \cite{AustrianCarnivoreLosses2025}. These numbers make no case for extermination. They cannot be written off as a rounding error either, because they are part of the bill.

The bill is larger than the market value of the carcasses. It includes fences, night pens, guard dogs, shepherds, vets, paperwork, lost breeding stock, and the risk that extensive grazing will simply be given up. A report for the Alpine Convention notes that during the decades without large carnivores the old practices of coexistence were lost and herds were left unguarded. It also notes that existing subsidies often do not cover the standing cost of a shepherd or a dog \cite{AlpineConventionCarnivorePrevention2024}.

Mountain pasturing is more than a postcard of cowbells and green slopes. Grazing keeps open the meadows on which many other species depend. If predator policy speeds the retreat of marginal farms, a program announced in the name of biodiversity may damage one kind of biodiversity while promoting another. That possibility does not dictate the answer, but it rules out the easy one.

\section{Compassion at somebody else's expense}

The political asymmetry is now plain. The benefits of a large predator are widely shared. People like knowing that it exists. They may come across a track, and a native species has been restored. The tourist office gets a symbol, and there are whatever ecological effects there may be. Most of the beneficiaries pay next to nothing. The burdens fall on a few places and a few trades. A shepherd cannot settle his feed bill with the satisfaction felt in the capital that the mountains are wild again.

It does not follow that distant citizens have no say. Wildlife is a legitimate public good, and democracies often protect things that lie far from most voters. What follows is a principle of payment. If the public chooses to have predators, the public must buy the conditions of coexistence. Full compensation for verified losses is where the account opens. It also covers fencing and its upkeep, extra labor, trained shepherds, guard dogs, veterinary bills, missing animals, and prompt help after an attack. These are costs of a public policy. They are no evidence that a farmer has failed a moral examination.

Compensation is not enough by itself, because money after the event does not prevent the suffering or give back the time and the breeding lines. Prevention can cut losses substantially, especially when several measures are combined, but no method works everywhere or completely. Steep ground, scattered pastures, hikers, the size of the herd, and the habits of the local wolves all matter \cite{ECLargeCarnivoreFacts2026,AlpineConventionCarnivorePrevention2024}. ``Just put up a better fence'' is the sentimentalist's ``let them eat cake'' when the speaker will neither build the fence nor pay the person who inspects it in a storm.

Symbolic compassion carries a moral hazard. The advocate collects the identity of a kind person and exports the dirty work. The farmer handles the carcasses, changes his methods, fills in the forms, and carries the remaining risk. A grown-up policy staples the invoice to the decision. Those who want the benefit must be prepared to pay the cost.

\section{Neither extermination nor a fairy tale}

Sentimentality is not confined to the friends of the wolf. Its enemies have nursery stories of their own: a pastoral world of eternal harmony, a bloodthirsty invader, and a population in constant danger of death. Europe did not lose its large carnivores through calm cost-benefit analysis. Fear and hatred played their part, as did competition and the prestige of mastering nature. In correcting the romance of the predator one should take care not to fall into a romance of extermination.

Human safety deserves a sense of proportion. The Commission describes attacks by large carnivores as rare and reports no fatal wolf attack in Europe in this century. It acknowledges isolated attacks by bears and the need for procedures to deal with animals that have lost their fear of people \cite{ECLargeCarnivoreFacts2026}. ``Rare'' does not mean ``impossible,'' and an average is no comfort to the person on the spot. It means that policy should distinguish ordinary presence from identified dangerous behavior, and treat no animal as either a plush toy or a murderer.

A rational policy on large carnivores would contain at least the following:

\begin{enumerate}
\item explicit targets for population and range, set at the proper cross-border scale and revised as evidence comes in;

\item open monitoring of the predators, the livestock losses, the ecological effects, the performance of prevention, and where the public money goes;

\item preventive measures suited to local terrain and farming practice, with the labor as well as the equipment paid from public funds wherever the predators are a public objective;

\item prompt compensation for direct loss and for documented indirect loss, under rules simple enough to make reporting worth a farmer's time;

\item thresholds and lawful procedures, fixed in advance, for scaring off, capturing, or killing animals that repeatedly attack livestock, have lost their shyness, or are dangerous; and

\item a real voice for the communities that live with the policy, without giving any single interest an unconditional veto over a species that belongs to everyone.
\end{enumerate}

European law has moved cautiously toward management. Directive (EU) 2025/1237 transferred the wolf from Annex~IV to Annex~V of the Habitats Directive. Member states gain room for maneuver. The wolf does not lose its protection, since any exploitation must remain compatible with a favorable conservation status \cite{EUWolfProtectionDirective2025}. The change is no order to hunt, and it does not declare every population secure. It opens the way out of the sterile choice between absolute protection and eradication, toward management that someone has to answer for.

\section{The discipline of complete compassion}

Reason does not tell us to love the lamb and hate the wolf, or the reverse. It shows up the incompleteness of a compassion that sees only one of them. It asks whether the benefit to biodiversity has been measured and whether another investment in conservation would do more. It asks whether the farmer has a voice and whether the full cost is being met. And it asks whether the advocates would keep their enthusiasm if part of that cost turned up on their own bill.

The same discipline applies to the household pet. Affection for a dog or a cat is real and valuable. It is no certificate of ecological wisdom. The cub's face is a good place for attention to begin and a bad place for policy to end. A grown-up ethics has to look past the face, to the prey, the habitat, the farm, the budget, and the uncertain chain of causes. It has to see the people who are expected to make an attractive ideal work on difficult ground.

Sentimentality says, ``I care deeply; therefore my conclusion is good.'' Rational compassion says, ``Because I care, I am obliged to discover what my conclusion does.'' The first enjoys its innocence. The second accepts an invoice. Only the second is serious enough for wolves, sheep, bears, farmers, and the Alps they have to share.


% ===================== Chapter file ch13 =====================
\chapter[Bring Your Whole Country]{Bring Your Whole Country: The Stories Good People Tell and the Bills They Export}
\label{chap:responsible-goodness}

\epigraph{``It is good that you brought a bottle of red wine. It would have
been better if you had brought your whole country with you.''}{A director at
the Lebedev Institute, Moscow, in the final years of the Soviet Union}

\section{The bottle in Moscow}

Toward the end of the Soviet Union I visited the P.~N.~Lebedev Physical Institute in Moscow, known by its Russian initials as FIAN. It was the proud face of a state that could train formidable scientists and launch spacecraft. It also stood in an economy where a bottle of wine could be a small diplomatic event.

Gorbachev's anti-alcohol campaign had made drink especially hard to come by. The campaign had its reasons. Alcohol was wrecking Soviet health, families, and work, and the measures of 1985 did at first reduce drinking and deaths. They also cut legal production, lengthened the queues, bred moonshine, and deprived a state that lived on vodka revenue of its income. Verdicts differ according to which of these one counts \cite{BhattacharyaGathmannMiller2013,Reitan2001}. Soviet wit put it all in a nickname. The General Secretary became the \emph{mineralny sekretar}, the Mineral Secretary, because the official answer to drink was mineral water.

I had brought a bottle of red wine from the West. My Moscow colleagues produced another, carried home from a conference in Georgia. Neither was a grand bottle. In that room the two of them stood for travel, access, and the outside world. One of the institute's directors thanked me, and then said the sentence I have never forgotten: ``It is good that you brought a bottle of red wine. It would have been better if you had brought your whole country with you.''

It was funny because it joined friendship to an impossibility. It was also a whole political economy in one sentence. Why stop at the wine? If one system has plenty and the other has shortages, let the first pour itself into the second, with its shops, its hard currency, its machines, its administration, and its habits of work.

Could my small country have done that? Its wealth looked large when it took the form of one bottle shared among a few colleagues. Measured against the shortages of an empire it was nothing. Suppose the instruction had been carried out, and the stock, the budget, the housing, and the hospitals of a small country had been emptied into a system whose rules went on producing scarcity. The result would not have been two rich countries. The small one would have acquired the shortages of the large one. A gift can relieve a want. It cannot reproduce the order that made the gift available.

This last chapter starts from that distinction. A humane impulse is not yet a humane system, and a beautiful intention does not describe its effects. The moral quality of a decision cannot be read off the story its author tells. One has to ask what the decision does at scale and over time, once other people have reacted to it, and once the invoice has reached somebody who was not present for the applause.

\section{The animal that tells itself a life}

Human beings do not simply remember what happened to them. They arrange it. They pick out beginnings and turning points, injuries and rescues, and they decide what sort of person the hero was turning into. Psychologists call the result narrative identity, an evolving life story that rebuilds the past and imagines the future so as to give a life unity and purpose \cite{McAdamsMcLean2013}. The story need not be a lie, and it is never a transcript. My account of the Moscow wine selects one dinner out of thousands of hours and makes a parable of a director's joke.

The plots come from stock. Family, church, nation, profession, and cinema supply them: the self-made man, the survivor, the rebel, the victim who speaks at last, the citizen who stood on the right side of history. And the story is always told to an audience, even in solitude, where the audience is a parent, a rival, posterity, or God. Goffman described social life as the management of face \cite{goffman1967}. Narrative identity is face-work with a plot.

It follows that the wish to appear good is no superficial vice. It is part of the fabric. A reputation for honesty and generosity is what allows strangers to trust each other, and people do good partly because they are watched and remembered \cite{WuBallietVanLange2016}. A civilization in which nobody cared what anyone thought would be much the poorer.

The danger lies at the seam between the performed virtue and the material world. Stories reward acts that can be seen. Consequences are scattered and delayed. Often they are merely statistical, and strangers bear them. The hero collects his moral credit now, and the cost turns up later, in a chapter he will never read. When that separation becomes systematic, the story of goodness can outlive conduct that makes the world worse.

Hardly anyone casts himself as the villain. Even cruelty gets a justifying plot, in which severity was necessary and the circumstances were exceptional. A movement admits ``mistakes'' and keeps the purity of its direction. An institution explains each failure by insufficient commitment to the policy that failed. Three pressures make the moral story especially hard to correct.

The first is \emph{moral grandstanding}, popularly called virtue signaling: moral talk used partly to gain standing by showing off one's purity or commitment. The charge is easy to abuse, and to hurl it at an opponent can itself be a way of showing off. Moral speech can be sincere, rewarded, and useful all at once. The question to ask is whether the incentive to produce the signal has drawn attention away from the result \cite{TosiWarmke2020}.

The second is \emph{sentimentality}, which is feeling protected from the information that would complicate it. It freezes the frame at the rescue, the welcome, the spared animal, or the clean national statistic. The torn sheep, the displaced tenant, the exhausted teacher, and the factory abroad are outside the frame and cannot compete.

The third is the wish to \emph{walk beautifully through a beautiful and good world}. This is an aesthetic demand dressed up as a statement of fact. The world has to be innocent enough for kindness never to face a tragic choice. Everyone can be welcomed, every predator protected, every worker secured, and every unit of energy saved, without any material conflict among these aims. When a cost does appear, it is put down to bad implementation or bad people. It is never allowed to be a property of the choice.

Conformity fuses the three. Asch showed how a unanimous group can extract public agreement with an obviously wrong judgment \cite{Asch1956}. In politics the facts are seldom as simple as the length of a line, so the pressure is easier to rationalize. People who privately doubt a norm infer from one another's compliance that everybody else accepts it \cite{PrenticeMiller1993}, and they profess preferences they do not hold \cite{Kuran1995}. What results is worse than hypocrisy. It is a failure of feedback. Officials hear an approval that the cost of dissent has manufactured, and the story grows more unanimous with every private doubt it silences.

\section{Why visible goodness sometimes works}

A critique of moral display becomes childish if it holds that only hidden virtue is real. Cooperation needs visibility. Promises have to be known, help remembered, and betrayal reported. The honest shopkeeper wins custom, and the scientist who corrects his own error at some cost becomes a better witness. A norm has force when keeping it is seen to be rewarded.

The same mechanism explains the corruption. When the reward attaches to the gesture and not the outcome, people optimize the gesture. A program advertises how many were admitted and omits how many completed. A company announces a donation smaller than the budget for publicizing it. A government counts its domestic emissions and imports the carbon-intensive goods. A campaign counts the predators restored and leaves out the livestock killed. The signal stays favorable because the accounting boundary was drawn to make it so.

Research on \emph{moral licensing} adds a twist. People who have earned moral credit in one matter allow themselves more latitude in another. The effect in the laboratory is small and real \cite{BlankenVanDeVenZeelenberg2015}. In institutions it matters more. A state that has declared itself humanitarian treats criticism with contempt. A movement sure of its anti-racism stops noticing that it sorts individuals by ancestry. A firm with a celebrated sustainability program takes questions about its suppliers as hostile acts. Moral credit is being used as collateral against scrutiny.

The remedy is to tie the signal to a reasonably complete account of consequences, and moral language can then stay. A hospital is not humane because it has hung out a banner about universal care. It is humane if it treats real patients competently and rations scarce attention by defensible rules. It must also keep its staff from collapse and tell the public what it cannot do.

\section{Kindness with a memory}

Robert Axelrod's computer tournaments for the repeated Prisoner's Dilemma are remembered for one charming result. A simple strategy called TIT FOR TAT, which cooperates first and then copies the other player's last move, did remarkably well. The popular moral is that niceness wins \cite{Axelrod-80a,Axelrod1390,Axelrod-85}.

That moral leaves out three quarters of the mechanism. Axelrod found that successful strategies were \emph{nice}, \emph{retaliatory}, \emph{forgiving}, and \emph{clear}. They never defected first. They answered exploitation and did not reward it. They went back to cooperating when the other side did. And their behavior was plain enough for the other side to learn the pattern.

Take away any one of the four and the lesson changes. Niceness that never retaliates is simply an invitation to exploiters. Where retaliation is never followed by forgiveness, the result is a feud. Forgiveness that is not clear looks like inconsistency and teaches nothing. And a clear strategy with no opening kindness is efficient predation. The achievement lies in the combination. Real societies are noisier than tournaments, and later work found a use for a little extra generosity to break accidental cycles of revenge. But the principle holds. Forgiveness is valuable when it restores a game of reciprocity. It does harm when it reliably subsidizes defection.

A boundary is therefore no enemy of cooperation. It is usually what makes cooperation believable. A parent who never says no teaches that demands are free. A welfare system that cannot tell contribution, misfortune, and deliberate extraction apart uses up the consent on which it lives. An asylum procedure that neither grants nor refuses promptly rewards staying power over the merits of the case. A conservation rule that protects every predator whatever its numbers and behavior has turned management into ritual.

A cooperative order has to be able to say four things: yes, no, not yet, and yes again. The sentimental order can say only yes, until it collapses into an indiscriminate no. Because discriminating limits were stigmatized, the limits come late, and they come as anger, panic, and collective punishment.

\section{The missing players}

The Prisoner's Dilemma has two players and a visible table of payoffs. Public morality usually hides both. The chooser gets the moral benefit, and the payer is somewhere else. A third party adjusts his behavior, and people not yet born inherit the outcome. A policy can thus be cooperation with the first player and defection against the third.

A town promises a generous service and does not fund the staff, the buildings, or the pensions. Today's official is kind to today's applicant. The nurse works the extra shift, the next applicant joins a longer queue, and a future council discovers the liability. Or a consumer boycotts a product made by an objectionable process at home and buys an imported substitute whose production he cannot see. If the harm has only moved, he has cleaned up a narrative and left the world as it was.

At least six roles need to be kept apart:

\begin{enumerate}
\item the \emph{decision-maker}, who holds authority and usually takes the moral credit;
\item the \emph{immediate beneficiary}, whose visible need motivates the act;
\item the \emph{direct payer}, who supplies the money, labor, land, time, or risk;
\item the \emph{displaced claimant}, who would otherwise have received the scarce good;
\item the \emph{strategic responder}, who changes his behavior because the rule has changed; and
\item the \emph{future bearer}, who inherits the debt, the institutions, and the precedent.
\end{enumerate}

One person can fill several roles. The migrant may become the nurse and the taxpayer, and the farmer may also be glad that a native predator is back. A complete account does not prejudge the answer. It stops a favorable answer from being bought by omission.

Max Weber's two ethics name the two questions involved \cite{Weber1919}. The ethic of conviction asks whether the act expresses a principle worthy of the actor. The ethic of responsibility asks what consequences can be foreseen, and to whom the actor is answerable for them. Conviction without responsibility is moral exhibitionism. The actor welcomed, protected, banned, or subsidized, and therefore his hands are clean, and if disorder followed, somebody else must have betrayed the principle. Responsibility without conviction is opportunism. The hard kind of goodness joins the two. It begins with a commitment and accepts liability for the means and the trade-offs. It will not treat the naming of a value as its delivery.

This matters most in public life, because the state spends money it has taken by compulsion and makes rules for people who never agreed to them. A private host who puts a stranger in his spare room lives with the consequences every day. A minister who announces hospitality is assigning housing, schooling, policing, medical care, and taxes across a whole population. Both acts can be defended. The second is not the first writ large. Public virtue comes with an invoice enforced by law, and whoever governs for the sake of his own purity will leave it unread.

\section{The grammar of scarcity}

``Zero-sum'' is said too easily. Much of economic life is not a fixed pie. Trade, knowledge, and investment let several parties gain at once, as Chapter~\ref{chap:inheritance-paradox} argued in the case of family advantage. A society that sees theft in every difference destroys the cooperation that makes the pie grow.

The opposite error is now the fashionable one. Because the economy can grow, every particular claim is treated as though it competed with nothing. But at a given hour there is one bed in the intensive-care unit, one teacher before the class, one flat falling vacant, one hectare of pasture, one official working through the files, and one euro that can be spent once. Capacity may expand tomorrow. Today's allocation still has a displaced claimant.

Four structures should be kept apart. In a \emph{positive-sum exchange} everybody gains: the skilled migrant fills a vacancy, earns, pays taxes, and supplies a service. In a \emph{rival allocation} a fixed resource goes to one claimant and therefore not to another. This need not be unjust, since triage is unavoidable, but justice starts with naming the queue. In a \emph{common-pool problem} each user takes a private benefit and spreads the depletion over all. Hardin stressed the tragedy, and Ostrom showed that communities can sometimes govern a commons through local rules, monitoring, and graduated penalties \cite{Hardin1968,Ostrom1990}. In a \emph{negative-sum cascade}, such as a bank run or a collapsing public service, nearly everyone ends up worse off. The task there is to keep the oven from being wrecked, and dividing the pie comes afterwards.

The story of goodness slides from one of these to another. It describes a rival allocation as though the new capacity already existed. It treats an open commons as though each further claim were free. It supposes that more generosity can stop a cascade. Arithmetic is felt as hostility, because arithmetic brings in the person outside the frame.

And a gift is not a production function. The bottle at FIAN was a real good and no mistake. The mistake would be to conclude that because the single gift was good, the universal transfer must be good too. This is the fallacy of scale. One new child is absorbed into a class and enriches it. Fifty who arrive at once without the language of instruction change the class. One wolf in a large territory is not a pack working the same flocks night after night. An effect can change its sign as the scale grows.

Wealth, moreover, is not a warehouse whose doors need only be opened. It is a flow, kept going by property rules, knowledge, energy, trust, incentives, and political stability. Transfer part of the output and leave the productive order intact, and generosity can go on for ever. Transfer claims faster than the order can reproduce the output, and the source dries up. Institutions cannot be poured from a bottle, because they consist largely of expectations about what other people will do. Will contracts be kept? Can officials be challenged? Is effort rewarded and failure admitted? Such expectations change through experience, and no decree will change them. Aid that ignores this can do damage. It can ruin local sellers, raise the prize for capturing a corrupt administration, and remove the reason to maintain local capacity. None of this argues against aid. It argues that kindness should be designed as an intervention in a system, where the usual story shows only a transfer between two still photographs.

\section{Oil not burned here}

Hans-Werner Sinn's ``green paradox,'' already met in Chapter~\ref{chap:moving-climate}, is the cleanest example of the missing system. Germany reduces its demand for oil. In the domestic books less oil is burned, and the sacrifice shows in higher prices, different cars, and industrial adjustment. But oil is traded worldwide, and its owners plan over decades. Lower demand in one region lowers the world price below what it would have been, and others buy more. The announcement of a future transition gives owners a reason to pump sooner \cite{Sinn2015}.

This is a warning and no iron law. The size of the effect, and even its direction, depend on extraction costs, expectations, the breadth of the coalition, and the alternatives available. The literature on rebound does not support the idea that efficiency normally backfires completely, though cheaper use does regularly eat up part of the saving \cite{GillinghamRapsonWagner2016}. Recent work finds carbon leakage from unilateral policy, varying widely by sector and country, and nowhere near a universal one-for-one offset \cite{HemmerleKrusePisu2025}.

The honest conclusion asks more of us than either slogan does. Germany's saving is not meaningless, since it can develop technology and help to build a wider coalition. But the sacrifice does not validate itself. The object is the atmosphere, and the national moral account is beside the point. The discipline consists in following the barrel after it has left the German story.

Leakage is the shadow of all moral geography. Ban an activity here and it moves next door. Raise the standard for visible suppliers and production passes to invisible subcontractors. The cost can also be exported in time, through deferred maintenance, debt, inflated school marks, or a ban imposed before any replacement exists. Time is one more foreign country to which costs can be sent. Three borders should be on every map: territory, time, and social class. A professional in the capital supports a policy whose price is paid in a rural district, by a manual trade, after the next election. His rural opponent resists a public good and keeps his private subsidy. Both tell flattering stories about themselves. Complete compassion begins when each has to include the other's ledger.

\section{The wolf seen from the desk}

The return of wolves and bears is a nearly perfect theater of selective storytelling, and Chapter~\ref{chap:sentimentality-rationality} went through the evidence. Here I am concerned only with the split between the person who collects the moral satisfaction and the person who does the daily coexisting.

My good friend Herbert, a professor of political science at the University of Vienna, used to come back with astonishment to a story from the United States. A mountain lion had killed a young mother. The authorities tracked the lioness and shot her. She was a mother too, and left cubs. Two collections followed, one for the dead woman's children and one for the dead cat's offspring. As Herbert remembered it, a huge sum came in for the cubs and almost nothing for the children. He did not tell the story because he wanted the cubs to starve. He told it because the imbalance seemed to him a miniature of an entire moral culture.

The story is true, though memory had sharpened the contrast. Barbara Schoener, a forty-year-old marathon runner and mother of two, was killed by a cougar on a trail in El Dorado County, California, on 23 April 1994. Trackers shot the lioness a week later. Searchers then found a single dehydrated cub, seven weeks old, which went to the Folsom City Zoo \cite{LATimesSchoenerCubs1994,DeseretSchoenerCub1994}. After a week the appeal connected with the cub stood at \$21,000. The trust fund for Schoener's children, Andrew, who was eight, and Anna, who was five, stood at \$9,000. That is not nothing, and it is less than half. Even this comparison needs its denominator. Of the cub's total, \$15,000 was a pledge from the zoo's own nonprofit society toward an enclosure \cite{ChicagoTribuneSchoenerFunds1994}. No thousands of private donors had preferred the cat to the children. Two unlike things were being compared, an institutional building pledge and a family trust, as though they were a referendum on the worth of animal and human life. The disparity was still real enough to make the news. When it did, indignation brought in new donations, and the children's fund more than doubled \cite{NYTimesSchoenerFunds1994,WashingtonPostSchoenerFunds1994}.

The correction teaches as much as the original imbalance. Nobody had solemnly decided that a cougar's cub was worth more than two bereaved children. People had met two stories that were differently written. The cub was the ideal object of rescue. It was tiny, photogenic, helpless, and innocent of its mother's deed, and one visible act could save it. The children's loss was deeper and made a poorer story. They had a father, relatives, a bank account, and a future whose needs could not be squeezed into the picture of an animal small enough to hold. Attention followed the clarity of the image and paid little regard to the size of the loss.

Then a third story appeared, about a society that gives more to the killer's cub than to the victim's children. It was stronger than the bookkeeping it was built on. It made the children visible and turned giving to them into a protest, and the donations moved. The sequence shows both the power of sentimental selection and how unstable it is. A new frame did more than add information. It changed who looked like the innocent party in need of a public show of concern.

Nothing here calls for contempt of animals. The cub killed nobody, and saving it was decent. The error begins when the irresistible face is allowed to fill the whole moral field. Herbert's anecdote survives its factual correction because its point was never that people care too much. His point was that their caring is captured by whichever story makes goodness easiest to feel and to display.

With the large carnivores it is the same. The case for protecting them is real, and so are the costs. The costs fall on the people for whom the wolf is no idea \cite{ECLargeCarnivoreFacts2026,KuijperCarnivores2016}. The central confusion is between protecting a species and protecting every individual animal in every place. Once a population is established, conservation means management: monitoring, prevention, compensation, a local voice, and the lawful removal of dangerous animals. To refuse management is to keep the pure story and export the bloody work. Extermination, for its part, is not realism. It merely chooses the sheep as the only face. Complete compassion may still decide for the wolf. It can no longer pretend that it decided against nobody.

\section{Migration begins with categories}

Migration is the subject on which the demand for a full account is most readily treated as an offense. The immediate beneficiary has a face and often a terrible story. The displaced claimant is a statistic. He is the family still waiting for a flat, the child whose teacher's attention is divided, the patient whose appointment slips, or the resident who takes a different way home after dark. Because one side can be seen and the other is scattered, the language of welcome takes moral precedence before anyone has talked about scale. As Chapter~\ref{chap:two-plus-two-five} showed, ``manageable'' tells us nothing until its categories, its denominator, and its limits are given.

The first discipline is to distinguish. ``Migrant'' covers the engineer hired for a vacancy, the EU citizen exercising free movement, the student, the spouse, the recognized refugee, the applicant awaiting a decision, the rejected applicant who cannot or will not leave, the seasonal worker, and the unauthorized entrant. They have different rights, skills, motives, and probable effects. ``Welfare migrant'' is just as crude if it treats every move toward a richer country as fraud. Safety, family, work, ambition, and access to public services can all figure in one decision. People have mixed motives, and foreigners are no different in this.

The receiving country nonetheless has a legitimate interest in incentives. Where admission, asylum decisions, deportation, access to work, and entitlement to benefits are slow and inconsistent, behavior adapts to the rule as practiced, whatever the rule as announced may say. A right meant for exceptional protection becomes a general channel of immigration if claims take years to decide and refusals are seldom enforced. Conversely, to forbid applicants to work for long periods is to manufacture dependency and then cite it against them.

The second discipline is time. A refugee costs money at first, and a recruited worker may pay in from the first month. Children consume schooling before they earn, as all children do. Some families integrate strongly over the generations and others stay weakly attached to school and work, and a single average hides both. The third is selection. Among immigrants in Germany aged twenty-five to thirty-four in 2025, the share with a university degree was close to that of the whole age group, while the share with no vocational qualification and not in training was more than twice as high \cite{DestatisMigration2026}. ``They are skilled'' and ``they are unskilled'' can each be made true by choosing a subgroup. Each is false as a description of the whole.

Moral language that refuses these distinctions is no kinder for it. It makes selection for the labor market impossible and neglects those who need intensive help. It also lets the visible failures color the public's view of those who are working and learning. Categories are no bureaucratic insult. They are how a state keeps promises made on different grounds from destroying one another.

A generous welfare state, likewise, is not a heap of money waiting to be divided. It is a repeated game between generations. People pay in when they can, draw out in childhood, sickness, unemployment, and old age, and accept that some stranger may receive more than they do. The game requires a bounded community of obligation. The boundary need not be ethnic or closed for good, since citizenship, residence, contribution, need, and humanitarian protection can each ground a claim. But a system that will not say when a claim begins and what it contains has turned solidarity into open access. ``There is enough money'' misses the scarce resource, which is trust. People will finance large transfers that they think fair, and rebel against small ones that they think arbitrary. Agitators can damage trust. So can official euphemism. Citizens who are told that every visible strain is imaginary learn to disbelieve true reassurances along with false ones.

Axelrod's four qualities transfer rather well. A welfare state should be \emph{nice}, helping the desperate without waiting for them to become useful. It should be \emph{retaliatory} in the sober, institutional sense that fraud and refusal of reasonable duties have consequences. It should be \emph{forgiving}, so that a temporary failure does not mean civil death. And it should be \emph{clear} about what is expected and how to appeal. Migration strengthens the game when newcomers work, found firms, and join in reciprocity. It weakens the game when arrivals outrun integration or when policy makes it rational never to participate. To make either outcome the essence of ``the migrant,'' as savior in one story and parasite in the other, is the same narrative error with the parts reversed.

\section{Capacity is not a metaphor}

``We can manage this'' states no moral principle. It is a claim about capacity, and it should be measured like one. How many school places are there, and how many language teachers, flats, asylum officers, doctors, interpreters, police, and court dates? How fast can more be added? Which is the bottleneck? And who goes short until expansion catches up?

\emph{Schools} show the mechanism. A child who does not yet understand the language of instruction is guilty of nothing and is entitled to suitable teaching. That teaching takes staff, time, and small groups. When many such children arrive quickly, to pretend that nothing extra is needed harms them and their classmates. German school enrollment has risen for several years, mainly through immigration, which is a fact for planners and no argument for or against the children \cite{DestatisSchools2025}.

\emph{Housing} is more rigid. A person can cross a border in a day, and an apartment block takes years. Where supply cannot respond, added demand means higher rents and more crowding. Immigration is not the only cause of the shortage. Zoning, building costs, interest rates, and decades of underbuilding all contribute. But it is just as false to say that added demand has no effect because the shortage had other causes.

\emph{Health care} cuts both ways. New residents are patients, and they are also doctors, nurses, cleaners, and taxpayers. Foreign citizens made up about a fifth of nursing staff in Germany in 2025, and recent growth of employment in the sector depended heavily on them \cite{BundesagenturPflege2026}. The picture of the migrant who only fills the waiting room leaves out the migrant on the night shift. The opposite picture, in which every arrival helps to solve the staff shortage, leaves out the language barrier, the dependents, and the ethics of taking nurses from poorer countries.

\emph{Public order} calls for the same exactness. German police statistics show some categories of recent migrants overrepresented among suspects for particular offenses. A suspect is not a convict. Raw comparisons are distorted by age, sex, city residence, reporting habits, and offenses that only foreigners can commit. Migrants are also victims. The Federal Criminal Police Office spells out these limitations itself \cite{BKAFluchtmigration2023}. To refuse to mention the overrepresentation is dishonest. To turn it into hereditary collective guilt is dishonest and unjust as well.

Analysis of capacity protects individuals from that kind of collectivization. If a district lacks police, youth work, housing, and language teaching, the failure belongs to mechanisms that can be repaired, and no mythical national character is responsible. If a particular offender is dangerous, the law should deal with him, and should neither sentimentalize nor demonize a category. That is what the rule of law means: nobody is judged by a story about his group.

\section{Germany as a demographic stress test}
\label{sec:germany-demographic-stress-test}

Two finished stories are told about Germany. In one, immigration is pure enrichment, which makes up for the missing births and supplies the workers an aging country needs. In the other it is demographic replacement, and every difference in employment, schooling, crime, or custom is proof of irreversible national loss. An earlier chapter stated the second story at full strength. Each of the two starts from observable facts and ends by pretending that its verdict was observed along with them. It is better to take the measurements one at a time.

\emph{Births and age.} In 2025 Germany recorded 654,241 births and a fertility rate of 1.32 children per woman. Deaths exceeded births by about 352,000. Net immigration of about 235,000 did not close the gap, and the population fell slightly \cite{DestatisFertility2026,DestatisPopulation2026}. In every variant of the official projection a quarter of the population will be sixty-seven or older by 2035, and the working-age population falls by at least four million by 2070 \cite{DestatisProjection2025}. A projection is a conditional statement, but the retirement of people who are already alive is not very uncertain. Long life is an achievement and should not be called misery. A low birth rate is no order to citizens to breed, and the state cannot produce families without invading the freedom it exists to protect. It can remove obstacles such as scarce housing, unreliable child care, and insecure early careers. Birth policy works over decades if it works at all, and pension policy cannot wait for workers who have not been born.

\emph{What migration can offset.} It can make the country younger, enlarge the workforce, and fill particular vacancies. It cannot cancel aging by arithmetic. Migrants grow old too. Some leave, and some arrive as dependents. The honest claim is conditional. Migration relieves a labor shortage when the entrants are of working age, find work, learn German, and stay. It worsens local shortages when admissions outrun housing, schools, and administration. Neither result follows from foreign birth. Both depend on selection, scale, timing, and policy.

\emph{Employment.} The record refutes instant optimism and also permanent condemnation. In the IAB--BAMF--SOEP panel, the employment rate of refugees who arrived between 2013 and 2019 reached 68\% after eight years. For the 2015 cohort in 2022 it was 75\% among men and 31\% among women. Faster asylum decisions, earlier permission to work, and completed language courses went with better results \cite{IABRefugeeEmployment2024}. That is substantial integration, with a very large gap between the sexes. Years of administrative waiting, on the other hand, manufacture the very dependency that is later cited as proof that integration cannot work. Obligations are legitimate as well. Anyone able to learn or work may be required to take part, on reasonable terms, as a condition of particular benefits. Reciprocity of this kind is no hostility, and it protects the consent on which generosity depends.

\emph{Schools and buildings.} German is no ethnic test. It is the instrument with which a child follows a lesson, a patient describes a pain, and a citizen contests an official decision. The IQB education survey found larger losses of competence among children with an immigrant background, above all the first generation, and it identified social background and the language spoken at home as major factors \cite{KMKIQBBildungstrend2022}. That calls for early and intensive German teaching and for standards that stay visible. It does not justify treating a child's origin as a diagnosis. Housing runs on a different clock. Germany completed about 206,600 dwellings in 2025, 18\% fewer than the year before, and the average time from permit to completion had lengthened to twenty-seven months \cite{DestatisHousing2026}. An empty house in a shrinking rural district is not a flat near a job in Munich. The federal government decides on admission, and the duties fall on municipalities. They need predictable numbers, money that follows the person, and the right to report a bottleneck without being called heartless.

\emph{Crime without a collective defendant.} The federal police statistics record that non-German suspects are overrepresented in several groups of offenses. That should be published and studied, since it bears on policing, prevention, and admission policy \cite{BKAPKS2024}. The same statistics limit what can be inferred. ``Non-German'' is a category of citizenship. It includes tourists and excludes the naturalized, and offenses against immigration law must be taken out of any comparison. Age, sex, poverty, and city life explain part of the excess. They do not make an assault unreal. The defensible rule is individual. Publish figures for each type of offense, with denominators. Prosecute the offender, and assign no guilt to a nationality or a religion. Count migrant victims as carefully as any others. Where the law allows, remove serious foreign offenders after an individual decision that can be reviewed. Equal law grants no immunity in the name of vulnerability and permits no inherited suspicion in the name of realism.

\emph{A common constitutional culture.} A country needs more in common than an address, and what it has in common need not be ancestry. The Basic Law binds everyone to human dignity, equality before the law, equality of men and women, freedom of belief and expression, and the rule of law \cite{Grundgesetz2025}. These are the corrigible public powers defended in Chapter~\ref{chap:government-and-liberty}. The consequences should be stated without embarrassment. A husband has no cultural exemption from his wife's equality. No family may force a marriage. Religious offense confers no right to threaten a speaker. Antisemitic intimidation is intimidation whether it comes from a German neo-Nazi or an Islamist newcomer. Anyone may change or abandon a religion. Courts decide rights, and clans and clergy do not. None of this requires the state to supervise dinner tables. Integration is sufficient when people can communicate, support themselves where they are able, obey the law, and accept the equal citizenship of others. The state has duties in return. It must decide claims in a credible time, enforce final refusals, and select where selection is the purpose. It must fund the language teaching and the schools, and it must never describe its choice as free of cost.

\emph{Correction without prophecy.} Germany cannot restore a photograph of the late twentieth century, and no honest policy promises to. Nor does every demographic change deserve to be renamed progress. Some communities have lost familiar institutions. Some migrants have failed to integrate, or refused to. Many others work, raise children, and belong to the ordinary life of the country. All of these observations hold together. The useful question is which rules, at which scale, produce employment, command of the language, safety, housing, and constitutional loyalty. Whether Germany has been ``abolished,'' ``replaced,'' ``saved,'' or ``enriched'' is not a useful question. The government can publish employment by cohort after three, five, and ten years, along with language results, school capacity, crime rates by offense, and the number of returns after final decisions. Policies can then carry review dates and conditions of failure. A country is not preserved by declaring it lost, and integration is not achieved by declaring it complete.

\section{Fear, and the temptation to ban}

Policy is experienced locally. The minister sees annual totals, and the resident sees the bus stop, the school gate, and the railway station at night. A group of idle young men in a public place may be harmless, annoying, or dangerous. A resident's fear proves nothing statistically, and a national average does not cancel what she saw. Fear is evidence about perceived order. It should be investigated. It is no oracle, and it is no sin either.

Paternalistic reassurance increases the fear it is meant to calm. When every worry is answered with a lecture on prejudice, residents conclude that officials do not know the neighborhood or do not care, and rumor takes over the job that official speech has abandoned. Politics follows. Voters who believe that the ordinary parties will not regulate migration turn to parties that promise to. If the whole field of restriction is declared morally off limits, the most reckless actors are handed a monopoly on a real concern. This is how beautiful goodness becomes a negative-sum cascade. The center displays tolerance and neglects limits, confidence falls, and protest parties grow. The center then treats their growth as proof that the public must be instructed, watched, or excluded.

At this point Popper's paradox is brought out as a weapon, in the form ``the tolerant must suppress the intolerant.'' As Chapter~\ref{chap:government-and-liberty} showed, his actual argument is far more restrained. Argument comes first, and the right of self-defense arises against movements that put intimidation and violence in the place of discussion \cite{Popper-TOSAIE}. A party can cross that line. It may organize violence, seek a legal hierarchy of ethnic groups, or prepare to abolish elections. That calls for evidence and a lawful response. But to oppose open-ended immigration, to demand deportation after a final court ruling, to criticize multiculturalism, or to propose a lower intake is not a refusal of democratic argument. To treat disagreement over policy as proof of intolerance destroys the distinction on which Popper's case depends.

German constitutional law offers a sober example. In 2017 the Federal Constitutional Court found the aims of the National Democratic Party to be unconstitutional and still refused to ban it, because there was no evidence that the party had any prospect of achieving them \cite{BVerfGNPD2017}. The court did not admire the party. It insisted that prohibition must meet strict criteria and must not become an act of symbolic hygiene. A ban can be necessary. It can also confirm a movement's picture of itself as the only forbidden teller of truth, drive its networks underground, and spare the established parties the labor of answering their voters. The paradox runs in both directions. Unlimited tolerance lets coercive movements end tolerance, and unlimited ``antifascism'' lets officials end pluralism by calling every opponent a fascist.

\section{Nietzsche, the Samaritan, and Shen Te}

Nietzsche's critique of Christian and European morality is uncomfortable because he will not accept a value's account of itself. Pity claims to serve the sufferer, and humility claims to renounce power. He asks other questions. Who gains authority when this value prevails? What does it do to vitality, excellence, and responsibility? \cite{Nietzsche-GM,Nietzsche-JGB} Chapter~\ref{chap:weaponized-original-sin} traced one grammar of inherited guilt. Nietzsche's question is what kinds of rule become possible once guilt can no longer be paid off.

The bearing on secular politics is not that environmentalists and advocates of migration are Christians in disguise. Moral vocabularies outlive their theology. Innocence attaches to the victim. Guilt attaches to the strong, the prosperous, the established, and the native. Sacrifice proves purity, and doubt is hardness of heart. A prosperous society demonstrates its goodness by renouncing, admitting, and paying, and whoever asks whether the sacrifice does any good has shown himself morally defective. Nietzsche teaches a useful suspicion, that altruism can be a will to power. The advocate who speaks for a distant victim gains authority over a nearby payer. The administrator of compassion decides who is guilty and who is deserving, and what penance is due. A politics that announces the end of hierarchy installs a hierarchy of moral interpreters.

But a genealogy refutes nothing. Compassion can come from strength, and equality before the law protects the exceptional person too. A society indifferent to weakness leaves everyone one accident away from abandonment. Nietzsche's hammer is good for testing idols and of no use for designing hospitals. He should be used for diagnosis. One may ask whether pity has become a status system and whether a policy needs the weakness it claims to cure. Then one should go back to the evidence, remembering that suspicion can be a performance of superiority too.

Christianity has its own answer, and it is a hard one for performers of virtue. The man who gives in order to be seen has had his reward. The steward is judged by what he did with what was entrusted to him. He gets no credit for having given it all away. The Good Samaritan is regularly enlisted as a parable of unlimited public admission, and the text is more concrete than that. A traveler finds an injured stranger, treats him, carries him to an inn, pays the innkeeper out of his own purse, and promises to settle any further cost. The parable destroys the excuse that membership of another group cancels an immediate duty of mercy. It lays down no border regime and no annual quota, and it creates no entitlement to be paid for by unnamed third parties. Its force comes from the helper's taking the burden on himself. Political communities may rightly pool costs that individuals could not carry. But public compassion should keep the Samaritan's contact with the bill. And love seeks the neighbor's recovery. A helper who needs the other to stay helpless has made an idol of helping.

Brecht's Shen Te, the heroine of Chapter~\ref{chap:small-wickedness}, puts the matter on the stage. She can remain good only by inventing the stern cousin who says no \cite{Brecht-GMvS}. Modern states hide Shui Ta inside the administration. Politicians proclaim rights, and anonymous officials define who is eligible. Campaigners denounce exclusion while the budget office does the rationing. Courts enlarge principles and the waiting lists grow longer. The public story stays with Shen Te while the institution does the refusing, and the author of the promise enjoys the applause and blames the clerk for the limit. An honest politics would put both characters on the stage together. Every new right should name the official who has to deliver it, the claimant who moves back in the queue, and the tax that pays for it. Nobody is asked to become a bad person. The point is to stop handing the necessary part of goodness to someone who will later be denounced for carrying it out.

\section{Four small disasters}

\emph{The second disaster.} After an earthquake or a flood people fill boxes with clothes, blankets, toys, and tins. The boxes prove that sympathy became action, and the photograph of the collection point makes a satisfying ending. For the relief organization that is where the work begins. Unrequested goods have to be shipped, cleared, unloaded, sorted, stored, and matched to needs, or thrown away. A container of unsuitable clothing takes up the place in a damaged port that medicine or water equipment needed. Red Cross societies ask for money and call the flood of unwanted goods ``the second disaster'' \cite{NZRedCrossSecondDisaster2026}. The donor is sincere. But a parcel maximizes his feeling that something concrete was done. Cash to a competent organization means giving up control. It makes a poor photograph, and it does far more good. The recipient should define the need, and the donor's wardrobe should have no say in it. Responsible compassion lets the sufferer interrupt the helper's favorite story.

\emph{The shirt given twice.} A bag of old clothes in a charity bin seems to do two good things. The donor's wardrobe is cleared without waste, and a poor person gets cheap clothing. Sometimes that is exactly what happens. But in the receiving country the shirt may be a commercial import sold by weight. It competes with local tailors, and it ends as rubbish that the local dumps cannot handle. Garth Frazer found that imports of used clothing explained a considerable part of the decline of clothing production and employment in African countries over the period he studied \cite{Frazer2008}. It does not follow that a ban would make the recipients richer, since consumers would pay more and traders would lose their living. What follows is that ``donation'' describes only part of a market process. The moral unit is larger than the moment at which the donor lets go.

\emph{The protected tenant and the absent apartment.} Rent control shields a sitting tenant from a sudden increase, and the benefit is concrete. A person keeps a home, a school, and a neighborhood. The person nobody sees is the future tenant. If owners respond by converting flats, withdrawing them, neglecting them, or not building, the protected flat stays in view and the absent one cannot be photographed. A well-known study of San Francisco's 1994 extension of rent control found that covered tenants gained stability and landlords cut the supply of rental housing, which probably pushed market rents up and worked against the aim of the policy \cite{DiamondMcQuadeQian2019}. Other rules in other places behave differently. The case illustrates an asymmetry. Sitting tenants are organized and have names, and future ones have neither. And rationing does not disappear when the price is fixed. It takes other forms, such as queues, connections, key money, and inherited tenancies.

\emph{The green fuel that ate a field.} Biofuel looked like an accounting miracle. Burn a liter grown in a field in place of a liter pumped from a well, and the carbon the plants absorbed offsets the carbon released. The boundary drawn round the fuel plant left out the land. If cropland is switched from food to fuel, food is grown somewhere else. If forest or grassland is cleared to grow it, the carbon stored there goes into the air. A prominent analysis of 2008 argued that this indirect change of land use could create a carbon debt large enough to reverse the expected benefit of certain crop fuels \cite{SearchingerEtAl2008}. The magnitudes are disputed, and the dispute makes my point, because the answer changes with the boundary of the system. The label ``renewable'' cannot do the analysis.

\section{Five ways to export an invoice}

The examples differ and the structure is the same. A decision exports its invoice when the actor who collects the moral credit is separated from the person, place, or time that bears a foreseeable cost. There are five common routes.

\emph{Fiscal.} The promise is paid for through general taxation, inflation, debt, or a mandate that someone else must fund. Collective finance is perfectly legitimate. The export occurs when the author of the promise presents the benefit as his personal virtue and hides the compulsory contribution.

\emph{Spatial.} The emissions, the waste, the production, the predation, or the social strain moves outside the district whose statistics are being celebrated. The clean city imports its electricity and its goods. The urban voter assigns the cost of coexistence to the pastoral fringe.

\emph{Temporal.} Debt, deferred repairs, pension promises, and eroded skills are left to those who come after. The building still stands and the diploma still bears the same title. The decline shows only after the decision-makers have taken their applause.

\emph{Status.} A professional class gains reputation from a norm whose daily enforcement falls on people of lower standing. The campaigner welcomes, and the clerk turns down the ineligible claim. The conservation office protects, and the shepherd finds the carcass. The university announces inclusion, and the junior teacher copes with the unfunded classroom. The dirty discretion is delegated downward and then condemned from above.

\emph{Epistemic.} The burden of uncertainty is placed on the critic. The advocate counts the intended benefit and requires opponents to prove every indirect cost beyond doubt. Uncertainty cuts both ways, though. Whoever proposes a large and irreversible intervention has a duty to look into its adverse consequences, and no right to shelter behind doubt about them.

None of this makes collective action impossible. Taxation always separates payer from beneficiary, and administration always delegates. The test is whether the separation is visible, authorized, proportionate, and open to correction. An invoice assigned openly under common rules is politics. An invoice erased so that its author can look innocent is propaganda, whatever he privately intended.

One further distinction is needed, between \emph{redistribution} and \emph{reduction}. A person who gives up a scarce good often creates an opening for someone else, and that can be exactly the purpose. Parents go without for their children, and taxpayers fund patients they will never meet. The sacrifice fails only where the declared aim was to reduce total use or total harm and the saving merely changed hands. If German restraint makes oil cheaper elsewhere, the fuel has been redistributed and little of it eliminated. If a strict standard drives the factory to a country with a lax one, virtue has changed the flag on the invoice. Sometimes the process goes on to common misery. Costs drive productive people and capital away, the tax base shrinks while the obligations remain, services decay, and more people leave. I am not arguing that the mobile rich should be appeased. I am arguing against the pretense that behavior does not respond. The moral wish writes the first move, and the later moves make the system.

\section{The pivot is the loss of correction}

A bad policy is not yet a turning point in the fate of a civilization. Policies are repealed, budgets revised, and officials replaced. The word \emph{pivot} should be kept for something narrower, the point at which correction becomes expensive because criticism is taken as an attack on the moral standing of the policy, its advocates, or its beneficiaries.

The bad sequence is familiar. A vivid problem is compressed into a story with a victim, a helper, and a villain. Institutions adopt targets that measure the visible act, such as admissions, money spent, predators counted, and domestic emissions cut. The costs turn up in other budgets, places, and years, and because they lie outside the target they are dismissed as anecdotes. Critics restore some of the missing context. They may be right, wrong, or half right, and the danger begins when their standing is judged before their evidence is. They are classified as selfish, hateful, anti-scientific, or extremist. Public agreement rises while private confidence falls. Meanwhile careers, budgets, and legal claims form around the existing rule, and revision grows dearer. At last protest returns in a cruder form. Those who would have discussed a limited correction lose ground to those who promise outright rejection, and the excesses of the latter are then used to discredit the earlier criticism too.

No step follows from the previous one by necessity, and a better sequence is just as possible. In it the vivid problem leads to a provisional commitment, and nobody makes an unlimited promise. The commitment comes with a budget, an estimate of capacity, a review date, and a statement of what would count against it. The people who meet the displaced costs report them early, and they are heard without first having to adopt the official moral vocabulary. Advocates can concede a cost without disowning the beneficiary. Critics can be proved right on a fact without acquiring a license for contempt. The rule is narrowed, widened, or ended according to results, and the correction is presented as part of keeping the promise. It is not treated as a confession that the promise was wicked.

So the pivot is no date after which truth is powerless. It is a change in the price of feedback. Before it, evidence can alter a policy without ruining everyone attached to the policy. After it, the same evidence has to fight identities, careers, budgets, and statutes as well. The task is to keep the price low enough for revision to remain an ordinary function of government and not become a crisis of legitimacy. This is the corridor of Chapter~\ref{chap:last-permissible-thought} seen from its practical side.

\section{An audit for goodness}

No formula removes tragedy, and no cost-benefit table prices every human value. A persecuted person may have a claim that exceeds any fiscal sum. A species may be worth keeping although no market pays for it. A constitutional right is not cancelled because it is costly to administer. An audit does not reduce morality to money. Its purpose is to stop moral language from conferring immunity on claims about cause and effect. Before any public act is called good, the following should be answered as concretely as the evidence allows.

\begin{enumerate}
\item \emph{What is the object?} Name the outcome and not the gesture. Is it lower global emissions or lower domestic fuel use? A viable species or every individual animal? Rescue from persecution or the largest number of arrivals? Secure tenants or affordable housing?
\item \emph{What happens otherwise?} Doing nothing is not the same as nothing changing.
\item \emph{Who are all the players?} Beneficiary, payer, displaced claimant, strategic responder, administrator, and future bearer.
\item \emph{What is rival now?} Beds, classrooms, flats, land, energy, time, legal capacity, and trust. Say how and when each bottleneck can be widened.
\item \emph{Where is the boundary?} Follow the effects across borders, supply chains, agencies, and generations.
\item \emph{How will behavior change?} Applicants, consumers, producers, landlords, owners of resources, officials, and voters will all respond. Do not assume that everyone but the policy-maker stands still.
\item \emph{What happens at scale?} Separate the single case from the general rule, and name the thresholds at which the effect changes.
\item \emph{What reciprocity sustains it?} What may be asked of recipients, contributors, and administrators, and what counts as defection?
\item \emph{Who has a voice?} Include those who live with the policy. Their voice need not be a veto, but distance should never confer authority without responsibility.
\item \emph{What will be measured?} Choose measures capable of revealing failure, and publish denominators, distributions, and uncertainty.
\item \emph{What is reversible?} Run pilots, set review dates, and keep alternatives alive.
\item \emph{What evidence would change the policy?} A measure that no possible result could count against is a creed and not a program.
\item \emph{May opponents speak without excommunication?} To dispute a policy's means is not to deny the human worth of its beneficiary. Punish threats and crimes, and leave inconvenient context alone.
\item \emph{Who signs the invoice?} Make the decision-maker name the tax, the labor, the risk, or the displaced use that pays for the good. Sentimental power thrives wherever the sacrifice is anonymous.
\end{enumerate}

An audit will not produce agreement. It will produce more honest disagreement. One person will accept a high cost for the sake of refuge, and another will give more weight to the continuity of institutions. One will value the wolf beyond any measurable service to the ecosystem, and another will put the mountain pastures first. The audit obliges them to argue about the same world.

\section{Bring the institutions with the gift}

Where correction has already become expensive, reform has to separate acknowledgment from surrender. A limit can be revised without declaring the beneficiary unworthy, and a critic can be partly right without being a saint. A few practices keep islands of responsibility above water. Attach advocacy to payment, so that whoever demands a public good also supports the tax, the building, the staff, and the enforcement it requires. Prefer bounded commitments to theatrical absolutes, since a defined number can be raised after success, while an unlimited moral claim can only be cut back amid cries of betrayal. Make policy conditional on capacity. Admission without housing and schools creates a queue, and nobody should call it hospitality. Protecting predators without prevention and compensation conscripts the rural population. Restraint in energy use with no theory of global substitution is domestic penance. Reward correction, because a culture that treats every revision as a confession of evil will be led by people who never admit a mistake. Keep judging individuals as individuals. And practice Axelrod's complete cooperation: open without hostility, answer exploitation, forgive those who return, and make the rule plain. A society that can only welcome has no means of defending the conditions of welcome, and it is a mistake to call that tolerance. A society that can only retaliate has turned order into revenge.

The director at FIAN did not really expect my country to arrive with the wine. His sentence joined gratitude to an exact sense of proportion. A bottle crosses a border in a suitcase. The productive order behind it does not. Shops, hospitals, laboratories, sound money, enforceable contracts, trained workers, and habits of cooperation are not a liquid that can be poured from one society into another.

That order was never morally pure. Prosperous societies contain inherited advantage, exclusion, waste, coercion, and luck. Their wealth does not justify every rule by which it was made, and their injustices do not make their capacities imaginary. The work is to tell the practices that generate trust and plenty from those that merely protect incumbents. Neither admiration nor resentment should be allowed to grow into a complete theory.

Moral stories belong to the mixture. A public promise can recruit effort, make suffering visible, and teach citizens what they owe one another. It becomes unreliable when its visible goodness is detached from the labor, money, restraint, and correction needed to keep it. The speaker's feeling of compassion does not decide the matter. Neither does a critic's discovery of hypocrisy. What decides it is whether the act is still connected to its consequences after other people have responded and the first invoice has come in. Numbers do not choose the good. But the good cannot claim exemption from numbers.

We cannot know in advance whether a strained institution will recover, or whether a generous policy will reproduce the conditions of further generosity. Large systems hide resilience as well as fragility. That is a reason to state limits, schedule reviews, and keep the ability to change course. It is no reason to refuse every gift until its remotest effects have been proved harmless.

Bring the wine. Share it, and enjoy the company it makes possible. When someone asks for the whole country, bring the arithmetic as well, along with the boundary, the institution, the right to say no, and the willingness to cooperate again. A lasting gift does not require giver or receiver to be innocent. It requires both to be still present when the consequences appear, and free to revise the terms on which they meet.


\bibliographystyle{plain}
\bibliography{references,svozil}

\end{document}
